\documentclass[11pt]{ctexart}
\usepackage[a4paper,margin=1in]{geometry}
\usepackage{graphicx}
\usepackage{xcolor}
\usepackage{hyperref}
\definecolor{refgray}{gray}{0.45}
\title{\textbf{市场最新观点汇总}\\[0.4em]{\large 通胀结构重塑、AI需求扩散与供给侧再定价成为三大主线，市场对供给约束的持续性定价不足。}}
\author{KC桌面 自动生成}
\date{260602}
\begin{document}
\maketitle
\tableofcontents
\newpage
\section{一页摘要}
\begin{itemize}
  \item 美国4月核心PCE同比升至3.29\%，核心CPI升至2.74\%，推动力来自地缘冲突、AI基础设施需求和关税传导三股结构性力量，市场低估了其持续性。
  \item 全球工业自动化周期出现罕见的“双顶”形态，中国以外地区的复苏正在加速，日本机床订单4月同比+45.1\%，机器人出口对欧洲+68.3\%。
  \item 韩国5月出口同比+53.2\%，但以吨位计出口量连续两个月同比下跌18\%，AI芯片价格暴涨制造了名义繁荣，非科技类出口集体走弱。
  \item 中国快递物流行业4月出现“量缓、收增、价升”的罕见组合，反内卷政策正从运动式转向制度化，行业定价权正在结构性回归。
  \item 美元定价的锚已从GDP切换到企业盈利，全球其他地区盈利增速与美国并驾齐驱，传统“好数据=强美元”逻辑失效。
  \item 亚洲新兴市场正从“增长韧性”切换至“通胀韧性”，韩国央行暗示最早7月加息，未来一年可能累计加息100个基点。
  \item 中国地产销售总量企稳但分化加剧，国企5月销售同比+6\%，民企同比-50\%，供给侧清洗正在重塑竞争格局。
  \item 全球资金5月加速涌入美国股票基金（86亿美元），但韩国散户连续两个月净卖出美国资产，机构与散户出现显著分歧。
\end{itemize}
\section{宏观与利率：通胀结构重塑，美元定价逻辑切换}
\textbf{美国通胀正被地缘冲突、AI需求和关税三股力量共同重塑，市场低估了其持续性；美元与经济增长的传统正相关关系正在断裂，企业盈利差距收窄削弱了美元走强的基础。}

\begin{itemize}
  \item 4月核心PCE同比升至3.29\%，核心CPI同比升至2.74\%，租金数据因去年10月政府关门出现技术性补涨，环比涨幅达0.55\%。
  \item 伊朗战争通过能源直接传导、非石油大宗商品（化肥、铝）和供应链成本重构三条路径同时推高通胀，预计全年推高核心PCE约0.35个百分点。
  \item AI相关需求推高电子元件价格同比+20\%，内存芯片、电池等产品涨价预计未来一年推高核心PCE约0.3个百分点。
  \item 关税对核心PCE的贡献已从峰值回落，预计年底仅剩0.1个百分点，但关税向消费价格的传导率已达90\%。
  \item 美国GDP增长预期重回G10前三，但全球其他地区企业盈利增速与美国相当（过去半年均超20\%），削弱了美元走强的盈利基础。
  \item 外国购买美国股票的速度在2026年明显放缓，此前2025年创纪录的买入节奏难以为继。
  \item 美国2年期利率近期大幅重定价，但美元基本未动，两者传统的正相关关系出现断裂。
  \item 韩国央行释放强烈鹰派信号，暗示最早7月加息，未来一年可能累计加息100个基点，此前市场预期仅50个基点。
  \item 菲律宾和印尼央行已因通胀冲击和汇率压力加息，印度央行仍在观望，但市场已在定价多次加息。
  \item 人民币中间价模型预测值降至6.7762，较前值大幅调降405个基点，韩元、卢布、欧元和澳元是主要贡献币种，定价逻辑回归贸易加权篮子框架。
\end{itemize}
\begin{figure}[htbp]
\centering
\includegraphics[width=0.88\linewidth]{figures/fig_001.jpg}
\caption{Exhibit 1 - Exhibit 1: Both Core PCE Inflation (+0.05pp to +3.29\%) and Core CPI Inflation (+0.14pp to +2.74\%) Accelerated on a Year-on-Year Basis in April}
\end{figure}
\begin{figure}[htbp]
\centering
\includegraphics[width=0.88\linewidth]{figures/fig_008.jpg}
\caption{Exhibit 8 - Exhibit 8: Crude Oil Prices Have Increased by Roughly 50\% Since the Start of the War in Iran; Our Oil Strategists Expect Brent Oil Prices to Fall to \textbackslash{}\$90 by 2026Q4, but See Two-Sided Risks to Prices Going Forward}
\end{figure}
\begin{figure}[htbp]
\centering
\includegraphics[width=0.88\linewidth]{figures/fig_010.jpg}
\caption{Exhibit 10 - Exhibit 10: We Estimate That the Boost from Commodity Price Increases to Year-Over-Year Core PCE Inflation Will Peak Around 0.35pp in 2026Q4 in Our Baseline Scenario for Oil Prices, Excluding the Impact via Imported Consumer Goods}
\end{figure}
\begin{figure}[htbp]
\centering
\includegraphics[width=0.88\linewidth]{figures/fig_012.jpg}
\caption{Exhibit 12 - Exhibit 12: AI-Related Demand Has Put Upward Pressure on Key Electronics Inputs}
\end{figure}
\begin{figure}[htbp]
\centering
\includegraphics[width=0.88\linewidth]{figures/fig_016.jpg}
\caption{Exhibit 16 - Exhibit 16: We Estimate That Passthrough from Tariff Costs to Consumer Prices Has Reached 90\% So Far Share of Tariff Costs Passed onto Core Consumer Prices, Based on CPI Data Through April}
\end{figure}
\begin{figure}[htbp]
\centering
\includegraphics[width=0.88\linewidth]{figures/fig_057.jpg}
\caption{Figure 8 - Figure 1: The US has returned to being a high growth country within G10}
\end{figure}
\begin{figure}[htbp]
\centering
\includegraphics[width=0.88\linewidth]{figures/fig_060.jpg}
\caption{Figure 4 - Figure 4: US earning are very strong, but so are RoW earnings}
\end{figure}
\begin{figure}[htbp]
\centering
\includegraphics[width=0.88\linewidth]{figures/fig_061.jpg}
\caption{Figure 5 - Figure 5: Foreign buying of US equities has slowed a lot in 2026, after record buying last year}
\end{figure}
\begin{figure}[htbp]
\centering
\includegraphics[width=0.88\linewidth]{figures/fig_063.jpg}
\caption{Figure 7 - Figure 7: There's been a big front-end repricing for the US vs peers, but the USD is unmoved}
\end{figure}
\begin{figure}[htbp]
\centering
\includegraphics[width=0.88\linewidth]{figures/fig_065.jpg}
\caption{Figure 1 - Figure 1: EMAX tanker imports}
\end{figure}
\begin{figure}[htbp]
\centering
\includegraphics[width=0.88\linewidth]{figures/fig_066.jpg}
\caption{Figure 2 - Figure 2: India tanker imports}
\end{figure}
{\small\color{refgray}\textbf{References:} [R001] 通胀正在被三股力量重塑，市场低估了它们的持续性; [R004] 市场真正低估的不是需求，而是美元定价的结构性断裂; [R005] 亚洲能源冲击的真正赢家不是增长，而是央行的鹰派转向; [R017] 人民币中间价模型正在发出一个被忽视的定价信号}

\section{AI与半导体：需求扩散至多节点，供给约束成为新定价变量}
\textbf{AI对半导体的拉动已从训练驱动切换至推理+存储+封装的多波次共振，WFE长期可见性前所未有地清晰，但洁净室产能和内存短缺成为真正的瓶颈。}

\begin{itemize}
  \item AI训练需求首先拉动先进逻辑/代工，AI推理爆发正驱动NAND和CPU需求上升，WFE驱动力从单一节点向多节点扩散。
  \item Applied Materials预计AI驱动的先进逻辑、DRAM（HBM）和先进封装将贡献2026年增量WFE支出的80\%。
  \item NVIDIA Rubin芯片按计划推进，市场流传的“重新流片”或“延迟”系误读，Blackwell生产反而加速，生命周期延长。
  \item 内存短缺是当前AI建设最大瓶颈，DRAM三季度价格预期涨20\%以上，供应商不急于谈价。
  \item HDD供需缺口正在加速扩大，CY26缺口从200EB上修至300EB，CY27和CY28均达400EB，供应商目标价到CY28可能翻倍。
  \item 大容量HDD混合均价当前约1.5美分/GB，供应商目标CY27达2.5美分/GB，CY28达3美分/GB，超大规模客户“只要求更多供应”。
  \item Broadcom 2026年下半年AI收入因机架部署推迟，约210亿美元机会大部分延至2027年，但需求未消失，客户基础正在扩大。
  \item Broadcom当前AI市场份额仅约10\%，预计将持续提升，2027年XPU份额的再定价空间是估值中枢的关键。
  \item 韩国5月半导体出口同比+169.4\%，DRAM出口+369.8\%，NAND+206.8\%，但以吨位计出口量连续两个月同比下跌18\%，增长由价格驱动。
  \item HBM相关材料进口大涨：海力士环氧树脂进口+57\%，三星塑料薄膜进口+96\%，可作为HBM出货量的前瞻信号。
\end{itemize}
\begin{figure}[htbp]
\centering
\includegraphics[width=0.88\linewidth]{figures/fig_323.jpg}
\caption{Exhibit 2 - Exhibit 2: Weekly stock performance}
\end{figure}
\begin{figure}[htbp]
\centering
\includegraphics[width=0.88\linewidth]{figures/fig_326.jpg}
\caption{Exhibit 8 - Exhibit 8: 2025-27e Revenue CAGR}
\end{figure}
\begin{figure}[htbp]
\centering
\includegraphics[width=0.88\linewidth]{figures/fig_327.jpg}
\caption{Exhibit 9 - Exhibit 9: 2025-27e EPS CAGR}
\end{figure}
\begin{figure}[htbp]
\centering
\includegraphics[width=0.88\linewidth]{figures/fig_328.jpg}
\caption{Exhibit 10 - Exhibit 10: Semiconductor Company inventory is at 111 days, up 3 days q/q which is 4 days behind of the seasonal decrease of 1 days. DOI is 22 days above the historical median and trailing 4-quarters slightly increased.. Semicondu}
\end{figure}
\begin{figure}[htbp]
\centering
\includegraphics[width=0.88\linewidth]{figures/fig_329.jpg}
\caption{Exhibit 11 - Exhibit 11: Semi Customer DOI decreased 6 days sequentially to 52 days, below the seasonal decrease of 3 days. DOI is below the historical median and trailing 4 quarters increased from last quarter.}
\end{figure}
\begin{figure}[htbp]
\centering
\includegraphics[width=0.88\linewidth]{figures/fig_330.jpg}
\caption{Exhibit 12 - Exhibit 12: Distributor inventory is at 63 days, down 2 days q/q, below the seasonal decline of 0 days. DOI is 9 days above the historical median while trailing 4-quarters decreased sequentially. Distributors (DOI)}
\end{figure}
\begin{figure}[htbp]
\centering
\includegraphics[width=0.88\linewidth]{figures/fig_356.jpg}
\caption{Exhibit 1 - Exhibit 1: Import trends of epoxide resins have been showing a strong correlation with Hynix's HBM shipments Imports trend for epoxide resins from Japan to Hynix}
\end{figure}
\begin{figure}[htbp]
\centering
\includegraphics[width=0.88\linewidth]{figures/fig_357.jpg}
\caption{Exhibit 2 - Exhibit 2: Imports trend for plastic film has been showing a strong correlation with SEC's HBM shipment Import trend of plastic film from Japan to SEC}
\end{figure}
{\small\color{refgray}\textbf{References:} [R021] 硬件市场真正稀缺的不是需求，而是被低估的供给议价权; [R022] 市场真正低估的不是需求，而是供给侧的再定价; [R025] 半导体CEO们的共识：AI需求正在重塑WFE定价权，但真正的约束是物理空间; [R027] 市场低估的不是Broadcom的短期业绩，而是其2027年XPU份额的再定价空间; [R032] 韩国半导体出口的真相：市场低估了供给约束的持久性}

\section{能源与大宗：供给恢复慢于预期，锂矿重启信号密集}
\textbf{霍尔木兹海峡供应恢复将比市场预期的更慢更曲折，油价近端有支撑；锂价反弹正激励澳洲矿山低成本重启，但累积效应可能在2027年后将市场推向小幅过剩。}

\begin{itemize}
  \item 布伦特原油当前约92美元/桶，经通胀调整后恰处20年中位数，但市场定价隐含了过于乐观的快速恢复假设。
  \item 近月布伦特CFD从5.7美元/桶降至2.6美元/桶，M1-M2价差从3.8美元降至不足1美元，市场在过去三周变得过于乐观。
  \item 供应恢复面临六大瓶颈：扫雷与保险、油轮错位、储油罐满、油田重启困难、基础设施损坏、人力短缺。
  \item MS维持三季度布伦特100美元/桶预测，上调四季度至95美元，2027年一季度至85美元。
  \item 全球海上钻井需求到2030年预计增长12\%，深水部分增长17\%，而供给端预计2027-2030年几乎无新增。
  \item 深水钻井平台利用率预计到2028年升至约85\%，历史上该水平支撑很强的定价权，日费率有望回到2023年峰值。
  \item 锂辉石精矿价格从不足800美元/吨回升至1500美元以上，澳洲矿山重启信号密集出现。
  \item Mineral Resources的Bald Hill重启将增加14万吨/年SC6精矿产能，首批货物最早7月发出。
  \item 2026年全球锂市场预计仍有约11.7万吨LCE缺口，但重启产能累积效应可能在2027-2029年将市场推向小幅过剩。
  \item 欧洲钢铁进口配额较2024年削减约47\%，超出配额关税提升至50\%，CBAM 2026年正式执行，利润池正向本土钢厂迁移。
\end{itemize}
\begin{figure}[htbp]
\centering
\includegraphics[width=0.88\linewidth]{figures/fig_242.jpg}
\caption{Exhibit 20 - Exhibit 1: Utilization reaches \textbackslash{}\textasciitilde{}85\%+ by '28 in our forecasts, a level that supports strong driller pricing power...}
\end{figure}
\begin{figure}[htbp]
\centering
\includegraphics[width=0.88\linewidth]{figures/fig_243.jpg}
\caption{Exhibit 2 - Exhibit 2: ...Pushing leading-edge dayrates for higher-spec rigs back to recent \textbackslash{}\textasciitilde{}'23 peak levels by \textbackslash{}\textasciitilde{}'27-'28.}
\end{figure}
\begin{figure}[htbp]
\centering
\includegraphics[width=0.88\linewidth]{figures/fig_244.jpg}
\caption{Exhibit 9 - Exhibit 3: We see rig counts increasing across global supply segments to support O\&G demand growth, supply diversification, and to build strategic reserves. We see the most growth offshore ('30E vs '25 \%Δ – OS DW: +17\% / OS Shelf:}
\end{figure}
\begin{figure}[htbp]
\centering
\includegraphics[width=0.88\linewidth]{figures/fig_249.jpg}
\caption{Exhibit 9 - Exhibit 9: Offshore rig tendering activity in the deepwater space has been robust YTD. The \textbackslash{}\textasciitilde{}75 rig years awarded as of late-May is \textbackslash{}\textasciitilde{}70\% / \textbackslash{}\textasciitilde{}30 rig years ahead of historical trend and represents the most cumulative rig years award}
\end{figure}
\begin{figure}[htbp]
\centering
\includegraphics[width=0.88\linewidth]{figures/fig_550.jpg}
\caption{Exhibit 1 - Exhibit 1: In inflation-adjusted terms, the current Brent price sits precisely at the 50th percentile of the price distribution over the last 20 years - the most-common price despite the most-unusual supply disruption Inflation-adj}
\end{figure}
\begin{figure}[htbp]
\centering
\includegraphics[width=0.88\linewidth]{figures/fig_551.jpg}
\caption{Exhibit 3 - Exhibit 3: Crude oil production from the Middle East remains under pressure...}
\end{figure}
\begin{figure}[htbp]
\centering
\includegraphics[width=0.88\linewidth]{figures/fig_553.jpg}
\caption{Exhibit 5 - Exhibit 5: The number of daily tanker transits via the Strait remains a fractions of its historic level Strait of Hormuz}
\end{figure}
\begin{figure}[htbp]
\centering
\includegraphics[width=0.88\linewidth]{figures/fig_555.jpg}
\caption{Exhibit 7 - Exhibit 7: On our new assumption, the cumulative supply loss vs a flat pre-conflict production would be 2.6 bn bbl in 2026 Production recovery assumptions}
\end{figure}
\begin{figure}[htbp]
\centering
\includegraphics[width=0.88\linewidth]{figures/fig_358.jpg}
\caption{Exhibit 1 - Exhibit 1: Australian supply response is building as lithium prices rise (LHS) which could bring upside risk to modest CY27-29 market surpluses on MSe (RHS)}
\end{figure}
{\small\color{refgray}\textbf{References:} [R023] 市场低估的不是需求，而是供给纪律的自我强化; [R033] 锂价反弹的真正考验：供给侧正在醒来，但市场低估了重启的节奏; [R034] 欧洲钢铁业正在经历一场被低估的供给侧重定价; [R035] 市场真正低估的不是需求，而是供给恢复的速度}

\section{中国地产与消费：销售分化加剧，储蓄率是消费真正拐点}
\textbf{新房销售总量企稳但国企与民企增速裂口超50个百分点，供给侧清洗不可逆；消费疲软核心驱动不是收入而是财富效应，储蓄率何时下降才是真正拐点。}

\begin{itemize}
  \item 5月百强房企销售额同比-4\%，与4月-5\%基本持平，但国企销售同比+6\%，民企同比-50\%，困境房企-30\%，国企与民企增速差距高达56个百分点。
  \item 华润置地5月同比+28\%，招商蛇口+20\%，越秀地产+18\%，中海地产+14\%，国企头部玩家逆势增长。
  \item 第22周新房成交面积环比+13\%，同比转正至+4\%，库存月数从4月均值29.3个月下降至27.6个月。
  \item 广州启动全市收储，中心区300万以内、70平以下房源优先，同时优化商贷转公积金政策。
  \item 4月社零增速仅0.2\%，城镇居民可支配收入增速维持在4\%左右，但预防性储蓄率高达33\%-34\%，消费倾向被系统性压低。
  \item 二手房价同比跌约12\%，负财富效应对消费意愿的压制短期内难以消除。
  \item 餐饮（+2.2\%）、化妆品（+4.7\%）、运动服饰等品类逆势增长，消费者仍愿为“悦己”和健康生活方式花钱。
  \item 城市更新计划目标改造50万套老旧住房、4000个城中村，但未提及现金补偿，对商品房销售拉动有限。
  \item 5月新能源车渗透率预计再次站上60\%，比亚迪海外销量16.06万辆同比暴增80\%，创历史新高。
  \item 蔚来5月交付3.77万辆同比+62\%，小鹏3.22万辆同比-4\%，理想3.34万辆同比-18\%，新势力内部分化明显。
\end{itemize}
\begin{figure}[htbp]
\centering
\includegraphics[width=0.88\linewidth]{figures/fig_136.jpg}
\caption{Figure 3 - Figure 3: Key developers' contracted sales Y/Y growth by type}
\end{figure}
\begin{figure}[htbp]
\centering
\includegraphics[width=0.88\linewidth]{figures/fig_137.jpg}
\caption{Figure 4 - Figure 4: Key developers' contracted sales vs. 2018-21 average by type}
\end{figure}
\begin{figure}[htbp]
\centering
\includegraphics[width=0.88\linewidth]{figures/fig_139.jpg}
\caption{Figure 6 - Figure 6: Key developers' May 2026 contracted sales Y/Y growth}
\end{figure}
\begin{figure}[htbp]
\centering
\includegraphics[width=0.88\linewidth]{figures/fig_140.jpg}
\caption{Exhibit 1 - Exhibit 1: Primary GFA sold last week was +13\% wow and +4\% yoy in c.75 cities}
\end{figure}
\begin{figure}[htbp]
\centering
\includegraphics[width=0.88\linewidth]{figures/fig_142.jpg}
\caption{Exhibit 3 - Exhibit 3: Secondary GFA sold last week was -2\% wow and +26\% yoy in c.20 cities Average weekly volume of secondary property sales}
\end{figure}
\begin{figure}[htbp]
\centering
\includegraphics[width=0.88\linewidth]{figures/fig_144.jpg}
\caption{Exhibit 5 - Exhibit 5: Average CSI was +0.2pp wow and +3.7pp yoy Weekly Centraline Salesman Index (CSI) tracker in 5 cities}
\end{figure}
\begin{figure}[htbp]
\centering
\includegraphics[width=0.88\linewidth]{figures/fig_145.jpg}
\caption{Exhibit 6 - Exhibit 6: Average CAI was -0.5pp wow and -3.6pp yoy Weekly Centraline Seller Asking Index (CAI) tracker in 6 cities}
\end{figure}
\begin{figure}[htbp]
\centering
\includegraphics[width=0.88\linewidth]{figures/fig_147.jpg}
\caption{Exhibit 8 - Exhibit 8: Inventory month was -0.8\% wow, representing -1.1\% from end-25 levels c.20 cities' inventory months (12mth rolling) breakdown by city tier}
\end{figure}
{\small\color{refgray}\textbf{References:} [R012] 市场真正低估的不是需求，而是供给侧的分化; [R013] 市场真正低估的不是需求，而是供给侧的再定价; [R014] 中国消费的真正拐点，不在于收入增速，而在于储蓄率何时开始下降; [R016] 中国车市真正的分化不在电动化，而在“谁的海外能扛”}

\section{全球贸易与供应链：补库信号已现，有效关税税率差异重塑格局}
\textbf{中国对美出口同比正增长已超过单纯基数效应，补库存行为正在发生；有效关税税率差异正在重塑亚洲各经济体对美出口的竞争格局。}

\begin{itemize}
  \item 截至5月28日当周，中国至美国重箱船舶数量环比-13\%，但同比仍+17\%，TEU同比+20.8\%。
  \item 洛杉矶港未来1-2周进口量预计同比增长54\%和46\%，补库行为正在发生。
  \item 亚洲不同经济体对美出口增速出现分化，中国增速领先，部分经济体因有效税率较低可能增加对美出口。
  \item 中国物流快递行业4月包裹量同比+3.2\%，但收入增速+6.3\%，单票收入恢复至7.65元同比+3\%，出现“量缓、收增、价升”的罕见组合。
  \item 反内卷政策正从运动式转向制度化，价格纪律、基层福利和网络稳定性三大支柱使重回价格战的可能性极低。
  \item 中通日处理退货包裹1000万件，逆物流能力成为护城河；韵达主动砍掉超低价订单，均价涨10\%。
  \item 4月线上实物商品GMV仅增长0.2\%，家电跌15\%，手机增速从27\%骤降至6\%，电商端需求明显降温。
  \item 韩国5月出口对东盟环比+17.4\%贡献全部增量，对美出口连续8个月增长，对欧盟出口环比-4.2\%唯一负增长。
  \item 韩国贸易顺差扩大到269亿美元创历史新高，但结构上越来越依赖半导体单一品类。
  \item 全球资金5月涌入美国股票基金86亿美元创1月以来新高，但新兴市场ETF流入从4月58亿美元骤降至25亿美元。
\end{itemize}
\begin{figure}[htbp]
\centering
\includegraphics[width=0.88\linewidth]{figures/fig_034.jpg}
\caption{Exhibit 1 - Exhibit 1: Container ships to the US remain volatile; levels are still down from post-Liberation day increases Daily Laden Container Ship Vessels from China to US}
\end{figure}
\begin{figure}[htbp]
\centering
\includegraphics[width=0.88\linewidth]{figures/fig_035.jpg}
\caption{Exhibit 2 - Exhibit 2: Laden container ships from China to US have declined on a YoY basis in the most recent week Daily Laden Container Ship Vessels from China to US, YoY}
\end{figure}
\begin{figure}[htbp]
\centering
\includegraphics[width=0.88\linewidth]{figures/fig_037.jpg}
\caption{Exhibit 4 - Exhibit 4: TEU growth was up in the most recent week Daily TEU from China to the US, YoY}
\end{figure}
\begin{figure}[htbp]
\centering
\includegraphics[width=0.88\linewidth]{figures/fig_038.jpg}
\caption{Exhibit 6 - Exhibit 6: Mainland China and Asia-Ex Mainland: vessels were up YoY for Mainland China and Asia Ex Mainland China Daily Laden Container Ship Vessels from Asia ex-Mainland China and Mainland China to US, YoY}
\end{figure}
\begin{figure}[htbp]
\centering
\includegraphics[width=0.88\linewidth]{figures/fig_039.jpg}
\caption{Exhibit 7 - Exhibit 7: TEU growth from Mainland China was up vs Asia-Ex Mainland China was down this past week Daily TEU from Asia ex-Mainland China and Mainland China to US, YoY}
\end{figure}
\begin{figure}[htbp]
\centering
\includegraphics[width=0.88\linewidth]{figures/fig_042.jpg}
\caption{Exhibit 10 - Exhibit 10: Planned TEUs are expected to increase sequentially next week and 2 weeks out Planned TEUs into the Port of LA}
\end{figure}
\begin{figure}[htbp]
\centering
\includegraphics[width=0.88\linewidth]{figures/fig_026.jpg}
\caption{Figure 3 - Figure 3: China express monthly industry volume (daily) Industry daily volume trend}
\end{figure}
\begin{figure}[htbp]
\centering
\includegraphics[width=0.88\linewidth]{figures/fig_027.jpg}
\caption{Figure 4 - Figure 4: China express monthly industry ASP Industry pricing trend}
\end{figure}
\begin{figure}[htbp]
\centering
\includegraphics[width=0.88\linewidth]{figures/fig_032.jpg}
\caption{Figure 9 - Figure 9: Online retail sales monthly YoY growth}
\end{figure}
\begin{figure}[htbp]
\centering
\includegraphics[width=0.88\linewidth]{figures/fig_033.jpg}
\caption{Figure 10 - Figure 10: Online physical goods sales monthly YoY growth}
\end{figure}
\begin{figure}[htbp]
\centering
\includegraphics[width=0.88\linewidth]{figures/fig_354.jpg}
\caption{Exhibit 1 - Exhibit 1: Sustained increases in May exports... Exhibit 2: Sustained increases in semiconductors partially offset by weakness in other exports}
\end{figure}
\begin{figure}[htbp]
\centering
\includegraphics[width=0.88\linewidth]{figures/fig_355.jpg}
\caption{Exhibit 3 - Exhibit 3: Exports to the US and EM increased further}
\end{figure}
{\small\color{refgray}\textbf{References:} [R002] 物流行业真正被低估的不是需求，而是定价权的结构性重塑; [R003] 市场真正低估的不是需求，而是供给侧的再定价; [R006] 全球资金正在重新定价AI，但亚洲散户已经提前离场; [R031] 韩国出口的“分裂式繁荣”：半导体孤军深入，而结构性风险正在积累}

\section{区域市场：亚洲央行鹰派转向，澳洲投资逻辑切换}
\textbf{亚洲新兴市场正从增长韧性切换至通胀韧性，央行关注点转向控通胀；澳大利亚投资逻辑从需求周期转向供给侧结构性再定价。}

\begin{itemize}
  \item 亚洲（除中国）通过多元化采购、动用战略储备和燃料替代，5月能源进口量已快速恢复正常水平。
  \item 4月IP数据表明区域科技势头依然强劲，配合部分经济体财政刺激，能源冲击影响“看得见但没那么疼”。
  \item 韩国央行暗示最早7月加息，未来一年可能累计加息100个基点，核心逻辑是科技驱动的增长势头太猛，担心需求侧通胀溢出。
  \item 印度5月制造业PMI 55.0，连续第三个月扩张，但新出口订单指数从56.3骤降至53.7，内需与外需呈现“一热一冷”格局。
  \item 印度制造业投入价格指数仍在四年高位附近，企业抱怨能源、燃料、原材料和运输成本上涨，产出价格指数下降说明提价能力有限。
  \item 澳大利亚资本支出驱动力来自国防（追加530亿澳元）、能源转型（150亿澳元）和数据中心建设，几乎不受短期利率影响。
  \item MS预计澳大利亚六大类资本开支从FY25的2250亿澳元增长到FY30的3260亿澳元。
  \item 澳大利亚市场共识预期FY26/FY27盈利增速超10\%，但MS认为FY27预期将大幅下调至约6\%。
  \item 日本IT服务业网络基础设施建设订单爆发，IIJ系统集成业务4Q订单同比+51\%，积压同比+37.7\%。
  \item 日本机床订单4月同比+45.1\%，来自欧洲订单增速从21.9\%跃升至35.6\%，日本本土从2.5\%飙升至43.4\%。
\end{itemize}
\begin{figure}[htbp]
\centering
\includegraphics[width=0.88\linewidth]{figures/fig_065.jpg}
\caption{Figure 1 - Figure 1: EMAX tanker imports}
\end{figure}
\begin{figure}[htbp]
\centering
\includegraphics[width=0.88\linewidth]{figures/fig_066.jpg}
\caption{Figure 2 - Figure 2: India tanker imports}
\end{figure}
\begin{figure}[htbp]
\centering
\includegraphics[width=0.88\linewidth]{figures/fig_111.jpg}
\caption{Exhibit 1 - Exhibit 1: Manufacturing PMI index edged up in May driven by domestic demand, while new export orders declined}
\end{figure}
\begin{figure}[htbp]
\centering
\includegraphics[width=0.88\linewidth]{figures/fig_112.jpg}
\caption{Exhibit 2 - Exhibit 2: Both the input and the output price index declined from their April highs, but remain elevated}
\end{figure}
\begin{figure}[htbp]
\centering
\includegraphics[width=0.88\linewidth]{figures/fig_331.jpg}
\caption{Exhibit 2 - Exhibit 2: The outlook for domestic demand conditions has deteriorated meaningfully}
\end{figure}
\begin{figure}[htbp]
\centering
\includegraphics[width=0.88\linewidth]{figures/fig_332.jpg}
\caption{Exhibit 3 - Exhibit 3: But elevated inflation is likely to limit how dovish the RBA is able to turn}
\end{figure}
\begin{figure}[htbp]
\centering
\includegraphics[width=0.88\linewidth]{figures/fig_333.jpg}
\caption{Exhibit 4 - Exhibit 4: We forecast that capex in six key categories will rise from A\textbackslash{}\$225bn in FY25 to A\textbackslash{}\$326bn in FY30}
\end{figure}
\begin{figure}[htbp]
\centering
\includegraphics[width=0.88\linewidth]{figures/fig_334.jpg}
\caption{Exhibit 6 - Exhibit 6: Consensus earnings expectations are for robust market earnings growth over coming years}
\end{figure}
\begin{figure}[htbp]
\centering
\includegraphics[width=0.88\linewidth]{figures/fig_335.jpg}
\caption{Exhibit 7 - Exhibit 7: Our top-down earnings models point to a much more cautious earnings outlook}
\end{figure}
\begin{figure}[htbp]
\centering
\includegraphics[width=0.88\linewidth]{figures/fig_185.jpg}
\caption{EXHIBIT 37 - EXHIBIT 37: Japan total machine tool order growth trend Japan Monthly Machine Tool Orders \& YoY Growth}
\end{figure}
\begin{figure}[htbp]
\centering
\includegraphics[width=0.88\linewidth]{figures/fig_186.jpg}
\caption{EXHIBIT 38 - EXHIBIT 38: Japan machine tool orders growth by region}
\end{figure}
{\small\color{refgray}\textbf{References:} [R005] 亚洲能源冲击的真正赢家不是增长，而是央行的鹰派转向; [R008] 印度制造业的真正信号：内需托底，但成本压力正在重塑利润分配格局; [R015] 全球工业技术周期的“双顶”正在确认，但市场对“中国需求之外的结构性加速”定价不足; [R028] 市场真正低估的不是需求，而是供给侧的再定价; [R037] 日本IT服务业的真正拐点，藏在两家非覆盖公司的财报细节里}

\section{医疗健康：双抗改写肺癌标准治疗，口服GLP-1竞赛参数被重新定义}
\textbf{PD-1xVEGF双抗首次在III期头对头试验中击败PD-1标准疗法，肺癌一线治疗格局可能被重新定义；口服GLP-1竞争焦点从减重幅度转向耐受性，中国生物科技公司正从跟随者变为参数定义者。}

\begin{itemize}
  \item HARMONi-6研究显示依沃西单抗（PD-1xVEGF双抗）联合化疗中位OS 27.9个月 vs 对照组23.7个月，HR=0.66，死亡风险降低34\%。
  \item 这是首个在III期试验中证明比PD-1联合化疗更优的方案，PD-L1阴性患者同样获益。
  \item KOL明确表示不会仅凭一个中国单中心研究调整临床决策，FDA批准+全球HARMONi-3阳性结果才是临床落地的必要条件。
  \item RET靶向药selpercatinib在RET阳性早期NSCLC辅助治疗中EFS的HR仅0.17，但约70\%的RET阳性患者是不吸烟者，不在常规肺癌筛查范围。
  \item 口服GLP-1到2030年市场规模预计超350亿美元，占肥胖药物总可寻址市场约35\%。
  \item 口服GLP-1竞争焦点不是减重幅度而是耐受性，Wegovy口服片和orforglipron已将减重门槛推至13.9\%和11.5\%。
  \item 信达生物将在ADA公布IBI3032（口服小分子GLP-1）1期数据，4周治疗达到10\%减重效果，比同类竞品晚约2.5年。
  \item 恒瑞医药展示全管线GLP-1数据，策略是做全谱系覆盖，包括口服GLP-1/GIP肽HRS9531的2期肥胖数据。
  \item 劳拉替尼7年随访数据进一步巩固ALK赛道“天花板”地位。
  \item 中国双抗在肠癌早期数据中ORR达70.8\%，DCR 100\%，支持正在进行的全球III期研究。
\end{itemize}
\begin{figure}[htbp]
\centering
\includegraphics[width=0.88\linewidth]{figures/fig_352.jpg}
\caption{Exhibit 2 - Exhibit 2: \$15\textbackslash{}\%\$ weight loss is becoming the bar for newer oral GLP-1 after the approval of Wegovy pills and orforglipron... Weight loss data of oral GLP-1 with different treatment duration (placebo-adjusted)}
\end{figure}
\begin{figure}[htbp]
\centering
\includegraphics[width=0.88\linewidth]{figures/fig_353.jpg}
\caption{Exhibit 3 - The chart includes a legend indicating 'GLP-1RA with China plan in obesity' and a dashed line representing the semaglutide curve. \textbackslash{}*treatment arm data (placebo unadjusted)}
\end{figure}
{\small\color{refgray}\textbf{References:} [R026] 肺癌治疗正在经历一场“分子分型”驱动的定价重构; [R030] 口服GLP-1的竞赛逻辑已经变了：中国资产正在定义下一轮竞争参数; [R036] ASCO 2026最值得关注的信号：双抗在肺癌领域首次击败PD-1标准疗法}

\section{工业科技：全球自动化周期出现“双顶”，中国以外复苏加速}
\textbf{全球工业自动化周期在2026年二季度出现罕见的“双顶”形态，中国以外地区的复苏正在显著加速，市场对后者的定价远不够充分。}

\begin{itemize}
  \item 中国5月高技术制造业PMI扩张至52.9，装备制造PMI 52.1，高端制造仍在加速。
  \item 4月中国工业利润同比增速进一步加快至24.7\%，前4个月高端制造利润增长44.8\%，明显跑赢整体工业。
  \item 日本工业机器人出口量4月同比+40.7\%，对欧洲出口增速从35.3\%加速至68.3\%，对亚洲（除中国）从24.0\%飙升至88.6\%。
  \item 日本伺服电机产量3月+25\%，减速器+35.9\%，气动元件+27.2\%，自动化产业链全线回暖。
  \item 美国5月PMI 55.3，日本54.5，欧元区51.4，全部处于扩张区间，全球制造业同步扩张。
  \item 中国软件行业4月收入同比+8.9\%较3月放缓，但净利率从9.0\%跃升至12.0\%，创近12个月相对高位。
  \item 半导体设计软件以18.5\%增速领跑中国软件行业，云计算和大数据+12.3\%，网络安全仅+5\%明显跑输。
  \item AI Agent落地场景从概念验证进入加速变现阶段，尤其在知识留存和生产流程优化领域。
  \item 中国零食行业量贩渠道开店速度超预期，年初至今净开店约2500家/系统，全年目标4000-5000家。
  \item 魔芋原料价格已从高点回落25\%，成本端迎来顺风，盐津铺子和卫龙美味下半年净利润增速预计可达27\%和12\%。
\end{itemize}
\begin{figure}[htbp]
\centering
\includegraphics[width=0.88\linewidth]{figures/fig_150.jpg}
\caption{EXHIBIT 2 - EXHIBIT 2: China manufacturing PMI trend}
\end{figure}
\begin{figure}[htbp]
\centering
\includegraphics[width=0.88\linewidth]{figures/fig_154.jpg}
\caption{EXHIBIT 6 - EXHIBIT 6: Value-add growth of high technology manufacturing and general manufacturing}
\end{figure}
\begin{figure}[htbp]
\centering
\includegraphics[width=0.88\linewidth]{figures/fig_155.jpg}
\caption{EXHIBIT 7 - EXHIBIT 7: Growth of industrial profit and inventory of finishd goods in China (yoy)}
\end{figure}
\begin{figure}[htbp]
\centering
\includegraphics[width=0.88\linewidth]{figures/fig_164.jpg}
\caption{EXHIBIT 16 - EXHIBIT 16: China factory automation shipment value growth trend by product (YoY)}
\end{figure}
\begin{figure}[htbp]
\centering
\includegraphics[width=0.88\linewidth]{figures/fig_165.jpg}
\caption{EXHIBIT 17 - EXHIBIT 17: Japan machine tool orders from China and Y/Y growth. Japan machine tool orders from China}
\end{figure}
\begin{figure}[htbp]
\centering
\includegraphics[width=0.88\linewidth]{figures/fig_176.jpg}
\caption{EXHIBIT 28 - EXHIBIT 28: Growth rate of Japan robot export to China Japan robot export to China}
\end{figure}
\begin{figure}[htbp]
\centering
\includegraphics[width=0.88\linewidth]{figures/fig_177.jpg}
\caption{EXHIBIT 29 - EXHIBIT 29: China industrial robot production and Y/Y growth China industrial robots production}
\end{figure}
\begin{figure}[htbp]
\centering
\includegraphics[width=0.88\linewidth]{figures/fig_128.jpg}
\caption{Exhibit 3 - Exhibit 1: China software industry growth was at \$11\textbackslash{}\%\$ YoY in 4M26 4M26 China Software industry revenue growth by segment}
\end{figure}
\begin{figure}[htbp]
\centering
\includegraphics[width=0.88\linewidth]{figures/fig_129.jpg}
\caption{Exhibit 2 - Exhibit 2: Software industry growth slowed in Apr 2026; SMB PMI was weak in May 2026}
\end{figure}
\begin{figure}[htbp]
\centering
\includegraphics[width=0.88\linewidth]{figures/fig_132.jpg}
\caption{Exhibit 6 - Exhibit 6: Net margin at 12.0\% in Apr (vs. 9.0\% in Mar 2026), leading to 4M26 NM at 11.4\% Software companies registered in China: aggregate NI YoY growth and net margin}
\end{figure}
\begin{figure}[htbp]
\centering
\includegraphics[width=0.88\linewidth]{figures/fig_350.jpg}
\caption{Exhibit 1 - Exhibit 1: Weilong/Yanker offline market share remain strong as of April-26 at \$52.9\textbackslash{}\% / 11.0\textbackslash{}\%\$ vs last year's \$51.9\textbackslash{}\% / 8.3\textbackslash{}\%\$ Konjac snack market share in offline}
\end{figure}
\begin{figure}[htbp]
\centering
\includegraphics[width=0.88\linewidth]{figures/fig_351.jpg}
\caption{Exhibit 2 - Exhibit 2: Konjac raw material costs spiked in 2024 by \$60\textbackslash{}\%\$ yoy and we now expect raw material costs to decline by \$15\textbackslash{}\%\$ yoy in 2026E in our base case}
\end{figure}
{\small\color{refgray}\textbf{References:} [R011] 中国软件行业4月增速放缓至8.9\%，但利润率改善揭示了一个被低估的结构性信号; [R015] 全球工业技术周期的“双顶”正在确认，但市场对“中国需求之外的结构性加速”定价不足; [R029] 市场低估的不是零食需求，而是成本周期与渠道红利的共振}

\section{风险偏好与资金流向：机构与散户出现显著分歧}
\textbf{全球资金5月加速涌入美国股票基金，但韩国和台湾散户投资者连续两个月净卖出美国资产，机构与散户之间的信息分层和预期差正在扩大。}

\begin{itemize}
  \item 5月最后一周外国资金净流入美国股票基金27.25亿美元，是前一周13.54亿美元的两倍。
  \item 5月全月流入美国股票基金的外资总额达86亿美元，创2026年1月以来最高单月纪录。
  \item 新兴市场股票ETF资金流入从4月58.32亿美元骤降至5月24.18亿美元，资金在“回归”发达市场。
  \item 韩国散户投资者5月26日至6月1日净卖出6.01亿美元美国组合资产，连续两个月净卖出，为2023年6月以来首次。
  \item 台湾投资者同期净卖出1.49亿美元美元计价债券基金，卖出力度加大。
  \item 美国债券基金同期流入仅1.2亿美元，资金流入高度集中在股票而非债券。
  \item 标普500自ChatGPT问世以来累计上涨84\%，其中近三分之二来自盈利增长而非估值泡沫。
  \item MS测算2026年标普500每股盈利预测340美元，较ChatGPT发布前远期预测高出52\%。
  \item Conviction List新加入Casella Waste Systems、TPG、Tyson Foods和Block，核心逻辑是“盈利独立性”而非主题一致性。
  \item 美国消费者通胀担忧从1月低点53\%回升至59\%，但62\%计划夏季出行，世界杯成为新的消费催化剂，44\%消费者计划参与相关活动。
\end{itemize}
\begin{figure}[htbp]
\centering
\includegraphics[width=0.88\linewidth]{figures/fig_206.jpg}
\caption{Exhibit 1 - Exhibit 1: Conviction List - Directors' Cut May Update All price targets are 12-months; Prices are as of 5/31; green indicates new}
\end{figure}
\begin{figure}[htbp]
\centering
\includegraphics[width=0.88\linewidth]{figures/fig_207.jpg}
\caption{Exhibit 4 - Exhibit 4: CL stocks as of 5/31: YTD Estimate Revision vs. Price Returns}
\end{figure}
\begin{figure}[htbp]
\centering
\includegraphics[width=0.88\linewidth]{figures/fig_209.jpg}
\caption{Exhibit 7 - Exhibit 7: Conviction List - Directors' Cut Performance, presented as a portfolio Total Returns}
\end{figure}
\begin{figure}[htbp]
\centering
\includegraphics[width=0.88\linewidth]{figures/fig_210.jpg}
\caption{Exhibit 8 - Exhibit 8: Conviction List idiosyncratic returns have outpaced coverage idiosyncratic returns Cumulative idiosyncratic returns by rating or designation}
\end{figure}
\begin{figure}[htbp]
\centering
\includegraphics[width=0.88\linewidth]{figures/fig_315.jpg}
\caption{Exhibit 1 - Exhibit 1: World Cup Engagement by Demographic Categories May 2026: Expect to Engage with the World Cup (Among Total)}
\end{figure}
\begin{figure}[htbp]
\centering
\includegraphics[width=0.88\linewidth]{figures/fig_316.jpg}
\caption{Exhibit 2 - Exhibit 2: World Cup Engagement by Activity May 2026: World Cup Engagement (Among Total)}
\end{figure}
\begin{figure}[htbp]
\centering
\includegraphics[width=0.88\linewidth]{figures/fig_317.jpg}
\caption{Exhibit 3 - Exhibit 3: World Cup Spending By Activity May 2026: Expect to Spend More than Usual Because of World Cup (Among Consumers Who Expect to Engage)}
\end{figure}
\begin{figure}[htbp]
\centering
\includegraphics[width=0.88\linewidth]{figures/fig_318.jpg}
\caption{Exhibit 4 - Exhibit 4: Leisure Travel Plans by Type}
\end{figure}
{\small\color{refgray}\textbf{References:} [R006] 全球资金正在重新定价AI，但亚洲散户已经提前离场; [R019] 市场真正低估的不是AI，而是“非AI”资产的盈利韧性; [R024] 世界杯与夏季旅行正成为消费分化的新分水岭，而非总量信号}

\section{结语}
综合来看，市场当前定价的核心矛盾在于对供给约束的持续性普遍低估——无论是中东原油的恢复速度、AI芯片的产能瓶颈、还是中国地产的供给侧清洗。与此同时，全球资金行为的分化提示我们，在AI主题之外，盈利独立性和供给侧再定价正在成为下一阶段超额收益的关键来源。
\newpage
\section{报告覆盖清单}
\begin{itemize}
  \item [R001] 宏观与利率：通胀结构重塑，美元定价逻辑切换 - 通胀正在被三股力量重塑，市场低估了它们的持续性
  \item [R002] 全球贸易与供应链：补库信号已现，有效关税税率差异重塑格局 - 物流行业真正被低估的不是需求，而是定价权的结构性重塑
  \item [R003] 全球贸易与供应链：补库信号已现，有效关税税率差异重塑格局 - 市场真正低估的不是需求，而是供给侧的再定价
  \item [R004] 宏观与利率：通胀结构重塑，美元定价逻辑切换 - 市场真正低估的不是需求，而是美元定价的结构性断裂
  \item [R005] 宏观与利率：通胀结构重塑，美元定价逻辑切换 / 区域市场：亚洲央行鹰派转向，澳洲投资逻辑切换 - 亚洲能源冲击的真正赢家不是增长，而是央行的鹰派转向
  \item [R006] 全球贸易与供应链：补库信号已现，有效关税税率差异重塑格局 / 风险偏好与资金流向：机构与散户出现显著分歧 - 全球资金正在重新定价AI，但亚洲散户已经提前离场
  \item [R007] 未被主题摘要引用 - 中国金属市场正在讲述一个“供给故事”，而非“需求故事”
  \item [R008] 区域市场：亚洲央行鹰派转向，澳洲投资逻辑切换 - 印度制造业的真正信号：内需托底，但成本压力正在重塑利润分配格局
  \item [R009] 未被主题摘要引用 - 全球奢侈品定价的核心矛盾：市场低估的不是需求，而是供给侧的再定价
  \item [R010] 未被主题摘要引用 - 市场真正低估的不是油价，而是炼化结构壁垒的重新定价
  \item [R011] 工业科技：全球自动化周期出现“双顶”，中国以外复苏加速 - 中国软件行业4月增速放缓至8.9\%，但利润率改善揭示了一个被低估的结构性信号
  \item [R012] 中国地产与消费：销售分化加剧，储蓄率是消费真正拐点 - 市场真正低估的不是需求，而是供给侧的分化
  \item [R013] 中国地产与消费：销售分化加剧，储蓄率是消费真正拐点 - 市场真正低估的不是需求，而是供给侧的再定价
  \item [R014] 中国地产与消费：销售分化加剧，储蓄率是消费真正拐点 - 中国消费的真正拐点，不在于收入增速，而在于储蓄率何时开始下降
  \item [R015] 区域市场：亚洲央行鹰派转向，澳洲投资逻辑切换 / 工业科技：全球自动化周期出现“双顶”，中国以外复苏加速 - 全球工业技术周期的“双顶”正在确认，但市场对“中国需求之外的结构性加速”定价不足
  \item [R016] 中国地产与消费：销售分化加剧，储蓄率是消费真正拐点 - 中国车市真正的分化不在电动化，而在“谁的海外能扛”
  \item [R017] 宏观与利率：通胀结构重塑，美元定价逻辑切换 - 人民币中间价模型正在发出一个被忽视的定价信号
  \item [R018] 未被主题摘要引用 - 韩国出口的真相：AI 推高了价格，但没有推宽复苏
  \item [R019] 风险偏好与资金流向：机构与散户出现显著分歧 - 市场真正低估的不是AI，而是“非AI”资产的盈利韧性
  \item [R020] 未被主题摘要引用 - 自动驾驶的真正拐点，不在于技术，而在于运营的规模经济学
  \item [R021] AI与半导体：需求扩散至多节点，供给约束成为新定价变量 - 硬件市场真正稀缺的不是需求，而是被低估的供给议价权
  \item [R022] AI与半导体：需求扩散至多节点，供给约束成为新定价变量 - 市场真正低估的不是需求，而是供给侧的再定价
  \item [R023] 能源与大宗：供给恢复慢于预期，锂矿重启信号密集 - 市场低估的不是需求，而是供给纪律的自我强化
  \item [R024] 风险偏好与资金流向：机构与散户出现显著分歧 - 世界杯与夏季旅行正成为消费分化的新分水岭，而非总量信号
  \item [R025] AI与半导体：需求扩散至多节点，供给约束成为新定价变量 - 半导体CEO们的共识：AI需求正在重塑WFE定价权，但真正的约束是物理空间
  \item [R026] 医疗健康：双抗改写肺癌标准治疗，口服GLP-1竞赛参数被重新定义 - 肺癌治疗正在经历一场“分子分型”驱动的定价重构
  \item [R027] AI与半导体：需求扩散至多节点，供给约束成为新定价变量 - 市场低估的不是Broadcom的短期业绩，而是其2027年XPU份额的再定价空间
  \item [R028] 区域市场：亚洲央行鹰派转向，澳洲投资逻辑切换 - 市场真正低估的不是需求，而是供给侧的再定价
  \item [R029] 工业科技：全球自动化周期出现“双顶”，中国以外复苏加速 - 市场低估的不是零食需求，而是成本周期与渠道红利的共振
  \item [R030] 医疗健康：双抗改写肺癌标准治疗，口服GLP-1竞赛参数被重新定义 - 口服GLP-1的竞赛逻辑已经变了：中国资产正在定义下一轮竞争参数
  \item [R031] 全球贸易与供应链：补库信号已现，有效关税税率差异重塑格局 - 韩国出口的“分裂式繁荣”：半导体孤军深入，而结构性风险正在积累
  \item [R032] AI与半导体：需求扩散至多节点，供给约束成为新定价变量 - 韩国半导体出口的真相：市场低估了供给约束的持久性
  \item [R033] 能源与大宗：供给恢复慢于预期，锂矿重启信号密集 - 锂价反弹的真正考验：供给侧正在醒来，但市场低估了重启的节奏
  \item [R034] 能源与大宗：供给恢复慢于预期，锂矿重启信号密集 - 欧洲钢铁业正在经历一场被低估的供给侧重定价
  \item [R035] 能源与大宗：供给恢复慢于预期，锂矿重启信号密集 - 市场真正低估的不是需求，而是供给恢复的速度
  \item [R036] 医疗健康：双抗改写肺癌标准治疗，口服GLP-1竞赛参数被重新定义 - ASCO 2026最值得关注的信号：双抗在肺癌领域首次击败PD-1标准疗法
  \item [R037] 区域市场：亚洲央行鹰派转向，澳洲投资逻辑切换 - 日本IT服务业的真正拐点，藏在两家非覆盖公司的财报细节里
\end{itemize}
\section{逐篇报告摘录}
\subsection{[R001] 通胀正在被三股力量重塑，市场低估了它们的持续性}
{\small\color{refgray}归类：宏观与利率：通胀结构重塑，美元定价逻辑切换}

四月的通胀数据，表面上是数字的跳动，背后却是三个长期变量的共振。某外资投行最新发布的美国通胀监测报告，提供了一个值得产业决策者和资产配置者仔细拆解的分析框架。

核心PCE同比升至3.29\%，核心CPI同比升至2.74\%。单看这两个数字，似乎只是通胀从高点回落过程中的一次颠簸。但报告揭示的更深层问题是：推动四月通胀上行的三股力量——地缘冲突、AI基础设施需求、以及关税传导——各自有不同的持续性、传导路径和政策含义。市场对其中任何一股力量的定价，可能都不够充分。

这不是一个“通胀会不会再加速”的老问题。这是一个“通胀的结构正在被谁重新定义”的新问题。

1. 伊朗战争对通胀的冲击不只是油价，而是通过三条供应链同时传导

报告显示，油价从二月平均每桶71美元飙升至四月平均每桶117美元，涨幅超过50\%。但真正值得关注的不是油价本身，而是战争通过三条供应链同时推高通胀的机制。

第一条是能源直接传导。报告预计，更高的能源价格今年将直接推高核心PCE通胀约0.35个百分点。这是最直观、也最容易被市场定价的渠道。

第二条是非石油大宗商品。战争推高了海湾地区其他出口品的价格，包括氮肥和铝。报告预计，化肥成本上升将使食品价格今年上涨约1.0\%，对整体通胀贡献约0.07个百分点。更关键的是，食品价格上涨还会通过餐饮价格传导至核心PCE，额外贡献0.05至0.10个百分点。这是一个间接但更持久的传导路径。

第三条是供应链成本的重构。战争导致化肥、铝、甲醇等工业原料价格上涨，这些成本最终会嵌入制造业和服务业的定价中，其影响周期可能远超油价本身。

报告给出一个关键数据：自战争爆发以来，已将2026年12月同比整体PCE通胀预测上调了1.4个百分点。这意味着，地缘冲突不是通胀的短期扰动，而是正在成为通胀结构的一部分。

对于企业而言，这意味着需要重新评估供应链的“安全成本”是否已经永久性上升。对于投资者而言，这意味着能源和商品相关的通胀溢价可能不会随着战争结束而完全消失。

这是这份报告最值得反复阅读的部分。AI对通胀的影响，不是通过需求拉动，而是通过供给侧的三个具体渠道。

第一个渠道是电子元件价格。AI基础设施的强劲需求推高了关键电子元件的价格，包括数字存储芯片和电池。报告显示，电子元件及配件的生产者价格指数四月同比上涨了20\%。这是自2022年以来最快的增速。

第二个渠道是软件定价。报告指出，多家公司近期提高了软件价格，而软件和配件在PCE和CPI中都没有进行质量调整。这意味着，AI新功能带来的价值提升，在统计上被完全计为价格上涨，而非数量增加。这是一个潜在的测量偏差，但其对实际通胀的影响是真实的。

第三个渠道是电力需求。数据中心的强劲需求正在推高电力价格。报告估计，未来几年，电力价格上涨将使整体PCE通胀每年增加约0.1至0.2个百分点。这个数字看似不大，但其持续性值得关注——数据中心建设周期长，电力需求增长具有结构性特征。

报告特别强调了一个被市场忽视的细节：软件和配件在PCE中的权重约为CPI的30倍。这意味着，同样的软件价格上涨，对核心PCE的推升作用远大于对核心CPI。四月软件和配件价格环比上涨5.0\%，已是连续三个月大幅上涨。报告预计，未来一年，电子和数字投入品价格上涨将使核心PCE通胀额外增加0.3个百分点，到2026年三季度达到峰值0.6个百分点。

对于关注通胀指标的投资者而言，这意味着PCE和CPI的分化可能持续，而美联储更关注PCE。对于科技行业从业者而言，这意味着AI带来的成本压力正在从基础设施向软件和服务传导。

3. 关税的滞后效应尚未完全释放，核心商品通胀的回落可能比预期更慢

报告对关税效应的分析提供了一个重要的时间维度。数据显示，关税成本向消费者价格的传导率在13个月后达到90\%。

[... middle omitted ...]

心商品价格同比平均下降0.5\%。这意味着，当前商品通胀的“正常化”水平已经系统性上移。

报告预计，核心商品通胀将从四月的2.8\%降至2026年12月的1.8\%，再到2027年12月的0.5\%。但这一预测的前提是关税成本传导结束。如果贸易政策出现新的变化，这一回落路径可能被打断。

对于零售和制造业企业而言，这意味着成本端压力缓解的时间窗口可能比预期更短。对于关注通胀路径的投资者而言，核心商品通胀的回落速度是判断整体通胀走向的关键变量。

4. 通胀预期的分化正在形成，这对美联储的政策框架构成新挑战

通胀预期是决定实际通胀走向的重要因素。报告呈现了一个值得警惕的分化格局。

短期通胀预期明显上升。密歇根大学调查的一年期通胀预期升至4.8\%，纽约联储调查升至3.6\%。这反映了消费者对近期物价上涨的直接感受。

但长期通胀预期出现了显著分化。密歇根大学调查的5-10年通胀预期升至3.9\%，上升了0.4个百分点。而纽约联储的五年期预期则降至3.0\%，下降了0.1个百分点。

\subsection{[R002] 物流行业真正被低估的不是需求，而是定价权的结构性重塑}
{\small\color{refgray}归类：全球贸易与供应链：补库信号已现，有效关税税率差异重塑格局}

当市场还在为燃油成本上涨和电商增速放缓而担忧时，一份由某外资投行在2026年全球中国峰会期间发布的研报，揭示了一个被普遍忽视的底层变化：中国物流快递行业正在经历一场从“规模竞赛”到“定价权回归”的范式转换。四月的数据提供了最直接的证据——行业包裹量增速放缓至同比+3.2\%，但收入增速却加速至+6.3\%，单票收入恢复到7.65元，同比上涨3\%。这种“量缓、收增、价升”的组合，在过去三年的行业历史中极为罕见。

这不是一个短期的季节性波动，而是政策驱动、竞争格局重塑和头部企业战略选择共同作用的结果。市场对需求的关注过度集中在电商数据的月度波动上，却低估了供给侧正在发生的结构性变化——反内卷已经从一场“运动”演变为可重复执行的运营规则。

1. 反内卷的可持续性正在被低估，政策逻辑已经从“运动式”转向“制度化”

市场普遍担心反内卷政策是否只是权宜之计，但四月的数据和峰会交流给出了更清晰的答案。政策当局已经将价格稳定和基层网点福利纳入常态化监管框架，这意味着行业重回“价格战”的制度性障碍已经形成。

投行研报明确指出，反内卷措施正在从“运动”转向更可重复的运营规则，包括价格纪律、基层福利和网络稳定性三大支柱。中通快递管理层在峰会上表示，虽然无法保证2027年政策延续性，但政府持续关注基层福利和网络稳定，使得回归激烈价格战的可能性极低。

这意味着什么？行业竞争的核心变量正在从“谁更能烧钱”转向“谁更能守住价格”。对于投资者而言，判断行业格局的关键不再是猜测下一轮价格战的起点，而是评估各公司在价格趋同环境下的差异化能力。这从根本上改变了估值逻辑——从基于量的增长预期转向基于利润率的稳定性溢价。

2. 四月数据揭示了一个关键信号：价格弹性正在取代量价博弈成为行业主逻辑

四月快递行业数据是理解当前拐点的最佳窗口。行业包裹量16.84亿件，同比+3.2\%，这一增速在历史序列中并不突出。但更值得关注的是收入端的表现——行业总收入达到1289.1亿元，同比+6.3\%，增速显著快于量端。背后的驱动力是单票收入回升至7.65元，同比+3\%。

这种“量缓价升”的组合，在过去的行业周期中通常意味着需求疲软下的被动涨价，但这一次不同。它反映的是行业主动选择的质量优先策略——企业不再为了冲量而牺牲价格，而是在政策护航下，有底气对低质量订单说“不”。

韵达股份的案例最具说服力。四月韵达主动压缩超低价流量，导致包裹量同比下降4\%，但单票收入同比大涨10\%。这种“舍量保价”的行为，在过去的行业环境中几乎不可想象。它证明，头部企业已经将价格纪律内化为经营策略，而非仅仅是应对政策的被动选择。

对行业格局的含义是：市场份额的争夺正在从“量的竞争”转向“质的竞争”。能够维持价格纪律并同时提升服务质量的企业，将在新一轮竞争中占据优势。中通和圆通的表现印证了这一点——两者在四月均实现了量价齐升，尽管幅度不同。

3. 电商需求减速是真实的，但市场过度放大了其对快递定价的负面影响

四月国家统计局数据显示，线上零售额同比增速放缓至2.3\%，线上实物商品GMV仅增长0.2\%，创下2024年12月以来最低月度增速。家电和通讯器材品类尤为疲软，家电销售同比下滑15\%，通讯器材增速从27\%骤降至6\%。

这些数字无疑令人担忧。但研报的深层洞察在于：需求减速与价格回升正在同时发生，这恰恰说明行业定价权正在从需求端向供给端转移。过去，需求放缓几乎必然导致价格战，因为企业需要通过降价来争夺萎缩的蛋糕。但现在，政策约束和企业理性使得价格具有了“刚性”。

更关键的是，四月的数据显示，电商需求减速对快递企业的冲击是不对称的。头部企业通过服务质量和逆向物流能力，正在从结构性增长中获益。中通的逆向物流能力（峰值日处理1000万件）已经成为难以复制的竞争优势。韵达的逆向包裹

[... middle omitted ...]

燃油价格上涨对利润率的侵蚀。但峰会交流显示，头部企业并非被动承受成本压力，而是建立了有效的成本传导机制。

极兔速递在东南亚市场已经建立了成熟的成本传导机制，能够将增加的成本有效传递给客户，且客户接受度良好。韵达通过差异化的价格传导策略（根据不同地区和客户类型调整），基本抵消了燃油成本上涨的影响。满帮集团则观察到，油价上涨加速了货主向线上平台迁移，这反而支撑了一季度订单量。

更值得关注的是，企业正在通过技术和自动化来消化成本压力。中通在最后一公里持续推进自动化，直接分拣率达到40-45\%，无人车部署也在推进中。满帮将人工智能应用于优化货主和司机的工作流程，实现了成本节约和效率提升。

这些案例说明，燃油成本压力虽然真实，但头部企业已经将其转化为倒逼运营效率提升的动力。对于投资者而言，关键问题是：哪些企业具备足够的技术和管理能力，能够在成本压力下仍保持利润率稳定？这将是未来一年行业分化的核心驱动力。

5. 报告尚未完全回答的关键问题：定价权能否在电商平台压力下持续？

\subsection{[R003] 市场真正低估的不是需求，而是供给侧的再定价}
{\small\color{refgray}归类：全球贸易与供应链：补库信号已现，有效关税税率差异重塑格局}

过去几周，全球贸易数据释放了一个看似矛盾但实则高度一致的信息：从中国出发驶向美国的集装箱船量，周环比在下降，但同比依然保持两位数正增长。某外资投行最新一期关税影响追踪报告显示，截至5月28日当周，中国至美国的重箱船舶数量环比下降13\%，但同比仍高出17\%；同期TEU（标准箱）环比下降16.5\%，同比则增长20.8\%。

这一组数字很容易被解读为“贸易摩擦冲击正在消退”或“抢运高峰已过”。但真正值得决策者关注的，不是环比波动的方向，而是同比数据的韧性背后所揭示的结构性变化——全球供应链的再配置，正在从“被动应对关税”转向“主动重构布局”。这份报告最核心的判断并非关于短期进口量，而是关于一个正在成形的、以“有效关税税率差异”为核心驱动力的新贸易格局。

为什么这个判断现在很重要？因为市场目前的定价逻辑，仍然停留在“关税冲击-需求波动”的线性框架里。投资者对物流、运输和消费类资产的估值，主要依据的是季节性需求变化和库存周期的节奏。但报告中的高频数据表明，真正影响未来6-12个月贸易流量的变量，已经不再是关税本身的高低，而是不同国家之间有效关税税率的差异，以及由此引发的产能再分配。

这份报告提供的新信号是：中国对美出口的同比正增长，并非简单的“抢运”或者“低基数效应”。报告明确指出，目前的数据已经“超过了单纯对等关税纪念日的基数效应”，背后存在明显的“补库存”行为。而在补库存的同时，亚洲其他经济体对美出口的增速却明显落后于中国。这种分化，才是真正值得深挖的结构性线索。

报告中最被市场忽视的一个数据点，是亚洲不同经济体对美出口增速的分化。以截至5月28日当周的数据为例，中国对美重箱船舶数量同比增长17\%，而亚洲其他主要经济体（越南、韩国、台湾、日本）合计仅增长3\%。在TEU维度上，中国增长20\%，亚洲其他地区则下降4.5\%。

这种分化并非偶然。报告在讨论近中期影响时指出，一个关键变量是“那些现在拥有较低有效关税税率的国家，可能会增加对美出口”。但数据同时显示，这种替代效应尚未充分显现——至少目前，中国仍然在维持甚至扩大其对美出口的相对份额。

这里的含义是：市场此前普遍预期的“中国出口份额系统性下降”，可能被高估了。有效关税税率的变化并非均匀分布，不同商品、不同供应链环节的税率差异，导致企业并不简单地“从中国搬到越南”，而是根据税率差、产能成熟度和物流成本做更精细的权衡。报告没有完全展开的一个问题是：这种分化能否持续？如果中国通过供应链调整（比如通过第三国中转或调整产品分类）维持了相对较低的“实际关税负担”，那么“去中国化”的速度可能会比市场预期的慢得多。

对于投资者而言，这意味着单纯押注“越南受益、中国受损”的贸易替代逻辑，可能忽略了企业端更复杂的成本优化策略。那些能够帮助客户在复杂关税环境下灵活调整供应链的物流企业，反而可能获得更高的议价权。

2. 洛杉矶港的“双周脉冲”揭示的不是需求，而是库存管理的范式转换

报告提供了洛杉矶港未来两周的进口预测：截至6月5日当周，TEU环比增长5\%；两周后，预计环比增长38\%，同比增长46\%。如果只看这些数字，很容易得出“需求强劲”的结论。但报告的真正洞察在于：这种脉冲式增长，是“补库存”与“低基数”叠加的结果，而非终端消费需求的真实映射。

更值得关注的是报告中的另一组数据：西海岸铁路多式联运量同比增长14\%，卡车可用运力同比下降1\%，但卡车现货运价（不含燃油）同比上涨25\%。运价上涨的同时运力并未显著收紧，这说明价格上涨更多来自成本推动（尤其是燃油）和结构性运力错配，而非需求过热。

报告在讨论运价时做了一个关键区分：“我们需要将运价上涨的驱动因素，从产能限制和附加费中剥离出来，与进口需求和关税影响分开看待。”这是一个非常重要的分析框架。目前市场对运输股的定价，往往

[... middle omitted ...]

性”的回归

报告在讨论2026年展望时，提出了一个被很多市场参与者低估的观点：关税不确定性的消除，比关税本身的降低更重要。报告指出，2026年4月2日将迎来“解放日”一周年，届时“可能为托运人提供一个更一致的规则手册来制定计划”。

这句话的背后是一个深刻的逻辑：过去一年里，企业无法做中长期供应链规划的根本原因，不是关税水平有多高，而是关税政策的变化频率和不可预测性。报告引用了一项关键进展——美国最高法院的裁决以及随后根据第122条宣布的新关税——但强调，这些政策的不确定性“可能对中期至长期的货运规划产生影响”。

这里有一个尚未完全展开但至关重要的推论：一旦政策路径变得可预测，即使关税水平没有下降，企业的库存策略和供应链布局也会发生根本性变化。从“被动抢运”转向“主动规划”，意味着贸易流量将从脉冲式波动回归到更平滑的季节性节奏。这种变化对运输类资产的影响是深远的——波动率下降意味着运费溢价的消失，但也意味着企业可以更准确地匹配运力与需求，从而降低空载率和运营成本。

\subsection{[R004] 市场真正低估的不是需求，而是美元定价的结构性断裂}
{\small\color{refgray}归类：宏观与利率：通胀结构重塑，美元定价逻辑切换}

过去两年，全球宏观投资者一直在追问同一个问题：美国例外论（US exceptionalism）到底还在不在？2023到2024年，这个叙事一度瓦解，不是因为美国变弱了，而是因为世界其他地区追了上来。但进入2026年，美国经济数据再度领跑，企业盈利上修，美联储鹰派重定价，美股相对表现也在回暖——一切似乎都指向“例外论回来了”。

然而，某外资投行最新发布的宏观策略报告给出了一个反直觉的判断：即便美国经济数据全面占优，美元却很难因此走强。真正值得关注的，不是美国是否重新领先，而是市场定价美元的方式已经发生了结构性变化。这份报告的核心洞察在于，传统的“好数据=强美元”逻辑正在失效，而投资者如果沿用旧框架，可能会错过更重要的资产配置信号。

以下是我们基于该报告内容提炼的五个关键层次，帮助你理解为什么这次“例外论回归”的含金量远低于表面看起来的样子。

1. 美国经济重回G10前列，但美元定价的“锚”已从GDP切换到企业盈利

报告明确指出，美国GDP增长预期在2026年重新回到G10前三，并超出同行均值约0.75个百分点。与此一致，美国经济数据惊喜指数在近几个月大幅转正，与欧元区形成鲜明对比（图2）。同时，美国企业盈利修正也显著优于其他地区（图3）。这些数据放在一起，表面上构成了经典的美元看多组合。

但报告提出了一个关键区分：S\&P not GDP。也就是说，驱动全球资本流动的，从来不是宏观经济增速本身，而是企业盈利的边际变化。而恰恰在这个维度上，美国与世界其他地区的差距正在消失。报告中的图4显示，过去六个月，美国与世界其他地区的企业盈利增速都超过了20\%，几乎并驾齐驱。这意味着，全球股权投资者并没有足够理由因为“美国盈利更好”而加仓美元资产。

这一层的含义是清晰的：美元定价的锚正在从宏观叙事转向微观盈利。如果全球企业盈利增速趋同，那么美国GDP的领先优势就不足以单独支撑美元走强。投资者需要关注的不是美国经济是否“更好”，而是美国企业盈利是否“更独特”。

2. 全球资金流向正在逆转：2025年的“买爆”之后是“撤退”压力

报告提到了一个容易被忽视的结构性变量：全球投资者在2025年创纪录地买入了美国股票，而2026年的TIC数据已经显示，外资购买步伐明显放缓（图5）。这背后可能有两个驱动因素：一是2025年的买入过于集中，仓位已经拥挤；二是全球投资者开始寻求分散化，尤其是当其他市场的盈利增速赶上来之后。

这里需要特别注意的是，报告并没有断言外资已经开始净卖出，只是指出了“放缓”这一事实。但考虑到2025年的买入量级，即便只是边际放缓，也意味着美元资产的边际买家减少。对于美元汇率而言，资金流入的增量变化往往比存量更重要。如果外资购买意愿持续降温，美元将失去一个重要的需求支撑。

这一层的推导是：美国例外论在资金流上的体现，可能已经进入“利好出尽”阶段。2025年的创纪录买入本身就是对未来预期的提前定价，而2026年的放缓则暗示市场正在重新评估这一叙事的可持续性。

3. 利率优势的“降级”：美国已跌出G10收益率前三，且利差重定价未能推升美元

报告指出一个令人意外的数据：尽管美联储在近一个月被鹰派重定价，美国两年期收益率仍然低于G10中的澳大利亚、新西兰和挪威，排在前三名之外（图6）。这个位置是美国一年前丢失的，至今未能收复。更关键的是，报告中的图7显示，美债前端利率与G10其他国家的利差已经大幅走阔了约40个基点，但美元汇率几乎纹丝不动。

这说明什么？说明利率差异对美元的驱动力正在减弱。过去，利差扩大几乎必然伴随美元走强，但当前市场似乎更关注其他因素——比如全球资金再平衡、风险偏好变化、或者对美联储后续路径的怀疑。报告将此描述为“price action may signal something”，即市

[... middle omitted ...]

场可能正在对美元的长期信用进行重新定价。这种“信用折价”可能来自于市场对财政可持续性的担忧、对地缘政治风险的重新评估、或者对美元作为唯一储备资产地位的质疑。报告没有直接讨论这一点，但它的数据实际上为这个假设提供了间接证据——如果传统定价因子都指向美元走强，而美元不动，那么一定是有一个未被纳入模型的负面因子在起作用。

这个因子的具体内容，报告没有给出明确答案。但投资者在构建自己的观察框架时，有必要将“美元避险溢价的衰减”作为一个独立假设纳入考量。如果这个假设成立，那么当前美元的弱势就不是暂时的，而是结构性的。

报告在文末提出了一个具体的交易思路：做空欧元兑日元（sell EUR/JPY）。理由是，日本经济数据正在跑赢欧元区，但利率市场的重定价却走了相反的方向——欧洲央行被鹰派重定价，而日本央行显得更为鹰派（图8）。这种“数据与利率背离”的状态通常不可持续，要么数据向利率靠拢，要么利率向数据修正。报告倾向于认为，日本央行更鹰派的态度将逐渐被市场定价，从而推动日元走强。

\subsection{[R005] 亚洲能源冲击的真正赢家不是增长，而是央行的鹰派转向}
{\small\color{refgray}归类：宏观与利率：通胀结构重塑，美元定价逻辑切换 / 区域市场：亚洲央行鹰派转向，澳洲投资逻辑切换}

市场最初的担忧是清晰的：霍尔木兹海峡的长期关闭将重创依赖该地区的亚洲经济体。然而，最新一期投行研报揭示了一个反直觉的图景——亚洲（除中国外）通过多元化采购、动用战略储备和燃料替代，已将能源进口量迅速拉回正常水平。但这份韧性是有代价的。代价并非经济增长的显著放缓，而是通胀压力持续累积，以及随之而来的央行政策立场的根本性转变。韩国央行已发出明确的鹰派信号，预计最早在7月启动加息周期。这比市场此前预期的要早得多。真正值得关注的不是亚洲能否扛过能源冲击，而是“扛过”之后，货币政策正常化的速度将如何重塑资产定价、利差交易和资金流向。

报告的核心判断是：亚洲新兴市场正在经历从“增长韧性”到“通胀韧性”的切换，这意味着央行的关注点已从保增长转向控通胀。对于投资者而言，这意味着需要重新评估亚洲利率曲线的陡峭化程度、本币汇率的压力来源，以及那些依赖宽松流动性支撑的高估值板块的脆弱性。

1. 能源进口恢复的速度超出预期，但供给侧的切换能力是分化的关键

报告中最具说服力的数据来自油轮进口的实时追踪。EMAX（新兴亚洲除中国外）和印度的5月能源进口量已快速逼近正常水平。这并非因为需求萎缩，而是因为采购来源的迅速切换、国家战略储备的释放，以及向煤炭和生物燃料等替代能源的转移。

这一事实意味着，市场此前对“供给中断导致经济断崖”的极端情景定价可能过度。但需要指出的是，这种切换能力并非均匀分布。大型炼油企业和拥有多元化供应链的国家（如韩国、印度）具备更强的缓冲能力，而小型炼厂和依赖单一来源的经济体（如部分东盟国家）可能仍面临压力。报告没有完全展开的是，这种切换成本——包括物流溢价、合同违约风险和替代燃料的额外排放成本——可能在未来几个季度逐步反映在工业品出厂价格中。

对于投资者而言，这意味着：能源进口量的恢复是短期利好，但切换成本的结构性上升可能成为下一阶段通胀的新来源。那些能够将成本转嫁到下游的行业（如半导体、化工）将相对受益，而成本敏感且竞争激烈的行业（如低端制造）可能面临利润率侵蚀。

报告中最具政策含义的观察来自韩国央行。尽管该行在最新会议上维持利率不变，但其前瞻指引发生了实质性变化——信号指向最早7月加息，且后续加息可能贯穿明年。报告将这一转变归因于韩国持续的增长动能（尤其是科技板块）以及对该动能将溢出至需求侧通胀的担忧。

这一信号的意义在于：韩国是全球科技周期的风向标，其央行的鹰派转向意味着，科技出口的强劲复苏正在转化为内需通胀压力。这不是一个孤立事件。菲律宾和印尼央行已先行加息，尽管驱动因素不同（菲律宾是通胀冲击，印尼是汇率压力）。报告暗示，一个“鹰派阵营”正在亚洲形成。

对于全球投资者而言，这意味着：亚洲利率曲线的定价需要系统性上调。此前市场对亚洲央行“跟随美联储”的预期可能需要修正——亚洲央行现在可能因为自身通胀压力而主动加息。这将影响利差交易策略，并可能推动资金从亚洲高收益债市场向发达国家回流，至少在加息初期如此。

3. 印度央行是关键的“例外”，但这份耐心能持续多久是最大的未解问题

报告明确指出，印度央行是当前亚洲主要央行中预计将保持耐心的一个。其逻辑是：印度央行行长的表态是，只有在能源价格压力溢出至核心通胀时才会加息，而目前核心通胀非常温和。报告因此预期印度央行将在下周的会议上按兵不动。

这一判断值得关注，但需要警惕。印度的通胀结构有其特殊性：食品和能源在CPI中权重较高，而核心通胀确实相对温和。然而，报告没有完全回答的问题是：如果能源冲击持续，印度卢比面临贬值压力，央行是否还能保持耐心？历史经验表明，新兴市场央行的独立性往往在汇率压力面前受到考验。此外，印度经济增长本身非常强劲（报告预测2025年GDP增长7.5\%），这种增长动能本身就可能逐步推高核心通胀。

对于投资者而言，印度卢比和印度债券

[... middle omitted ...]

自政策节奏的“主动刹车”。报告认为，这意味着上半年的疲软将减轻下半年的基数效应压力，即“预期中的下半年回摆不会那么剧烈”。同时，4月出口的强劲反弹（集中在记忆芯片、电动车、太阳能和电池等高端品类）为经济提供了支撑。

对于投资者而言，这意味着：中国经济的短期叙事可能从“政策刺激”转向“内生修复”。出口的韧性是关键变量，但其集中度在高端品类也增加了脆弱性——如果全球科技周期或贸易政策出现变化，这些板块的波动性可能放大。5月PMI数据（报告预测NBS制造业PMI为50.1，较4月微降）将是检验这一判断的近期窗口。

5. 报告尚未完全回答的三个关键问题，决定了投资组合的再平衡方向

尽管报告提供了丰富的洞察，但仍有几个关键问题悬而未决，这些正是投资者需要持续跟踪的焦点。

第一个问题是：亚洲央行的鹰派转向是否已充分定价？韩国央行预计在一年内加息100个基点，但市场利率曲线的定价是否已完全反映这一路径？如果实际加息节奏更快，短端利率的上升可能超出预期，进而压制风险资产估值。

\subsection{[R006] 全球资金正在重新定价AI，但亚洲散户已经提前离场}
{\small\color{refgray}归类：全球贸易与供应链：补库信号已现，有效关税税率差异重塑格局 / 风险偏好与资金流向：机构与散户出现显著分歧}

这份来自某外资投行5月底的资金流周报，揭示了一个被市场喧嚣掩盖的结构性信号：全球资金正在以近半年来最快的速度涌入美国股票基金，但与此同时，韩国和台湾的散户投资者却在连续两个月净卖出美国资产。这不是一个简单的“风险偏好回升”故事，而是一个关于信息分层、预期差和资金行为分化的关键窗口。

报告显示，截至5月29日当周，外国资金净流入美国股票基金高达27.25亿美元，是前一周13.54亿美元的两倍有余。5月全月，流入美国股票基金的外资总额达到86亿美元，创下2026年1月以来的最高单月纪录。与此同时，新兴市场股票ETF的资金流入却从4月的58.32亿美元骤降至5月的24.18亿美元。

但真正值得注意的，不是这些总量数字，而是两个亚洲市场的反向操作。韩国散户投资者在5月26日至6月1日期间净卖出6.01亿美元的美国组合资产，这是自2023年6月以来首次出现连续两个月净卖出。台湾投资者同期也净卖出1.49亿美元的美元计价债券基金。

这意味着什么？全球最活跃的两批亚洲散户，正在用真金白银告诉市场：他们对美国资产的态度正在发生转折。而机构资金却在加速涌入。这种分歧本身，就是最重要的信号。

报告的数据最引人注目的部分，是流入美国股票基金的外资节奏。5月全月86亿美元的净流入，不仅显著高于4月的30.94亿美元，更是在3月净流出1.18亿美元的背景下实现的强势反转。从周度数据看，加速趋势同样明显：5月最后一周的27.25亿美元，几乎是前一周13.54亿美元的两倍。

但需要警惕的是，这种流入并非均匀分布。同一时期，美国债券基金的流入仅有1.2亿美元，虽然高于前一周的2.52亿美元，但与股票基金的规模完全不在一个量级。5月全月债券基金流入也仅为6.88亿美元，远低于股票基金的86亿美元。

这意味着，当前外资对美国资产的偏好高度集中在权益端。这不是一次全面的“美国资产再配置”，而是一次有明确指向性的“AI主题交易”。报告虽然没有直接点名AI，但从时间窗口看，5月正是全球AI产业链新一轮叙事发酵的时期。资金不是在买“美国”，而是在买“美国的AI”。

这种结构性的集中度，本身就暗藏风险。一旦AI叙事出现任何扰动，这些快速涌入的资金也可能以同样快的速度撤离。3月美国股票基金净流出1.18亿美元就是前车之鉴。

与美股基金的火爆形成鲜明对比，新兴市场ETF的资金流入正在快速降温。5月全月，外资净买入新兴市场ETF仅24.63亿美元，远低于4月的62.85亿美元。更值得注意的是，流入结构发生了根本性变化：4月新兴市场股票ETF和债券ETF分别流入58.32亿美元和4.53亿美元，而5月股票ETF流入骤降至24.18亿美元，债券ETF更是几乎停滞，仅有0.44亿美元。

这不是简单的“风险偏好下降”。如果全球投资者真的在全面避险，他们不会同时加速买入美国股票基金。更合理的解释是：资金正在从新兴市场抽离，集中配置到美国——尤其是美国的AI相关资产。

对于依赖外资流入的新兴市场而言，这意味着两个层面的压力。第一，流动性层面，外资的边际增量正在减少。第二，估值层面，当全球最活跃的资金都在追逐同一个主题时，被冷落的市场面临的不仅是资金流出，更是相对估值的系统性下修。

但报告没有完全展开的是，新兴市场内部也在分化。中国、印度、东南亚等不同区域在AI产业链中的位置不同，资金流出的程度和节奏也会不同。这部分仍需要更细颗粒度的数据来验证。

报告中最具“逆势”意味的数据，来自韩国散户。5月26日至6月1日当周，韩国散户净卖出6.01亿美元的美国组合资产，其中美股5.01亿美元、美债1.01亿美元。5月全月净卖出5.68亿美元，4月净卖出9.1亿美元。这是自2023年6月以来，韩国散户首次连续两个月净卖出美国资产。

回顾历史数据，韩国散户在202

[... middle omitted ...]

国散户的净卖出发生在全球机构资金加速流入的背景下。这意味着，机构正在接盘散户的筹码。这本身不是一个危险的信号——机构资金通常比散户更“长线”。但问题在于，当机构资金流入的节奏放缓，而散户的卖出意愿还在持续，市场的边际买家就会减少。

台湾投资者的行为与韩国散户有相似之处，但标的和节奏不同。5月23日至29日当周，台湾投资者净卖出1.49亿美元的美元计价债券基金，高于前一周的1.06亿美元。不过，从月度数据看，5月全月台湾投资者净买入1.86亿美元，部分抵消了4月3.42亿美元的净卖出。

这意味着，台湾投资者的行为更像是一次“战术性调整”，而非趋势性转向。他们在4月大幅减持后，5月又重新小幅建仓。这与韩国散户的“连续两个月净卖出”形成对比。

背后的原因可能在于，台湾投资者对美元债券的配置更多是基于“息差交易”逻辑，而非单纯的资本利得预期。当美债收益率波动加大时，他们会进行仓位调整，但不会完全退出。而韩国散户对美股的交易更多是基于“动量”，一旦趋势转弱，离场意愿更强。

\subsection{[R007] 中国金属市场正在讲述一个“供给故事”，而非“需求故事”}
{\small\color{refgray}归类：未被主题摘要引用}

过去六周，市场目光几乎全部聚焦在中国需求端——PMI回落至50.0，固定资产投资中的基建部分同比下滑18.6\%，出口新订单指数跌破48.6。这些数据似乎都在指向同一个叙事：中国金属消费正在经历疫后复苏的阶段性回落。

这份来自某外资投行的周度高频库存追踪报告，却给出了一个更值得深思的信号。铜的周度去库速度确实在放缓，从三、四月的异常强劲回归到季节性正常水平。铝的去库刚刚启动，节奏甚至略慢于历史同期。锌的库存横盘不动。如果只看这些，结论似乎与市场共识一致。

但真正打破共识的，是隐藏在库存数字背后的另一组数据。全球铁矿石周度发运量环比飙升约19\%，同比同样接近19\%的增幅，其中澳大利亚环比增长20\%，巴西环比增长30\%。与此同时，中国港口的铁矿石库存却纹丝不动，维持在1.71亿吨的绝对高位。

这意味着什么？供给端正在发生比需求端更剧烈的结构性变化。而这份报告真正值得读透的部分，不是“需求好不好”，而是“供给正在如何被重塑”。

1. 铜的“弱去库”不是需求崩塌的信号，而是回归正常化的确认

铜的周度去库从3-4月的异常高位（单周去库超过40kt）回落至上周的约3kt，乍看之下像是消费动能大幅衰减。但将时间轴拉长到五年均值来看，当前的去库节奏几乎完全贴合季节性。

更关键的是绝对库存水平。中国铜总库存目前约220kt，仍处于历史季节性区间的底部。即便相比去年同期高出约60kt，但去年的低库存是特殊事件驱动的——当时全球铜贸易流因美国关税预期发生扭曲，大量铜被提前拉入美国，导致中国到货量偏低。

所以，铜的库存故事不是“需求变差”，而是“需求从异常强劲回到了正常”。这本身算不上利好，但也不构成利空。真正需要警惕的，是如果这种“正常化”持续到传统消费旺季结束后，库存是否会出现反季节性的累积。报告目前尚未对此给出明确判断，这恰恰是值得持续追踪的变量。

对投资者的含义：铜价的下行风险更多来自宏观预期转弱，而非基本面已经恶化。当前库存结构不支持趋势性做空，但也不支持追高。

铝的数据与铜形成鲜明对比。中国铝库存总量高达1.4Mt，处于2019年以来同期的最高水平。尽管上周出现了约16kt的去库，但绝对库存水位仍比五年均值高出约40\%。

市场通常将高库存解读为需求疲弱。但这份报告揭示了一个更复杂的图景：铝的高库存并非单纯因为消费不足，而是因为供给端在过去一年经历了产能释放的集中期。云南水电铝产能的逐步恢复、以及部分新增合规产能的投产，共同推高了社会库存的基数。

更值得注意的是，去库已经启动。虽然节奏略慢于季节性，但方向是明确的。这暗示着需求端并非毫无承接力，只是消化存量需要时间。如果后续铝的去库速度加快，市场对铝板块的定价逻辑可能从“库存压力”转向“成本支撑+供给约束”。

对投资者的含义：铝的库存压力是已知信息，但去库拐点的确认时间窗口才是真正的交易信号。当前估值中已经包含了悲观预期，超预期的去库节奏可能带来修复机会。

这是整份报告中信息密度最高的一个片段。全球铁矿石发运量单周暴增约19\%，但中国港口库存毫无变化。与此同时，运费在同步飙升——从图2可以看出，图巴朗-青岛航线运费从2月底的26美元/吨已经上涨到37美元/吨，涨幅超过40\%；黑德兰-青岛运费同期从9美元涨到16美元。

这三个数据放在一起，指向一个被忽略的传导链：矿山在增加发货，但中国钢厂在谨慎接货；运费的上涨正在侵蚀矿山的FOB利润空间，而港存的不增反稳说明实际到港量并未完全转化为可用库存。

更深一层的含义是：铁矿石的CFR价格虽然维持在110美元/吨附近的高位，但FOB价格（澳矿约95美元/吨）已经基本持平于2月中旬的水平。这意味着，矿山并没有享受到名义价格上涨的全部红利，大部分被运费吃掉了。

对投资者的含义：铁矿石的供给故事正在从“量”转向“

[... middle omitted ...]

。这是一种明确的政策引导：鼓励行业整合，通过并购而非新建来优化产能结构。

第三，氢基冶金和电炉等绿色高端项目允许1:1置换，不锈钢中频炉首次被纳入监管。

Mysteel的测算结论是：在1.5:1的全国统一减量置换叠加僵尸产能清退后，2027年之后中国钢铁行业将基本没有新增净产能，转而进入净削减阶段。供需平衡将收紧，为钢价提供更强的底部支撑，行业盈利稳定性有望改善，低价格内卷的问题将得到缓解。

这份报告的判断如果成立，将是过去五年中国钢铁行业最重要的供给侧改革。它不是在“限制”产能，而是在“重置”产能的定价逻辑。过去市场对钢铁股的估值长期受困于“产能过剩-价格战-利润微薄”的负向循环，而新规试图从制度层面打断这个循环。

但这里有一个报告尚未完全回答的关键问题：执行力度和过渡期安排。新规设有两年过渡期，跨企业、跨集团产能指标转移在此期间仍被允许。这意味着，2027年之前的实际产能出清速度取决于企业的主动淘汰意愿和地方政府的执行力度。报告给出了方向，但没有给出概率。

\subsection{[R008] 印度制造业的真正信号：内需托底，但成本压力正在重塑利润分配格局}
{\small\color{refgray}归类：区域市场：亚洲央行鹰派转向，澳洲投资逻辑切换}

印度制造业的真正信号：内需托底，但成本压力正在重塑利润分配格局

5月的印度制造业PMI数据，看起来是一个温和的改善。从54.7上升到55.0，连续第三个月扩张，产出指数更是跳升至三个月高点58.0。如果只看这个数字，很容易得出“印度制造业景气度持续向好”的结论。但这份来自某外资投行的研报，真正值得关注的不是PMI本身，而是藏在细项里的结构性分化：内需强劲，但外需走弱；产出扩张，但成本压力仍处四年高位。这组矛盾信号指向一个更本质的追问——当前的增长，有多少是可持续的，又有多少是在消耗企业的利润空间？

这份报告的核心判断是：印度制造业正在经历一轮由国内需求驱动的复苏，但复苏的质量正在被上游成本侵蚀。对于投资者和产业决策者而言，真正需要跟踪的，不是PMI的绝对值，而是企业能否将成本压力传导至下游——这决定了利润分配的方向，也决定了哪些行业、哪些公司能在下一阶段跑赢。

5月PMI的改善，几乎完全由国内需求拉动。产出指数从56.9升至58.0，受访企业明确将增长归因于“具有韧性的国内需求”。更细致地看，中间品和资本品行业的新订单增速最快，消费品则出现了放缓。这个结构本身就透露了关键信息：资本品和中间品的扩张，通常意味着下游企业正在加大投资或备货，指向的是对未来需求的乐观预期，而非当下消费的即时反馈。

但硬币的另一面是，新出口订单指数从4月的56.3骤降至53.7，创下三个月新低。这并不是一个孤立的下滑——在过去的几个季度里，印度出口订单的波动一直大于内需。全球贸易放缓、主要经济体需求降温，正在对印度出口形成实质性压制。这意味着，如果内需出现任何超预期的走弱，制造业将缺乏外部缓冲。

对于读者而言，这个分化带来的观察框架是：接下来需要重点追踪资本品订单的持续性。如果资本品扩张只是短期补库行为，而非趋势性投资周期启动，那么整个制造业的景气度将面临回调风险。

PMI数据中最容易被忽视，但可能最具长期含义的，是价格分项。5月投入价格指数从56.4微降至56.2，产出价格指数从53.8降至53.0。表面上看，通胀压力似乎有所缓和。但报告明确指出，投入价格仍处于“近四年高位”附近，企业普遍反映能源、燃料、原材料和运输成本在上升。

这里有一个重要的结构性细节：资本品行业面临最高的投入成本通胀，消费品行业则最低。这意味着，不同行业的企业，承受的成本压力完全不同。资本品企业——通常是机械、设备、工程类公司——对上游原材料价格最为敏感，而消费品企业因为有更强的品牌溢价和终端定价能力，成本传导相对顺畅。

但更关键的问题是，产出价格指数在下降。这说明，尽管成本高企，企业并没有（或无法）将全部压力转嫁给下游。投入价格与产出价格之间的“剪刀差”在收窄，但并未消失。对于制造业企业来说，这意味着利润空间仍在被压缩。

这个信号对资产定价的含义是：当成本压力持续而定价能力不足时，市场将开始区分“有定价权的企业”和“被动接受价格的企业”。前者的利润率韧性更强，估值应该获得溢价；后者则可能面临盈利下修风险。报告没有展开这一层，但这是投资者必须自己做的功课。

5月就业指数从4月的高位回落。报告没有给出具体数值，但趋势明确。这可能是整个PMI数据中最容易被低估的信号。

在PMI的五大分项中，就业通常是最滞后的指标。企业通常在确认需求可持续后才会增加招聘。因此，就业指数的回落，可能意味着企业对未来需求的信心出现了边际动摇。结合出口订单的下滑，这个信号值得警惕。

但也要注意，单月的数据波动并不构成趋势。印度制造业近年来经历了多次就业数据的短期波动，最终都回到了扩张区间。关键在于，接下来几个月就业指数是否会持续低于50的荣枯线。如果出现连续收缩，那就不是波动，而是需求走弱的确认。

对于产业决策者而言，就业数据是观察“增长是否惠及劳动力市场”的核心

[... middle omitted ...]

集中审批和采购，后续可能放缓。

报告没有提供资本品订单的同比增速或与历史周期的对比，也没有讨论政策刺激的边际递减效应。这是读者在阅读完整报告时需要自行追问的问题：如果剔除政策因素，内生性的资本开支需求究竟有多强？

这个未解问题，恰恰是判断印度制造业能否从“复苏”走向“扩张”的关键。如果资本品扩张是可持续的，那么整个制造业的景气周期将更长；如果只是政策脉冲，那么市场可能高估了未来几个季度的增长动能。

综合以上分析，对于关注印度市场的读者，可以建立一个简单的观察框架，用来评估制造业的真实健康状况，而不仅仅是PMI的读数。

第一层，看需求结构。内需和出口的比值是否在恶化。如果出口持续走弱，内需必须加速才能维持总量增长。但内需本身也有结构问题——消费品放缓意味着终端消费可能已经见顶。

第二层，看价格传导。投入价格与产出价格的差值，是衡量企业利润压力的直接指标。如果这个剪刀差持续扩大，意味着成本无法转嫁，企业盈利将承压。反之，如果产出价格开始回升，说明定价权正在回归。

\subsection{[R009] 全球奢侈品定价的核心矛盾：市场低估的不是需求，而是供给侧的再定价}
{\small\color{refgray}归类：未被主题摘要引用}

全球奢侈品定价的核心矛盾：市场低估的不是需求，而是供给侧的再定价

过去二十年，全球奢侈品行业以约8\%的年复合增长率跑赢全球GDP三倍，投资者涌入这个赛道，本质上是在为“持续增长”支付溢价。但2025年以来，全球奢侈品板块估值大幅回撤：以NTM EV/Sales计算，板块交易在3.7倍，较2025年均值折价18\%，较十年均值折价9\%。与此同时，行业ROIC仍维持在15.0\%的十年均值之上。某外资投行最新研报指出，当前市场最大的误判，并非对中国需求的过度悲观，而是对奢侈品行业“固定成本结构”下经营杠杆的重新定价尚未充分展开。

这份研报的核心贡献，是搭建了一个从收入增长到估值扩张的完整传导链条：收入超预期驱动经营杠杆，经营杠杆推动利润率扩张，利润率变化直接映射为EV/Sales倍数变动，最终通过ROIC改善传导至股东总回报。在这个框架下，当前估值与基本面的背离，恰恰指向了一个被忽视的投资窗口。

1. 奢侈品行业真正的护城河，藏在70\%-80\%的毛利率和固定成本结构里

研报揭示了一个常被误解的事实：奢侈品行业的竞争优势，并非来自品牌溢价本身，而是来自其独特的成本结构。行业平均毛利率高达70\%-80\%，因为原材料成本仅占零售价的极小部分。而运营支出的主体是零售门店租金、销售人员和营销费用——这些几乎全部是固定成本。

这意味着什么？规模本身就是最深的护城河。大型品牌可以用更低的销售费用率覆盖同样的固定成本，而小型品牌却要在同样昂贵的地段、同样稀缺的客户资源和人才市场上与巨头竞争。这种“规模驱动效率，效率强化竞争力”的正向循环，是奢侈品行业集中度持续提升的根本动力。

从研报提供的2025年行业成本结构来看，EBIT利润率平均为21\%，而运营支出占收入的46\%。在这样一个固定成本占比接近一半的行业里，收入增速的微小变化，都会通过经营杠杆被显著放大。

2. 收入增速的“超预期”才是股价的终极驱动力，而非利润本身

研报用一张长达17个季度的矩阵表格，清晰地揭示了奢侈品行业股价运行的核心规律：股价反应的是“收入增速是否超预期”，而不是利润是否增长。在统计的17个季度中，当公司季度有机收入增速超预期时，次日股价上涨的案例占比显著高于下跌；反之亦然。

这一发现挑战了传统投资框架中“利润为王”的直觉。但逻辑链条其实非常清晰：在固定成本结构下，收入超预期意味着经营杠杆的释放，而经营杠杆的释放会直接推动利润率扩张。更重要的是，奢侈品行业的利润表结构决定了，EBIT以下的项目（财务费用、税收）相对稳定，因此EBIT的变化几乎可以完整传递到净利润。这意味着，收入增速的“惊喜”或“失望”会以接近1:1的比例映射到利润端。

对投资者而言，这意味着盯住季度收入增速与市场预期的差距，比分析利润表细节更具前瞻性。

3. EBIT利润率每变动100个基点，估值倍数将同向变动约0.2倍

研报进一步量化了利润率与估值之间的传导关系：EBIT利润率的年度变化与EV/Sales倍数的变化之间存在稳定的正相关。从2006年至2025年的历史数据来看，EBIT利润率每扩张100个基点，EV/Sales倍数平均提升约0.2倍。

这个关系的稳定性来自于奢侈品行业利润结构的可预测性。由于财务费用和税率相对固定，EBIT的变化能够“干净地”传导至净利润，进而影响P/E估值。因此，有机收入增速通过影响利润率，最终成为估值倍数变化的最有效代理变量。

当前行业EV/Sales倍数处于十年均值以下9\%的水平，而ROIC仍高于十年均值。这意味着，如果收入增速能够企稳甚至回升，估值修复的空间将是可观的。

研报还建立了一个更长期的观察框架：ROIC的变化是预测中期股东总回报最有效的指标之一。这个关系在行业上行期和下行期都保持稳定。背后的逻辑是，ROIC的提升通常

[... middle omitted ...]

经济进一步走弱，那么ROIC的支撑基础就会动摇。

5. 报告尚未完全回答的关键问题：中国需求的“结构性”还是“周期性”？

研报坦承，当前市场最大的分歧在于：中国奢侈品消费放缓是结构性还是周期性的？中国房地产危机和消费者信心崩塌导致行业增速显著放缓。研报的历史视角提供了一种安慰：过去30年，每当行业增速放缓，投资者都会担心“这是否是奢侈品的终点”，但每一次都不是。

然而，这一历史规律是否适用于当前的中国市场？中国奢侈品消费者的行为模式正在发生变化：从“炫耀性消费”向“体验式消费”迁移，从“品牌崇拜”向“性价比敏感”转变。这些变化究竟是经济周期导致的暂时调整，还是消费价值观的长期重塑？研报没有给出确定的答案，而是将这个问题留给了读者。

这恰恰是当前市场定价分歧的核心。如果中国需求是周期性的，那么当前估值折价就是过度反应；如果是结构性的，那么奢侈品行业的长期增长中枢可能下移，估值折价将具有合理性。

6. 投资者的观察框架：从“增长故事”转向“防御+自愈”的双轨配置

\subsection{[R010] 市场真正低估的不是油价，而是炼化结构壁垒的重新定价}
{\small\color{refgray}归类：未被主题摘要引用}

当多数投资者还在关注原油价格波动对炼化企业的影响时，一份来自某外资投行的亚太精选名单更新，给出了一个值得深思的信号：该机构将一家韩国炼油企业——S-Oil——新增至其“Conviction List”（确信买入名单），并给出了37\%的上行空间。

这份报告的核心判断并非关于油价走势，而是关于一个被市场系统性低估的结构性变量：中间馏分（柴油与航空煤油）的供应紧张正在重塑炼化行业的利润分配格局，而拥有特定产品结构优势的企业，将在这一轮“炼化超级周期”中实现显著的现金流拐点。

这不是一个短期的交易信号。报告所指向的，是一个可能持续至2027年甚至更长时间的价值重估过程。对于关注亚太权益市场的投资者而言，理解这一判断背后的逻辑，比追逐任何单一标的更为重要。

1. 这份名单揭示了一个超越地域和行业的共同主题：从“增长”到“价值实现”的叙事切换

该投行的亚太精选名单覆盖了24只股票，横跨消费、金融、医疗、工业、科技、自然资源等多个板块。表面上看，这是一份分散化的组合推荐。但如果深入分析其构建逻辑，会发现一个清晰的底层框架。

报告将这些股票归入六大主题：转型、增长+价值、利润率扩张、全球布局、拐点。这并非简单的标签分类，而是一种投资哲学的表达。在当前的宏观环境下，纯粹依赖收入增长的叙事正在让位于“如何将规模、技术或品牌优势转化为可持续的现金流和股东回报”。

以新加入的S-Oil为例，其入选理由并非产量增长，而是自由现金流的急剧拐点——报告预计其2027年自由现金流收益率将达到18\%。同样，名单中的其他公司，如日本的安川电机（Yaskawa Electric）和中国的汇川技术（Inovance），其核心逻辑也是AI资本开支与自动化复苏带来的利润率上行，而非简单的营收扩张。

这意味着，投资者需要调整其评估框架：在利率环境维持高位、增长稀缺的背景下，市场愿意为“确定性”和“转化效率”支付溢价。 那些能够证明自己正在从“投入期”进入“收获期”的企业，将获得更积极的估值重估。

2. 炼化行业的结构性紧张并非周期性波动，而是供给侧的“硬约束”正在形成

S-Oil的核心投资逻辑，根植于某投行对全球炼化行业的一个深层判断：中间馏分（柴油和航空煤油）正在成为炼化系统中最关键的瓶颈，且这一瓶颈在中期内难以消除。

报告指出，S-Oil约56\%的产品产出集中在中间馏分，这种产品结构使其在市场中占据有利位置。而这一判断的背景是：全球范围内，新增炼化产能的投产周期漫长且成本高昂，同时，环保法规和能源转型压力正在加速老旧炼厂的退出。这意味着，即便需求增速放缓，供给侧的收缩速度可能更快，从而导致产品价差（crack spread）维持在高于历史均值的水平。

这里需要特别注意的是，报告并未假设需求端的爆发式增长。其逻辑的坚实之处在于：它押注的是供给端结构性收紧，而非需求端的周期性繁荣。 这是两种截然不同的投资叙事。前者具有更强的可持续性和可预测性，而后者往往伴随着较大的波动风险。

对于投资者而言，这意味着需要重新审视炼化板块的定价逻辑。如果市场仍然用传统的周期性框架来估值这些公司，那么当自由现金流持续超预期时，估值修复的空间可能相当可观。

3. “超级周期”的验证节点在于现金流能否兑现，而不仅仅是利润表

报告对S-Oil最引人注目的预测，是2027年18\%的自由现金流收益率。这一数字如果实现，意味着该公司的估值将从当前的周期性低谷，跃升至一个完全不同的水平。

这一判断的关键支撑在于资本开支的显著下降。S-Oil正在完成一项约70亿美元的产能扩张项目，预计于2026年年中完工。资本开支高峰过后，自由现金流将出现“断崖式”改善。这是典型的“投入期结束、收获期开始”的叙事。

然而，这里存在一个需要投资者自行判断的风险点：新产能——即名

[... middle omitted ...]

，但自由现金流的真正改善，才是估值重估的最终确认。 投资者需要关注的不仅是季度盈利，更是资本开支节奏、营运资本管理以及股东回报政策（如分红和回购）的实际执行情况。

4. 报告尚未完全回答的关键问题：当“炼化超级周期”遭遇“能源转型”的长期叙事

尽管报告对S-Oil的中期前景充满信心，但有一个关键问题并未得到充分展开：“炼化超级周期”的可持续性，与全球能源转型的长期趋势之间，存在怎样的张力？

如果能源转型加速推进，电动汽车渗透率提升导致汽油需求见顶，航空业的脱碳努力（如可持续航空燃料SAF）取得突破，那么对传统炼化产品的长期需求将面临结构性下行。报告所依赖的“供给紧张”逻辑，在长期可能被“需求萎缩”所削弱。

当然，报告的隐含假设可能是：转型过程将是渐进的，且中间馏分的需求韧性更强（尤其是在航空和货运领域）。但这一假设的验证，需要更长时间的数据积累。对于投资者而言，这意味着需要将S-Oil的投资视为一个“中期机会”，并密切关注能源政策、技术突破和消费者行为的变化。

\subsection{[R011] 中国软件行业4月增速放缓至8.9\%，但利润率改善揭示了一个被低估的结构性信号}
{\small\color{refgray}归类：工业科技：全球自动化周期出现“双顶”，中国以外复苏加速}

中国软件行业4月增速放缓至8.9\%，但利润率改善揭示了一个被低估的结构性信号

4月中国软件行业收入同比增长8.9\%，较3月的11.6\%明显回落。这个数字本身并不意外——季节性波动、PMI走弱、AI基础设施投入对传统软件预算的挤压，都是市场已经讨论过的因素。但这份最新研报真正值得关注的判断，不是增速放缓，而是利润率的结构性改善和AI Agent商业化路径的加速，正在重塑中国软件行业的价值分配逻辑。市场目前过度聚焦于总量增速的短期波动，而低估了行业内部“谁在赚钱、靠什么赚钱”的深层变化。

为什么现在这个判断重要？因为4月数据同时出现了两个看似矛盾的信号：收入增速下行，但净利率从3月的9.0\%跃升至12.0\%，创下近12个月来的相对高位。这不是简单的成本压缩可以解释的。结合研报对AI产品迭代、SMB PMI波动以及重点公司新品发布的跟踪，我们认为，中国软件行业正在经历从“规模驱动”向“效率驱动+AI变现”的换挡期。那些能够将AI能力转化为可定价产品、并优化客户支出结构的公司，正在获得超出行业平均的利润率溢价。

这份报告提供的核心新信号有三个：第一，IT服务仍是收入主力，但其内部结构正在向云计算和大数据倾斜；第二，AI Agent的落地场景从概念验证进入加速变现阶段，尤其是在知识留存和生产流程优化领域；第三，海外收入占比小幅提升至3.1\%，虽然绝对值不大，但边际变化值得关注。

1. 4月收入增速放缓并非系统性风险，而是结构性换挡的必经阵痛

4月8.9\%的同比增速，叠加4M26累计10.9\%的增长，确实低于1Q26的11.6\%。但如果我们只看总量，会错过关键信息。从细分领域看，半导体设计软件以18.5\%的增速领跑，云计算和大数据紧随其后增长12.3\%，而传统软件产品和网络安全分别只有8.0\%和7.5\%的增长。这意味着，增速放缓的主要拖累来自传统品类，而非AI相关的高增长领域。

同时，4月收入环比下降12\%，这符合季节性规律——3月通常有年初集中交付效应，4月回落是常态。真正需要警惕的不是环比波动，而是SMB PMI在5月进一步降至48.5，连续两个月低于荣枯线。这预示着中小企业客户的IT预算可能在未来1-2个季度继续承压。但值得注意的是，研报明确指出“客户软件支出仍偏弱，部分项目转向AI基础模型和定制化AI模型”。也就是说，预算并没有消失，而是在重新分配——从通用软件采购转向AI基础设施和定制化开发。这对传统软件厂商是压力，但对具备AI能力的平台型公司是机会。

所以，这个数据组合的含义是：总量增速放缓正在加速行业分化，那些无法提供AI价值增量的供应商将面临更大的收入压力，而能够承接AI预算的公司反而可能获得更高的客户粘性和定价权。

4月净利率12.0\%，不仅远高于3月的9.0\%，也高于4M26的累计均值11.4\%。这个改善并非一次性事件。回顾过去12个月的数据，净利率在2025年6月曾低至8\%左右，此后呈现震荡上行趋势。2026年1-2月累计净利率达到12.5\%，3月回落至9.0\%后，4月又迅速反弹。

这种波动背后的结构性因素是什么？我们认为至少有三个值得关注的解释，但报告并未完全展开：

第一，收入结构变化。IT服务占比稳定在67\%，但其内部子项——云计算和大数据的增速（12.3\%）显著高于整体IT服务（11.9\%）。云服务的边际利润率通常高于传统系统集成，这种结构微调可能对整体利润率产生积极影响。

第二，AI产品的商业化落地开始体现为利润贡献。报告提到，AI Agent正在加速变现，应用场景包括“将资深员工的隐性知识留在企业内部”以及“以更高效率分散生产基地”。这些场景的特点是高附加值、定制化，因此定价能力和利润率都优于标准化软件。

第三，成本端的优化。4月利润率提升是否来自一次性因素（如项目验

[... middle omitted ...]

 Tour”报告中强调，AI Agent的变现正在加速，且应用场景已经从简单的聊天助手扩展到企业核心知识管理和生产流程优化。具体来说，两个场景值得深入讨论：

第一个场景是“知识留存”。很多制造型企业面临资深员工退休或离职导致的技术诀窍流失问题。AI Agent可以通过记录、结构化、检索这些隐性知识，实现“人走知识留”。这不是简单的文档管理，而是将经验转化为可调用的数字资产。这个场景的付费意愿很强，因为解决的是企业最核心的“人才风险”。

第二个场景是“生产站点多样化”。企业将AI Agent部署到多个生产基地，实现标准化操作和远程管理，从而降低对单一工厂或单一管理者的依赖。这本质上是用AI替代部分现场管理职能，直接提升运营效率。

报告提到，美图将在6月17日举办年度产品发布会，聚焦AI Agent和生成式AI产品。广联达也发布了PMLead、PMSmart等AI产品，强调平台化战略以更好地整合数据。这些动作表明，头部公司正在将AI从“功能插件”升级为“平台核心”。

\subsection{[R012] 市场真正低估的不是需求，而是供给侧的分化}
{\small\color{refgray}归类：中国地产与消费：销售分化加剧，储蓄率是消费真正拐点}

5月销售数据出来了。100强开发商合约销售同比下降4\%，与4月的5\%降幅基本持平。表面看，这只是一个“没有变得更差”的信号。但如果你只看到这一点，就错过了这份报告最核心的判断：行业总量的底部已经不再是关键变量，真正决定投资回报的，是开发商之间不可逆的销售分化。这份分化，正在重塑中国房地产行业的竞争格局。

某外资投行最新发布的研报给出了一个清晰但容易被忽视的结论：在整体销售同比微弱下滑的表象之下，国企开发商的销售已经连续第二个月实现正增长，而民企和困境开发商的销售仍在深度下跌。这不是一个临时性的波动，而是结构性力量在发挥作用。对于产业决策者和投资者而言，理解这个分化的原因、可持续性以及尚未被定价的隐含风险，比猜测政策何时放松重要得多。

1. 总量企稳的背后，是国企与民企之间超过50个百分点的增速裂口

5月的数据最值得关注的不是-4\%的同比，而是国企与民企之间的增速差距。根据报告数据，国企开发商5月销售同比增长6\%，而民企销售同比下降50\%，困境开发商下降30\%。国企与民企之间的增速差距高达56个百分点。即便与 年的四年均值相比，国企的销售已经从4月的-10\%收窄至5月的-3\%，几乎回到了历史正常水平的边缘；而民企和困境开发商仍停留在-90\%的深渊。

这种裂口不是一个月形成的。从报告提供的图表趋势看，国企销售自2024年下半年以来就持续优于民企，且差距在2025年之后进一步拉大。这意味着市场正在经历的不是一个周期性的调整，而是一个供给侧的清洗。国企凭借融资优势、土地储备质量以及品牌信任度，正在从一个萎缩的市场中攫取份额。民企的销售萎缩不仅是需求不足的结果，更是供给侧信用收缩的自我强化。

5月合约销售环比增长11\%，看起来是一个积极信号。但报告明确指出，这种改善很大程度上源于季节性因素——5月历来是比4月更强的销售月份。从 年的月度分布均值看，5月占全年的8.7\%，而4月仅为8.1\%。换言之，如果5月没有出现环比正增长，那才是一个值得警惕的信号。

更关键的参照系是四年均值。5月销售较 年平均水平仍低71\%，与4月的降幅完全一致。这说明销售的绝对水位并没有改善，只是季节性波动掩盖了趋势。报告预计，6月作为开发商冲刺半年目标的传统旺季，环比将继续正增长，但同比大概率仍将维持正负5\%的窄幅波动。对于依赖总量反弹来驱动估值的投资者，这个判断意味着“等待V型反转”的策略可能持续失效。

3. 都市更新计划的政策信号有限，真正驱动销售的是企业层面的执行力

5月28日，国务院发布了第十五个五年都市更新计划，内容包括将城市危旧房改造数量从十四五的25万套翻倍至50万套，以及计划改造4000个城中村。报告发布当天，地产板块上涨4\%，跑赢恒指约3个百分点。但报告对这个政策的评价是“并不令人兴奋”。

核心原因有三：第一，该计划未提及现金补偿，这意味着拆迁居民不会获得用于购买商品房的现金，对住宅销售的直接影响有限；第二，政策方向与以往叙述高度相似，缺乏增量信息；第三，政策目标更偏向改善居住条件，而非刺激楼市。报告隐含的判断是：市场对政策的短期反应可能过度，而真正能带来销售增长的，是开发商自身的项目执行力、产品定位和推盘节奏。

这一点在部分国企开发商身上得到了验证。5月单月，华润置地销售同比增长28\%，招商蛇口增长20\%，越秀地产增长18\%，中海地产增长14\%。这些数字并非来自政策红利，而是来自精准的推盘和有效的去化。相反，中国金茂5月销售同比下降16\%，但报告指出这主要是因为推盘数量从去年的6个骤降至1个，而非需求问题——前5个月金茂累计销售仍增长11\%，是少数保持正增长的开发商之一。

4. 报告尚未完全回答的关键问题：国企的份额增长能否持续转化为利润

这是整份报告最值得深入追问的部分。国企销售份额的持续扩大已经是一个

[... middle omitted ...]

。报告没有对这一问题展开定量分析，但这恰恰是投资者需要自行判断的核心变量。如果国企的销售增长只是以牺牲利润率为代价，那么当前的估值溢价可能已经透支了未来的回报空间。

对于产业决策者和投资者，当前阶段最有效的观察框架不是猜测政策何时转向，而是追踪三个信号：第一，国企与民企的销售增速差是否继续扩大，如果差距从当前的50个百分点以上开始收窄，可能意味着民企的流动性压力边际缓解，但这也可能只是基数效应；第二，国企内部的分化是否在加剧——中海、华润、招商蛇口等头部国企能否持续跑赢同行，这决定了行业集中度的最终形态；第三，6月半年冲刺后的7-8月淡季表现，如果淡季销售同比降幅没有显著恶化，那么总量企稳的假设将获得更多支撑。

这些信号不会在同一天出现，但会在未来2-3个月内逐步清晰。报告给出的核心框架是：在总量没有系统性改善之前，选股逻辑应从“行业beta”转向“公司alpha”。那些能够持续实现正增长、且增长并非来自一次性推盘而是来自可持续竞争力的开发商，才是值得关注的标的。

\subsection{[R013] 市场真正低估的不是需求，而是供给侧的再定价}
{\small\color{refgray}归类：中国地产与消费：销售分化加剧，储蓄率是消费真正拐点}

过去几周，中国房地产市场的讨论焦点大多集中在需求端：政策还能不能更松？居民购房意愿有没有真正回暖？二手房成交量是否可持续？

这些问题的答案，在最新一期某外资投行的中国地产周报中有了初步的、但被市场普遍忽略的指向：成交量在恢复，但更关键的结构性变化发生在供给侧——库存正在被系统性消化，而市场定价尚未完全反映这一过程。

这份周报覆盖了截至第22周的数据。新房成交面积周环比上升13\%，同比转正至+4\%，创下清明假期以来的周度新高。二手房成交虽周环比小幅回落2\%，但同比仍高达+26\%。更重要的是，库存月数从4月均值29.3个月下降至27.6个月。这不是一个简单的季节性波动。它是一个信号：供给侧的出清正在加速，而市场的定价逻辑仍停留在“需求疲软”的旧框架里。

本文基于这份报告的核心数据与判断，提炼出五个值得产业决策者与投资者关注的层级。它们共同指向一个结论：当前资产定价的折价，更多反映的是对流动性的恐慌，而非对资产本身价值的理性评估。

第22周新房成交面积环比增长13\%，同比转正至+4\%，将年初至今的累计同比跌幅收窄至-14\%。与2024年和2023年同期相比，这一成交量分别低15\%和47\%。单看同比数据，似乎并不惊艳。但关键在于趋势的连续性。

报告显示，5月新房成交面积中位数环比上升5\%，同比持平。这意味着4月下旬以来的政策放松——包括广州在4月底的全面松绑、5月中旬南沙区的试点——正在转化为实际成交，而非停留在看房阶段。看房活动量周环比持平，说明需求并未出现“脉冲式释放后迅速冷却”的典型特征，而是在稳步释放。

更值得关注的是库存的消化速度。第22周库存总量周环比下降0.5\%，较2025年底已下降4.4\%。库存月数降至27.6个月，低于4月均值29.3个月。这个数字仍然偏高，但下降的趋势正在形成。如果这一趋势持续，开发商面临的去化压力将从“生死存亡”级别降至“可以承受”级别。

对行业的含义很清晰：新房市场最危险的时刻——库存高企、成交冻结、现金流断裂——正在成为过去。但市场定价尚未完全计入这一变化。

2. 二手房市场的“量升价稳”正在重塑交易结构，但卖家预期开始分化

二手房市场是这份报告中最值得细读的部分。第22周成交量周环比下降2\%，但同比仍高达+26\%。年初至今累计成交量同比持平，但较2024年和2023年同期分别高出23\%和7\%。这说明二手房市场已经完成了从“深度调整”到“新常态”的切换。

报告中的一个关键信号来自两个指数的分化：CSI（经纪人价格预期指数）周环比上升0.2个百分点，同比上升3.7个百分点，仍处于50以上的乐观区间；而CAI（卖家报价指数）周环比下降0.5个百分点，同比下降3.6个百分点，已跌破20。

两者的背离意味着：经纪人仍然看好价格，但卖家开始松口。这通常是市场从“卖方定价”转向“买方定价”的过渡阶段。如果卖家降价意愿增强，成交量有望进一步放大。但这也意味着，价格底部的确认需要更多时间——不是所有人都能等到反弹。

对于投资者而言，二手房市场的“量先行、价滞后”格局，为判断行业底部提供了更精细的观察窗口。 当CAI停止下行、CSI维持在50以上，价格的拐点才会真正到来。

3. 广州的“全市回购”是一次政策实验，其效果将决定其他城市的跟进速度

报告详细分析了广州的最新政策：在全市范围内推出二手房回购计划，目标为市中心总价低于300万元、面积小于70平方米的房源，收购后转为社会保障性住房和人才住房。同时，广州优化了商业贷款转公积金贷款的通道，将最高贷款成数从70\%提升至80\%，并扩大了借款人资格范围，包括多套房持有者、此前使用过公积金贷款的借款人，以及大湾区或广州都市圈内的公积金缴存者。

这不是一个孤立的城市行动。报告将其与上海模式对标，暗示这可能是全国性“以旧换新”

[... middle omitted ...]

思的估值对比：覆盖的离岸开发商目前交易价格较2026年底预期净资产价值（NAV）折价24\%，对应0.6倍2026年预期市净率（P/B）。而在2008年下半年、2011年下半年和2014年上半年的历史底部，折价幅度分别为39\%、73\%和58\%，对应P/B分别为0.7倍、0.9倍和0.9倍。

在岸开发商的折价同样显著：较2026年底预期NAV折价20\%，对应0.5倍2026年预期P/B，而历史底部的折价分别为67\%、64\%和61\%，对应P/B分别为1.6倍、1.5倍和1.2倍。

*当前的P/B估值已经低于2008年金融危机时的水平，但折价幅度却小于历史底部。** 这意味着市场定价隐含的假设是：资产价值还将继续缩水。但如果供给侧出清加速、库存消化持续，这种假设可能过于悲观。

报告特别指出，覆盖的强信用国企开发商在第22周股价平均上涨2\%，其中中国金茂和招商蛇口分别上涨5\%。而民营开发商股价持平，其他国企开发商上涨2\%。分化依然明显，但强信用国企正在获得市场的重新定价。

\subsection{[R014] 中国消费的真正拐点，不在于收入增速，而在于储蓄率何时开始下降}
{\small\color{refgray}归类：中国地产与消费：销售分化加剧，储蓄率是消费真正拐点}

中国消费的真正拐点，不在于收入增速，而在于储蓄率何时开始下降

过去两年，市场对中国消费的讨论几乎全部围绕一个核心变量：居民收入。收入增速放缓，消费自然承压，这个逻辑链条直白且直观。但这份某外资投行在2026年6月发布的研报，提供了一套更值得深挖的观察框架。它揭示了一个被普遍低估的事实：中国消费的真正拐点，可能并不由收入增速的边际改善驱动，而是取决于居民储蓄率何时从当前33\%-34\%的高位开始趋势性下降。

报告的核心数据表给出了一个令人警醒的信号。从2025年1月到2026年4月，中国社会消费品零售总额同比增速从6.4\%一路滑落至0.2\%。同期，城镇居民可支配收入增速从4.4\%降至4.2\%，虽有波动但并未出现断崖式下跌。也就是说，收入端并没有恶化到足以解释零售增速近乎归零的程度。真正在起作用的，是居民部门持续高企的预防性储蓄意愿——收入没有大幅减少，但人们更不敢花钱了。

这组数据放在一起，指向一个更本质的判断：中国消费市场当前面临的不是周期性疲软，而是结构性储蓄陷阱。居民资产负债表修复、就业预期不稳、财富效应缺失三重压力叠加，使得消费倾向被系统性压低。而这份报告的价值，正在于它没有止步于“消费疲软”这个共识，而是拆解了疲软背后的分层逻辑，并给出了一个关键的时间锚点——2026年下半年。

如果把中国消费下滑归因于“大家都没钱了”，这个判断过于粗糙。报告中的宏观数据追踪表显示，城镇居民可支配收入增速在2025年至2026年间维持在4\%左右，虽然低于疫情前水平，但并非负增长。真正拖累消费的是财富效应——房产价值缩水和资本市场波动对居民资产负债表造成的冲击。

报告特别列出了“二手房价同比增速”这一指标。从2025年1月至2026年4月，该数据持续徘徊在-12\%至-14\%的区间，没有任何一个月份出现收窄迹象。这意味着中国家庭最重要的资产——住房，在过去16个月里持续以两位数的速度贬值。对于持有房产的中产家庭而言，账面财富的缩水直接压制了消费意愿，尤其是对可选消费和耐用品的支出。

与此同时，报告中的“新经济平均月薪同比增速”数据揭示了另一个结构性问题。2025年1月至4月，该指标连续为负，最低达到-11\%。虽然2025年下半年至2026年初有所转正，但2026年3月和4月再次回落至3\%和几乎为零的水平。新经济领域曾是高收入岗位和消费升级的重要引擎，其薪资增速的波动直接影响了年轻高收入群体的消费信心。

这两组数据合在一起意味着：中国消费疲软的本质，不是所有人都没钱，而是“有钱人”的资产在缩水，“年轻人”的收入预期在恶化。消费市场正在经历一个财富效应缺失与收入预期下移叠加的“双杀”阶段。

2. 零售数据分化的背后，是必需品和可选品正在走向两个完全不同的市场

报告中的零售销售细分表揭示了一个被总量数据掩盖的重要现象：消费内部正在发生剧烈的结构性分化，而不是简单的全面萎缩。

2026年4月，整体零售增速仅为0.2\%，但扣除汽车后的增速为1.8\%。更值得关注的是分品类的表现。食品、饮料、烟酒类增速高达5.4\%，其中粮油食品类增长4.1\%，饮料类增长3.6\%，烟酒类更是达到11.7\%。而可选消费则全面承压：金银珠宝类暴跌21.3\%，家用电器和音像器材类下降15.1\%，家具类下降10.4\%，体育娱乐用品类下降8.0\%。

这种分化不是短期的消费降级，而是居民消费函数的结构性重置。当财富效应缺失且收入预期不稳定时，家庭会优先压缩“可推迟的支出”——珠宝、家电、家具、运动装备都属于这一类。而食品饮料等必需消费品，由于需求刚性，不仅没有萎缩，反而因为消费场景的回归和渠道下沉而保持正增长。

报告还提供了一个更有洞察力的对比：计算与2019年相比的年复合增长率（CAGR）。2026年4月，整体零售的CAGR为2.9\%，但线上零

[... middle omitted ...]

支持可能只在就业市场进一步恶化时才会推出。

这组判断实际上在说：市场对政策刺激的期待可能需要下调。过去两年，每当消费数据走弱，市场就会预期更强力的财政或货币刺激。但报告的数据表明，政策制定者更倾向于通过投资和基础设施来稳定经济，而非直接刺激消费。2026年耐用消费品以旧换新的预算并未扩大，这意味着家电、汽车等品类可能面临政策红利退坡的风险。

更关键的是，报告中的“居民储蓄率”数据始终维持在33\%-34\%的高位。这个数字意味着，即使政府发放消费券或提供补贴，居民部门也可能选择将这部分额外收入转化为储蓄，而非消费。预防性储蓄动机的根源在于就业预期和社保体系的不确定性，这不是短期补贴能够解决的。

报告还特别指出了“青年失业率”这一指标。2025年至2026年间，16-24岁青年失业率（不含在校学生）持续在15\%-18\%之间波动，从未低于15\%。青年群体是边际消费倾向最高的群体，也是未来消费升级的主力。当他们面临持续的就业压力时，整个消费市场的长期增长动力都会受到侵蚀。

\subsection{[R015] 全球工业技术周期的“双顶”正在确认，但市场对“中国需求之外的结构性加速”定价不足}
{\small\color{refgray}归类：区域市场：亚洲央行鹰派转向，澳洲投资逻辑切换 / 工业科技：全球自动化周期出现“双顶”，中国以外复苏加速}

全球工业技术周期的“双顶”正在确认，但市场对“中国需求之外的结构性加速”定价不足

这可能是当前最值得关注的宏观产业信号：全球工业自动化周期不仅没有见顶回落，反而在2026年二季度出现了罕见的“双顶”形态。某外资投行最新发布的《Asian Industrial Technology Barometer》显示，中国工厂自动化需求在4-5月维持近峰值增速，而中国以外地区的复苏正在显著加速。

这份报告的核心价值不在于罗列数据，而在于揭示了一个被市场普遍低估的结构性变化——本轮周期并非简单的库存回补，而是由“物理AI驱动的自动化渗透率加速”与“全球制造业资本开支重启”共同推动。市场目前对后者的定价远不够充分。

以下是我们从这份长达数十页的月度跟踪报告中提炼出的五个核心洞察，以及一个尚未被完全回答的关键问题。

报告中最容易被忽视的信号，来自“中国以外”的数据。日本机床订单在4月同比增长45.1\%，其中来自欧洲的订单增速从3月的21.9\%跃升至35.6\%，来自日本本土的订单更是从2.5\%飙升至43.4\%。这不是单一市场的脉冲，而是全球制造业同步扩张的证据。

同期，美国、欧元区和日本的制造业PMI全部维持在扩张区间（分别为55.3、51.4和54.5）。更重要的是，日本工业机器人出口量在4月同比增长40.7\%，其中对欧洲出口增速从35.3\%加速至68.3\%，对亚洲（除中国）出口增速从24.0\%飙升至88.6\%。

这些数据叠加在一起，意味着什么？全球工业资本开支的“第二波”正在形成。第一波是 年疫情后的补库周期，第二波则是基于产能利用率回升和自动化长期趋势的主动投资。中国作为全球最大的自动化市场，其需求维持高位是“定海神针”，但真正超预期的边际增量，来自欧洲和亚洲其他地区的加速复苏。对于布局全球市场的自动化零部件和机器人企业而言，这可能是未来12-18个月最重要的业绩驱动力。

报告指出，4月中国工业利润同比增长24.7\%，较3月的15.8\%显著加速。驱动因素有两个：营收增长加速和运营利润率改善。更值得关注的是，高端制造业在4个月的累计利润增速高达44.8\%，远超整体制造业的20.4\%。

利润增速远超营收增速，说明企业正在从“以量换价”转向“量价齐升”。这与PPI数据相互印证——4月PPI指数回升至101.3，结束了此前连续多月的通缩状态。上游价格回暖叠加下游需求稳定，意味着制造业企业的议价能力正在增强。

这一变化对投资者和产业决策者的含义是：过去几年市场对中国制造业的定价逻辑是“产能过剩、利润承压”，但当前的数据正在挑战这一假设。如果利润改善的趋势能够持续，那么中国自动化龙头企业将获得“估值修复”和“盈利上调”的双重驱动。报告没有明确讨论这一点，但从数据推演，这可能是最值得深入研究的投资主题。

报告提到日本机床订单正在形成类似以往周期的“双顶”形态。4月订单同比增长45.1\%，环比仅微降2.3\%，这已经是连续第二个月维持高位。所谓的“双顶”，是指订单在经历一轮强劲增长后，并未快速回落，而是在高位震荡后再次上冲。

从历史经验看，这种形态通常意味着周期被拉长。原因在于，第一轮需求来自库存回补和应急采购，第二轮则来自真实的产能扩张和自动化升级。当前全球制造业产能利用率仍在回升通道中（美国4月为75.7\%，欧元区1季度为77.9\%），尚未达到触发大规模资本开支的阈值。如果后续利用率继续提升，那么自动化设备的需求可能还没有见到真正的峰值。

对于产业决策者而言，这意味着供应链备货和生产排期的节奏需要调整——不能按照传统周期“见顶即收缩”的逻辑来规划。对于投资者而言，这意味着自动化板块的盈利预期可能持续上修，市场对周期见顶的担忧可能被证伪。

4. 物理AI正在成为自动化需求的结构性加速器，但量化影响仍需验证

报告在

[... middle omitted ...]

即自动化渗透率曲线将变得更加陡峭。

然而，报告并未提供具体的量化证据来区分“周期因素”和“AI加速因素”各自贡献了多少。这是全文最值得关注的一个未解问题。如果物理AI的加速效应被后续数据证实，那么当前市场对机器人及相关零部件企业的估值框架可能需要系统性重估——从“周期股”切换为“成长股”。反之，如果AI的影响被夸大，那么当前的高增速可能只是周期性的，后续将面临均值回归。

5. 一个尚未完全回答的关键问题：中国自动化需求的“质量”如何？

报告显示，中国4月工业机器人产量同比增长15.1\%，低于3月的24.4\%，也低于年初至今的25.7\%。同时，日本对中国机器人出口增速虽在加速（32.8\%），但绝对值低于对其他地区的增速。这些数据表面上看是“仍在增长，但动能边际放缓”。

但更深层的问题是：中国当前的需求增量，是来自传统制造业的自动化升级，还是来自新能源、半导体等新兴产业的产能扩张？如果是前者，那么需求持续性可能更强；如果是后者，则可能面临政策周期和补贴退坡的风险。

\subsection{[R016] 中国车市真正的分化不在电动化，而在“谁的海外能扛”}
{\small\color{refgray}归类：中国地产与消费：销售分化加剧，储蓄率是消费真正拐点}

北京车展的喧嚣刚落，市场交出的答卷并不令人振奋。某外资投行最新发布的研报指出，5月中国乘用车市场整体批发销量可能再度录得超过20\%的同比下滑。这不是一个意外，而是缺乏增量政策刺激下，本地需求持续疲软的延续。

但在这片低迷中，一个结构性的信号值得所有产业决策者关注：新能源汽车的渗透率预计将再次站上60\%的高位。电动车企继续显著跑赢燃油车玩家。市场真正在发生的，不是“电动化”与“燃油车”之间的胜负已分，而是电动车企内部，一场围绕“海外增量”与“本地份额”的残酷分化正在加速。

这份报告的真正洞察在于：在一个整体需求缺乏弹性的市场里，只有两类玩家能够获得超额表现——那些海外业务贡献持续扩大的公司，以及那些即便在国内市场，也能凭借技术和产品升级实现份额逆势增长的公司。前者决定了增长的“天花板”，后者决定了生存的“安全垫”。

1. 整体市场低迷，但渗透率新高背后是“结构性”而非“周期性”的驱动力

5月的数据再次确认了一个判断：中国乘用车市场的本地需求，已经进入了一个缺乏政策刺激的“存量博弈”阶段。北京车展未能带来实质性的订单转化，CPCA前24天的初步数据也指向了又一个双位数下跌的月份。

然而，电动车渗透率持续维持在60\%以上，这并非简单的消费偏好转移。更深层的原因在于，电动车企，尤其是头部玩家，正在通过产品和技术迭代，主动创造需求，而非被动等待市场回暖。这种“供给创造需求”的能力，是燃油车厂商在当前环境下难以复制的。

这意味着，对于投资者和行业观察者而言，关注月度总量已经失去意义。真正的信号，藏在每个玩家的“结构性”变量里：海外占比、技术升级节奏、以及新车型的订单转化效率。

2. 比亚迪的“止跌”信号：海外扛起半壁江山，但国内复苏仍需验证

比亚迪5月交付37.7万辆，同比持平，环比增长20\%。这是自去年9月以来，比亚迪首次实现同比不下跌。但仔细拆解数据，更能看清真相：海外销量达到16.06万辆，同比增长80\%。剔除海外，其国内批发量约为21.63万辆，同比下降25\%。虽然相比4月的-39\%有所收窄，但国内市场的压力依然巨大。

研报的积极信号在于，比亚迪的“技术牌”开始重新吸引消费者。其超快充车型的市场反馈积极，而5月28日发布的自主研发4nm智能驾驶芯片，也展示了其垂直整合的深度。随着刀片电池产能的进一步爬坡，叠加海外市场的稳健增长，研报判断比亚迪有望在2026年下半年实现国内出货量的同比改善。

这里的关键洞察是：比亚迪的“止跌”并非国内市场的反转，而是海外增量对冲了国内的下行。对于一家市值庞大的公司而言，海外业务的持续高增长，才是其估值获得支撑的核心逻辑。但国内能否真正企稳，仍需后续几个月的订单数据来验证。

3. 新势力阵营的分水岭：不是“谁卖得多”，而是“谁在增长”

5月的数据清晰地划出了一条分水岭。零跑汽车以8.16万辆的批发量，同比增长81\%，环比增长14.3\%，再创历史新高，成为新势力中当之无愧的“增长之王”。蔚来汽车交付3.77万辆，同比增长62.3\%，环比增长28.4\%，其中主品牌NIO连续7个月销量破万，子品牌乐道和萤火虫也在贡献增量。

而理想汽车交付3.34万辆，同比下滑18.4\%，环比下滑2.2\%。小鹏汽车交付3.22万辆，同比下滑4.1\%，环比微增3.7\%。小米和鸿蒙智行（HIMA）则分别维持在3万辆和4.61万辆的水平。

这份数据揭示了一个残酷的事实：在渗透率已经很高的市场里，增长不再是普惠的。零跑和蔚来的高增长，背后是产品定位的精准和订单的持续转化。零跑凭借极致性价比和产品力，正在收割对价格敏感但对技术有要求的用户。蔚来则通过多品牌战略（NIO、乐道、萤火虫）实现了对20万至60万价格带的全面覆盖，ES8的持续热销和ES9的积极反馈，证明了其高端定位的护城河。

理想

[... middle omitted ...]

主要市场的关税和政策风险，是否会在下半年集中释放？这些因素将直接影响海外业务的盈利能力和可持续性。

同样，零跑和蔚来的海外布局尚在早期，其海外销量在总销量中的占比仍然较低。对于这些公司而言，国内市场的增长仍然是其当前估值的核心支撑。海外市场是“第二曲线”，但这条曲线的斜率有多陡，还需要更长时间的数据来验证。

这一未解问题，恰恰是投资者在阅读完这份报告后，最需要深入探讨的议题。它决定了我们对该行业长期增长潜力的判断。

5. 给决策者的观察框架：从“销量排名”转向“三个核心指标”

基于以上分析，对于关注中国电动车产业的决策者和投资者，传统的月度销量排名已经失去了参考价值。真正有效的观察框架，应该聚焦于三个核心指标：

*第一，订单转化效率与“在手订单”的可见性。** 研报中反复提及“根据我们的行业渠道调研，订单势头良好”。这意味着，月度交付量是滞后指标，而“订单”是领先指标。关注公司管理层是否提供订单指引，以及第三方渠道的订单追踪数据，比关注已交付数据更具前瞻性。

\subsection{[R017] 人民币中间价模型正在发出一个被忽视的定价信号}
{\small\color{refgray}归类：宏观与利率：通胀结构重塑，美元定价逻辑切换}

市场参与者普遍将人民币汇率走势归因于中美利差、贸易摩擦和央行干预意愿。但某外资投行最新发布的中间价模型预测显示，一个更底层、更技术性的变量正在发生结构性变化，其信号强度可能被当前主流叙事所低估。

这份截至2026年6月2日的模型预测显示，在没有逆周期因子调整的情况下，美元兑人民币中间价模型预测值为6.7762，较前一日6.8167的预测值下降了405个基点。即便计入逆周期因子，模型预测值仍为6.7988，较前一日中间价低179个基点。

这不是一个简单的短期波动。这是模型对一篮子货币权重、隔夜市场波动以及政策干预框架三者之间动态关系的重新定价。理解这个信号，比猜测央行下一步动作更有意义。

405个基点的单日模型预测下降，在近一年的数据中并不常见。回顾模型误差的历史走势，2025年初模型曾出现高达1800个基点的负向误差，而进入2026年后，误差范围显著收窄至600个基点以内。这意味着模型的预测能力正在增强，或者说，市场定价与模型隐含的均衡水平之间的偏离正在收敛。

但更重要的是，这次下降并非孤立事件。模型预测的下降建立在隔夜一篮子货币变动的加权贡献之上。根据报告披露的权重数据，韩元、俄罗斯卢布、欧元和澳元是贡献最大的四个币种，分别贡献了27.5、23.2、20.1和8.3个基点。这些币种的集体走强，直接推动了人民币中间价模型预测的下行。

这意味着什么？人民币的定价逻辑正在从单纯的中美双边博弈，回归到更广泛的贸易加权篮子框架。任何忽略这一框架的汇率判断，都可能高估或低估了中间价的真实弹性。

模型预测值与计入逆周期因子后的预测值之间的差值，是理解政策意图的关键窗口。当前，两者相差226个基点。这个数字本身并不惊人，但结合近期中间价的实际走势来看，逆周期因子的使用频率和幅度正在发生微妙变化。

报告中的日度中间价变动图显示，2025年11月至2026年5月期间，中间价绝大多数时间保持稳定，仅在2026年3月19日出现过约150个基点的调整。这种“长期稳定、偶发调整”的模式，本身就是逆周期因子在发挥作用——它平滑了市场波动，维持了中间价的预期稳定性。

但模型预测的下降幅度明显大于近期中间价的实际调整幅度。这意味着，如果模型预测准确，央行要么需要通过逆周期因子来吸收这部分压力，要么允许中间价向模型预测值靠拢。无论哪种选择，逆周期因子的角色都在从被动的“对冲波动”转向主动的“信号管理”。投资者需要关注的不是逆周期因子是否存在，而是它在什么阈值下被触发，以及触发后的调整幅度是否超出市场预期。

报告提供的模型误差历史数据值得深挖。2025年初，模型误差曾达到-1800个基点，这意味着模型大幅高估了中间价。此后误差逐步收窄，但到2026年4月，误差仍维持在-600个基点左右。

误差的收窄本身是好事，但它也意味着模型正在更紧密地跟踪实际中间价。一旦模型预测与实际中间价出现系统性偏离，修复的速度和幅度可能会比过去更快。当前模型预测较前一日下降405个基点，而实际中间价尚未充分反映这一变化。如果误差持续为负，市场可能会面临中间价加速向模型预测靠拢的修正压力。

这里有一个关键问题：模型误差的收窄是否意味着中间价定价效率的提升？如果是，那么未来中间价的波动性可能会增加，因为模型捕捉到的信号将更快地传导至实际定价。如果不是，那么误差的收窄可能只是暂时的，模型与实际之间的差距仍可能重新扩大。报告没有给出明确的结论，但这个问题本身值得持续跟踪。

报告附带的2026年下半年事件日历，为理解汇率走势提供了一个超越技术分析的框架。7月底的政治局会议、11月的APEC会议、以及年底中国领导人访美，构成了三个重要的政策窗口期。

这些事件对汇率的影响路径并非线性。政治局会议可能释放稳增长信号，从而支撑人民币；APEC会议可能涉及贸易议程

[... middle omitted ...]

未完全展开。

从逻辑上推断，如果韩元和卢布的权重贡献持续上升，这可能意味着模型正在重新校准对新兴市场货币的敏感度。而如果欧元的权重保持稳定，则说明美元指数篮子依然是核心锚。这些权重的变化，本质上是对全球贸易格局和资本流动结构的隐性映射。

另一个未解问题是：逆周期因子的参数是否在变化？报告给出了计入逆周期因子后的预测值，但没有披露逆周期因子的具体计算方式或参数调整逻辑。这意味着，市场无法独立验证逆周期因子的“真实”影响。对于机构投资者而言，这既是风险，也是信息优势的来源——谁能更准确地拆解逆周期因子的行为模式，谁就能在汇率定价中占据先机。

这些未解问题，恰恰是这份报告最有价值的部分。模型预测给出了一个明确的信号，但真正的判断力来自对信号背后机制的理解。如果你对如何拆解逆周期因子的行为模式、如何跟踪模型权重的动态变化、以及如何将这些信号融入实际的汇率风险管理框架感兴趣，欢迎来星球微信群里继续讨论。我们会围绕这份报告的核心假设和未解问题，做更深入的框架推演和情景分析。

\subsection{[R018] 韩国出口的真相：AI 推高了价格，但没有推宽复苏}
{\small\color{refgray}归类：未被主题摘要引用}

韩国5月出口同比增长53.2\%，贸易顺差在2026年前五个月就超过了2017年全年纪录。这组数字表面上看是经济强劲的信号，但某外资投行最新研报揭示了一个更微妙的判断：这场出口繁荣的核心驱动力是AI芯片价格暴涨，而非出口量的普遍扩张。更关键的是，韩国央行使用的价格指数可能高估了实际出口对GDP的拉动作用，这意味着市场对韩国经济复苏的乐观程度可能超出了基本面支撑。

这份报告的核心洞察不在于“韩国出口很好”，而在于“出口好得不够均匀”——AI相关的高价芯片撑起了大部分增长，而汽车、机械、钢铁等传统部门仍在萎缩。这种K型复苏的结构性特征，对理解韩国央行未来的加息路径、以及对全球科技供应链的定价逻辑，都有重要的启示。

1. 芯片价格暴涨正在制造“名义繁荣”，但实际出口量已经连续两个月下滑

5月韩国半导体出口同比增长169.4\%，创下月度新高。其中DRAM出口增长369.8\%，NAND出口增长206.8\%。DDR5 16Gb固定价格从2025年5月的4.80美元飙升至2026年5月的37.50美元，涨幅高达682.1\%。这些数字看起来像是一场全面的出口爆发，但报告揭示了一个关键的分化：以吨位计算的海关出口量在5月同比下跌18\%，连续第二个月出现两位数下滑。

这意味着韩国出口的增长几乎完全由价格驱动，而非数量驱动。对于依赖“量价齐升”逻辑来支撑经济复苏预期的投资者来说，这是一个需要警惕的信号。当芯片价格周期进入平稳或下行阶段，出口增速将面临急剧回落。

报告进一步指出，韩国央行使用的出口价格指数可能未能及时反映近期HBM和AI相关芯片价格的暴涨。由于央行指数基于2020年基准结构，它对新产品、高单价产品的价格变化反应较慢，这可能导致其对实际出口量的估算偏高，进而高估出口对GDP的贡献。这个技术细节对理解韩国央行的政策立场至关重要：如果央行基于被高估的出口数据维持增长乐观，它可能会更早、更鹰派地加息。

2. AI相关出口已经占到总出口的42\%，但其余部门仍在挣扎

半导体出口在5月占总出口的比重超过42\%，加上计算机和SSD（同比分别增长290.7\%和337.4\%），AI基础设施相关的出口已经成为韩国出口的绝对主力。值得注意的是，SSD出口中很大一部分是面向AI数据中心的企业级产品，这反映出AI投资正在从芯片本身向存储和计算基础设施扩散。

然而，报告清晰地描绘了一幅“冰火两重天”的画面。汽车出口同比下降5.9\%，其中内燃机汽车出口下降14.4\%，虽然混合动力和电动车有所增长，但整体仍被供应链中断、中东地缘政治、美国关税以及海外生产转移所拖累。石油产品名义出口增长46.6\%，但出口量下降23.8\%；石化产品名义增长11.1\%，出口量下降25.5\%。钢铁出口同比下降20.3\%，通用机械下降6.0\%。

这些数据合在一起意味着：韩国经济的“引擎”是高度集中的。AI芯片及其相关硬件是唯一强劲的发动机，而传统的制造业支柱——汽车、钢铁、机械——要么停滞，要么萎缩。这种结构意味着出口增长对国内就业、工资和消费的传导效应将十分有限。报告对此的表述是“K-shaped exports”，即出口本身也是K型分化。

报告中最具技术含量、也最可能被市场忽视的洞察，是海关出口量与韩国央行出口量指数之间的显著分歧。自2025年以来，这两个指标开始走向不同方向：海关量指数在2026年5月同比下跌18\%，而央行指数仍维持两位数的正增长。

分歧的根源在于价格折算方法的不同。海关基于产品层面的单位价值，而央行使用基于2020年基准结构的出口价格指数进行折算。当芯片价格在短期内出现数倍暴涨时，央行的价格指数因为基准结构滞后，无法充分反映价格变化，导致实际出口量被高估。

这个误差的潜在影响是巨大的。如果出口对GDP的贡献被高估，那么韩国央行基于增

[... middle omitted ...]

增长。半导体是对华出口的主要贡献者，同比增长超过200\%。这组数据表面上看是“中国需求复苏”的信号，但报告暗示了更深层的结构性变化：中国正在成为韩国半导体供应链中的“需求方”，而非传统意义上的“组装出口方”。

对美出口中半导体增长651\%、计算机增长675\%，同样反映出AI基础设施投资对美国市场的拉动。而对东盟出口增长58.4\%，创下月度新高，半导体、石油产品、显示器和石化产品全面增长，显示韩国在区域IT和中间品供应链中的枢纽地位正在加强。

这些区域分化背后有一个共同的逻辑：AI投资正在重塑全球半导体贸易流向。韩国作为全球最大的存储芯片出口国，成为这一轮AI资本支出的最大受益者之一。但报告也提示，部分对东盟出口的增长是价格驱动而非数量驱动，石油和石化产品尤其如此。这意味着如果AI投资节奏放缓或芯片价格回落，这些区域的出口增速可能比预期更快地降温。

这份报告清晰地揭示了韩国出口“量缩价涨”的结构性特征，但有一个关键问题没有完全展开：芯片价格的上涨周期还能持续多久？

\subsection{[R019] 市场真正低估的不是AI，而是“非AI”资产的盈利韧性}
{\small\color{refgray}归类：风险偏好与资金流向：机构与散户出现显著分歧}

当某外资投行在2026年6月发布的“Conviction List Directors‘ Cut”三周年特刊中，将Casella Waste Systems、TPG、Tyson Foods和Block同时纳入其最高确信度买入名单时，传递出的信号远比名单本身更值得关注。这份更新并非简单的股票轮换，而是对当前市场“单一交易”格局的一次系统性对冲。

自2022年11月ChatGPT问世以来，标普500指数累计上涨84\%。但这份报告揭示了一个被多数投资者忽略的事实：在这84\%的涨幅中，近三分之二来自实实在在的盈利增长，而非估值泡沫。报告团队测算，2026年标普500每股盈利预测为340美元，较ChatGPT发布前的远期预测高出52\%。这意味着，即使估值不再扩张，仅靠盈利兑现，标普500就能达到其年底8000点的目标。

然而，真正值得深思的不是AI主题能否持续，而是当市场将所有注意力集中在“一个交易”上时，那些与AI无关、却拥有独立盈利动能的资产，正在被系统性地低估。这正是本期Conviction List的核心主张：在“怀疑滋养信仰”的市场环境中，盈利增长的质量和独立性，才是下一个超额收益的真正来源。

这份名单的生成机制本身就值得关注。它不是分析师各自推荐股票的简单汇总，而是由一个高级研究经理组成的子委员会，从各行业分析师处筛选出最高确信度、最具差异化观点的买入评级。委员会的核心筛选标准是：高确信度、独特的差异化观点、可见的催化剂路径以及高风险调整后回报潜力。

从最新调整来看，新加入的四只股票——Casella Waste Systems、TPG、Tyson Foods和Block——分布在工业、金融、消费和科技四个截然不同的领域。它们唯一的共同点是：分析师对其盈利预测均高于市场共识。Casella Waste Systems的分析师预计其有机EBITDA将实现中等个位数百分比的持续增长，驱动因素来自固废定价能力以及SG\&A节约和并购协同的额外上行空间。TPG的分析师则预计其基础管理费增长在2026和2027年均将实现约20\%的同比增速，2027年盈利预测高出市场共识8\%。Tyson Foods的分析师预计其盈利增长将高于市场共识，受益于多元化蛋白组合和运营改善。Block的分析师则预计2026和2027年每股盈利分别增长64\%和36\%，2027年预测高出市场共识4\%。

这些公司的共同特征，并非属于同一个热门主题，而是都拥有分析师认为市场尚未充分定价的独立盈利动能。在一个被AI主题主导的市场中，这种“盈利独立性”本身就是一种稀缺的alpha来源。

报告用相当坦诚的篇幅回顾了过去一年的表现。自上次年度报告以来，标普500上涨29\%，而Conviction List的等权重指数化回报跑输标普500 180个基点。但深入来看，这一表现并非没有亮点。

报告披露了一个关键数据：过去12个月，约80\%在列入名单后发布了首季财报的公司实现了盈利超预期。这意味着选股委员会在识别盈利拐点方面的能力是可靠的。问题在于，在一个“主题往往压倒数字”的市场中，处于资金流出板块的股票即使业绩超预期，股价也未必能获得应有的回报。

这正是当前市场结构的核心矛盾。AI主题吸引了绝大部分资金流入，导致其他板块即使基本面改善，也难以获得估值重估。但报告同时指出，许多个股的独立成功案例恰恰来自那些分析师关注的关键基本面驱动因素最终实现的公司——沃尔玛、Hershey、Capital One、Woodward、Philip Morris等。在这些案例中，向上的盈利修正最终确实推动了股价上涨。

这意味着，市场并非完全忽视基本面，而是需要更长的时间窗口来验证。对于能够耐心持有的投资者而言，当前的非AI资产可能正处于一个“盈利改善但估值压缩”的窗口

[... middle omitted ...]

PU而非GPU。这意味着，即使AI基础设施投资继续增长，其受益结构也可能发生根本性转变。

这些未解问题指向一个核心判断：当前市场对AI主题的定价，隐含了太多尚未验证的假设。价值在供应链中的分配路径、代理式AI的商业化时间表、以及非AI企业的盈利增长能否持续，都是需要持续跟踪的关键变量。这也解释了为什么选股委员会倾向于押注那些盈利驱动因素更独立、与AI相关性更低的资产。

4. 一个更实用的观察框架：区分“周期性的AI”与“结构性的盈利”

基于报告的框架，投资者可以建立一个更清晰的分析框架。当前市场面临的核心问题是：如何区分哪些盈利增长是周期性的AI驱动，哪些是结构性的公司基本面改善？

报告的十项主题回顾中，前五项聚焦于AI与电力、代理经济、数据中心建设等热门话题。但值得注意的是，报告团队在“半年度更新”中特意将十项主题压缩为五项，理由是“在一个‘单一交易’的市场中，投资者可能不需要十个主题来关注”。这本身就是一种信号：市场的注意力集中度已经达到了需要警惕的水平。

\subsection{[R020] 自动驾驶的真正拐点，不在于技术，而在于运营的规模经济学}
{\small\color{refgray}归类：未被主题摘要引用}

当市场还在争论“特斯拉的FSD何时能实现L4”或“Waymo的传感器成本何时降到可商业化”时，一份某外资投行的最新研报给出了一个更务实的判断框架：自动驾驶赛道的竞争，正在从“谁的技术更酷”转向“谁能把安全数据转化为运营规模，再把运营规模转化为用户习惯”。

这份研报没有停留在技术参数对比，而是用一组关键运营数据——事故频率、车队规模、月活用户、用户使用频次——构建了一个更接近产业现实的坐标系。在这个坐标系里，一些被市场忽视的信号正在浮现。

报告的核心洞察可以浓缩为一句话：行业正在经历从“技术验证”到“运营验证”的切换。这个切换过程中，安全数据的稳定性、车队的实际运营规模、以及用户的使用粘性，比任何技术路线图都更能决定哪家公司能率先跨越商业化的鸿沟。

1. 安全数据正在成为最硬的竞争壁垒，但可比性远没有市场想象的那么高

研报基于美国国家公路交通安全管理局的事故数据，以及Waymo和特斯拉各自披露的运营里程，估算出两家公司的每英里事故率：特斯拉大约每10万至12万英里发生一次事故，Waymo大约每15万至17.5万英里发生一次事故。

乍看之下，Waymo的事故率更低。但研报谨慎地指出，这个对比并不直接。特斯拉在奥斯汀的车队是混合状态——部分车辆有安全监控员，部分没有。Waymo则在更多城市进行商业化运营，且报告的事故包括那些人类驾驶员通常不会上报的轻微事件，比如压过路沿。

这里有一个更值得关注的信号：特斯拉在3月份实现了约20万英里才发生一次事故，并且在8月、2月和4月上半月都没有报告任何事故。Waymo在4月前半月的约1000万英里运营里程中，仅报告了2起事故。两家公司的安全表现都在快速提升，但波动性依然存在。

对于投资者而言，真正需要追问的不是“谁的事故率更低”，而是“谁能在扩大运营规模的同时，保持事故率的持续下降”。安全数据不是静态的比较，而是动态的压力测试。当Waymo的车辆数从目前的约3800辆扩大到数万辆时，每英里事故率能否维持在当前水平？当特斯拉的无人驾驶车队从42辆（得克萨斯州注册数据）扩展到更多城市时，其事故分布是否会发生变化？这些才是决定估值溢价的真正变量。

2. 车队规模的真相：注册数不等于运营数，运营数不等于有效供给

研报披露了一组容易被忽视的硬数据：根据得克萨斯州交通局和NHTSA的备案文件，Waymo的车队规模约为3800辆（其中577辆在得克萨斯州），而特斯拉在得克萨斯州的车队仅为42辆。

这个数字对比本身就说明了问题。特斯拉的Robotaxi服务虽然在奥斯汀、达拉斯和休斯顿已经以无人驾驶模式运营，但实际可用的车辆数量极为有限。媒体报道也佐证了这一点：在已开通无人驾驶服务的城市，车辆的可获得性仍然受限。

Waymo的3800辆也并非全部在运营。注册车辆可能包括测试车辆、备用车辆，以及配备安全监控员的车辆。但即便如此，其运营规模与特斯拉之间的差距是数量级的。

这引出了一个更深层的问题：自动驾驶的商业化，本质上是规模经济学。没有足够数量的车辆，就无法积累足够多的里程数据来迭代算法；没有足够多的里程数据，就无法证明安全性的稳定性；没有安全性的稳定性，就无法获得监管批准以扩大车队。这是一个正反馈循环，而起点是“把车造出来、开出去”。

Waymo在亚利桑那州的工厂正在向年产数万辆的产能迈进。特斯拉则计划在2026年底前将FSD v15版本推向市场，CEO表示这个版本可能是规模化扩张的关键。两家的路径截然不同：Waymo走的是“先有规模，再优化成本”，特斯拉走的是“先有技术，再扩大规模”。哪种路径更有效，目前尚无定论，但研报的数据显示，Waymo在“规模”这个维度上领先了一个身位。

3. 用户数据的悖论：月活增长强劲，但使用频次远低于传统网约车

SensorTower的数据揭示了一个有趣的现象：Waymo的月活跃用户在过去一年增长了20\%，占Uber美国月活用户的比例从年初的2\%上升到了约4\%。这个渗透率的提升是稳步的，而不是爆发式的。

但与此同时，Waymo用户的月均使用次数仍然显著低于Uber和Lyft。从研报提供的图表来看，Uber用户的月均会话次数稳定在23-25次，Lyft在19-21次，而Waymo只有7-12次。

这意味着什么？用户愿意尝试自动驾驶，但还没有形成高频使用的习惯。可能有几个原因：一是服务覆盖范围有限，用户只在特定区域使用；二是等待时间或接驾效率不如传统网约车；三是用户对自动驾驶的信任还需要时间来积累。

对于平台型公司而言，这个数据既是机会也是挑战。机会在于，一旦用户习惯了自动驾驶，使用频次有可能向传统网约车靠拢，这意味着巨大的增量空间。挑战在于，如果频次长期偏低，单位经济模型将难以跑通——每辆车每天的接单量不足，固定成本无法被摊薄。

这里有一个尚未被充分讨论的二阶效应：如果Waymo的用户频次始终低于Uber，那么对于Uber和Lyft来说，自动驾驶可能不是颠覆性的威胁，而是补充性的供给。Uber和Lyft可以通过接入Waymo、Mobileye等第三方自动驾驶车队，在不增加自有车辆的情况下扩大供给，同时保持平台的流量优势。这就是Uber和Lyft目前正在做的事情——它们把自己定位为“自动驾驶车队的需求聚合平台”，而不是自动驾驶技术的开发者。

\subsection{[R021] 硬件市场真正稀缺的不是需求，而是被低估的供给议价权}
{\small\color{refgray}归类：AI与半导体：需求扩散至多节点，供给约束成为新定价变量}

过去一个月，某外资投行在台湾与主要OEM、ODM及供应链厂商密集会面后，得出一个与市场共识存在显著温差的核心判断：硬件行业当前最大的结构性变量，并非需求端的波动，而是供给侧正在发生的、被多数投资者低估的定价权转移。这份行程纪要揭示了一个正在加速分化的市场格局——不同品类、不同厂商之间，议价能力的差距正在从量变走向质变。

报告最值得关注的判断在于：HDD（硬盘）正在成为整个硬件产业链中最具定价权的品类，其供需缺口和价格弹性远超市场预期；而PC则在真实需求走弱与记忆体成本上行的双重挤压下，面临利润率侵蚀的实质性风险。在这两个极端之间，服务器是唯一需求趋势相对清晰、但同样面临拉货前移和利润率压力的市场。

这一判断之所以重要，是因为它挑战了当前市场对硬件板块“普涨普跌”的认知框架。投资者习惯于用终端需求强弱来线性推演厂商表现，但报告显示，真正决定未来12-18个月利润走势的，是各厂商在供应链中的议价地位——谁能在稀缺中拿到更优的供应条件，谁就能在价格上行周期中实现超额收益。

1. HDD供需缺口正在加速扩大，定价权从买家转向卖家的拐点已经到来

这可能是整份报告中最具冲击力的数据点。三周前，分析师团队还认为HDD行业在CY26的供需缺口约为200EB，CY27约为250EB。而仅仅三周后的台湾之行后，这一数字被大幅上修：CY26缺口扩大至300EB，CY27和CY28均达到400EB。驱动这一变化的，是AI推理工作负载对通用服务器出货量的超预期拉动，进而带动对大容量HDD的需求激增。

更关键的是定价信号。当前大容量HDD的混合均价约为15美元/TB（1.5美分/GB），而供应商正在将CY27的目标价定在25美元/TB（2.5美分/GB），CY28进一步上调至30美元/TB（3美分/GB）。这意味着从当前到CY28，HDD价格有望实现近乎翻倍的增长。更为重要的是，超大规模客户对此几乎没有议价抵抗——“他们只是要求更多的供应”。这一现象表明，HDD的定价权已经从买方转移至卖方。

从供给端看，行业没有新增的绿地产能规划，供应增长几乎完全依赖HAMR（热辅助磁记录）技术的良率爬坡和产品组合优化。这意味着，未来几年的供给弹性极为有限。报告甚至指出，行业普遍认为HDD供应商有望在CY27实现超过60\%的毛利率，CY28可能触及70\%以上。对于一家成熟硬件公司而言，这样的利润率水平在过去十年中几乎不可想象。

这一判断对投资者的含义是：HDD正在从“周期性大宗商品”转变为“结构性稀缺资产”。那些拥有HDD供应敞口的公司，尤其是Seagate和Western Digital，其盈利能力和估值逻辑需要被重新审视。市场目前尚未充分定价这一结构性变化。

2. PC需求正在实时走弱，而记忆体成本上行正在侵蚀本已微薄的利润

与HDD的强势形成鲜明对比的是PC市场。报告传达的信号相当清晰：需求正在实时恶化，且利润率风险将在2026年下半年和2027年集中暴露。

一个ODM的出货指引变化尤为说明问题。正常情况下，笔记本出货的上下半年比例为45\%比55\%，但当前指引显示2026年上半年可能占到60\%，下半年仅40\%。这意味着大量需求已被提前拉货，下半年面临的是真实需求不足和渠道库存积压的双重压力。一家PC OEM甚至表示，过去四周内消费端和中小企业销售出现了两位数的同比下滑，欧洲、中东和非洲市场尤为明显。

与此同时，记忆体成本的持续上行正在压缩PC厂商的利润空间。报告明确指出，低端PC几乎无法销售，因为记忆体价格上涨已使该价位段的产品陷入负毛利。高端PC尽管需求相对坚挺，但其背后的逻辑并非消费升级，而是低端市场被成本挤出后的被动集中。

值得注意的是，不同OEM在记忆体采购条件上的差异正在成为利润分化的关键变量。按每比特成本从低到高

[... middle omitted ...]

力以及规模优势，正在“变现其记忆体关系”，成为PC板块中相对防御性最强的标的。

服务器是三个主要硬件品类中需求趋势最为清晰的，但报告揭示的细节远比“AI驱动增长”这个笼统叙事要复杂。

通用服务器方面，需求强劲的主要驱动力来自两个方面：一是去年被压抑的需求释放，二是对价格进一步上涨的担忧引发的拉货前移。ODM普遍反映，企业级CPU服务器订单在3月出现了显著跳升，其中DELL的通用服务器全年增速有望达到20\%。但报告同时强调，目前尚无法判断拉货前移的幅度和持续时间——这是一个尚未被充分量化的风险变量。

AI服务器方面，供应链对2027年以后的增长前景持积极态度，但实时可见性有限。Vera Rubin平台按计划在2026年下半年推出，DELL有望成为首批客户之一，但初期出货量极小。更重要的是，报告指出AI服务器机架设计正在向L10层级标准化演进，这可能导致OEM/ODM的利润率承压。虽然时间表尚不确定，但标准化趋势对OEM的负面影响正在被越来越多的供应链参与者所确认。

\subsection{[R022] 市场真正低估的不是需求，而是供给侧的再定价}
{\small\color{refgray}归类：AI与半导体：需求扩散至多节点，供给约束成为新定价变量}

过去几周，围绕AI芯片的叙事被两种声音主导：一种是关于“需求是否见顶”的争论，另一种是关于“产品延迟是否意味着周期拐点”的担忧。某外资投行刚刚发布的台湾供应链调研，给出了一组值得认真对待的反直觉信号。

调研的核心判断是：需求侧没有出现任何值得担忧的转折，但供给侧的结构性变化——从产品迭代节奏到竞争格局的重新划分——正在重塑整个半导体投资框架。市场当前定价中，隐含了太多对“需求放缓”的焦虑，却远远没有反映供给端正在发生的、具有持久影响力的再定价。

这份报告最值得关注的，不是它重申了对某家头部GPU公司的看好，而是它揭示了三个被普遍误解的关键变量：产品周期的节奏其实比市场感知的更健康、CPU市场的竞争格局正在被一家非传统玩家改写、以及AI基础设施的瓶颈已经从芯片本身转移到了存储。

最近一段时间，关于某头部GPU公司下一代产品Rubin“重新流片”或“延迟超预期”的传闻在市场上流传甚广。这些故事听起来符合半导体行业过往的“产品跳票”剧本，也迎合了部分投资者对AI投资热潮过热的警惕心理。

但报告通过对供应链多个节点的交叉验证，给出了完全相反的结论：Rubin的研发节奏完全符合既定时间表。此前确实存在几周的实施延迟，但这对芯片本身的量产节奏几乎没有影响。报告特别强调，没有看到任何证据表明Rubin会出现类似Grace Blackwell时期那种机架级别的量产挑战。相反，芯片性能表现强劲，尽管首代产品在优化上仍有空间，但这属于正常的技术迭代过程。

更值得关注的是，上一代产品Blackwell的生产不仅没有放缓，反而出现“再加速”。其生命周期被延长了几个月，产量也高于最初预期。报告将此解读为需求的真实体现——当Rubin在初期供给受限时，客户对即将成为“上一代”的Blackwell仍有强劲需求。

这意味着什么？市场习惯用“延迟”来推断需求疲软，但这里的链条是反过来的：正因为需求太过强劲，客户愿意为现有产品支付溢价，才使得供应链有动力和空间去延长成熟产品的生命周期。这不是周期见顶的信号，而是需求持续超配的证明。

这家GPU公司在财报电话会上抛出了一个让市场普遍感到意外的数字：今年CPU业务收入将达到200亿美元。这个数字如果属实，意味着它已经成为CPU市场的领导者之一。但市场的第一反应是怀疑——很多人认为这个数字一定包含了存储或其他产品的收入。

报告通过供应链调研，给出了一个令人信服的拆分：其中约100亿至120亿美元来自“头节点”（head node）——这部分CPU与GPU、NVSwitch和存储打包销售，定价本身难以拆分。但剩下的70亿至100亿美元，则来自独立的Grace或Vera CPU销售。这个数字已经通过多个独立信息源得到交叉验证。

更关键的是，报告指出，目前市场上除了这家公司之外，几乎没有足够的商用ARM服务器CPU供给。Grace和Vera的性能已经得到验证，而供给的稀缺性使其成为客户在CPU短缺背景下的自然选择。

这里需要特别关注一个二阶效应：CPU市场正在经历一场由“供给短缺”引发的结构性变化。报告提到，多家服务器部署中出现了“内存少于理想配置”的情况，这指向一个比DRAM短缺更隐蔽但同样重要的CPU短缺。而这家GPU公司正在利用自己的ARM CPU设计能力和晶圆供给优势，填补这个空白。

这对传统x86厂商意味着什么？报告给出一个意味深长的提醒：CPU厂商在涨价时必须小心，过度的价格压力可能会把更多业务推向ARM生态。

3. Intel的18A进展：一个“技术上可行、战略上存疑”的转折点

Panther Lake芯片在18A节点上的良率持续改善，已经达到“可靠”水平——尽管尚未达到经济目标。这无疑是Intel在先进制程上取得的一个技术里程碑。但报告同时指出，更关键的问题在于：这个

[... middle omitted ...]

权价值”。但有趣的是，由于当前服务器CPU存在普遍短缺，Intel有足够的能力在短期内“超出预期”——因为即便产品竞争力没有根本性提升，只要供给存在，就能拿到订单。

这意味着什么？Intel的“beat and raise”短期叙事是成立的，但长期竞争格局并没有因此变得清晰。市场对Intel的定价中，可能高估了代工业务带来的估值弹性，而低估了其核心CPU业务面临的竞争压力。

在x86服务器CPU市场，AMD正在执行一个被市场低估的竞争策略：它不是通过降价来抢夺份额，而是通过“涨得比对手少”来获得价格优势。报告指出，AMD的服务器CPU实现了同等产品上的价格上涨，但涨幅小于Intel和ARM替代方案。在客户对成本敏感的背景下，这种“相对便宜”本身就是一种竞争力。

即将在2026年下半年量产的2纳米产品Venice，虽然价格昂贵但性能非常强劲。更重要的是，5纳米和7纳米产品仍然具有竞争力。报告提到，大部分2026年的供给已经分配完毕，客户正在为明年的需求建立可见性。

\subsection{[R023] 市场低估的不是需求，而是供给纪律的自我强化}
{\small\color{refgray}归类：能源与大宗：供给恢复慢于预期，锂矿重启信号密集}

过去几年，市场对海上钻井的讨论始终绕不开一个核心焦虑：一旦日费率回升，船东会像上一个周期那样大规模重启冷停堆、甚至下单新船，把刚刚收紧的供需再次推向过剩。这种“供给侧必将失控”的假设，构成了投资者对海上板块长期保守定价的底层逻辑。

但某外资投行最新发布的北美能源服务与设备深度报告，给出了一个值得认真对待的反向判断。报告的核心推演并不复杂，但结论的分量在于：支撑这一轮海上钻井复苏的结构性因素，可能比市场目前定价的要更可持续。真正重要的不是需求增长有多快，而是供给侧的纪律能否被保持——而报告认为，当前的市场结构正在让这种纪律自我强化。

这份报告覆盖了从深水浮式钻井平台到自升式平台的全面供需模型，时间跨度延伸到2030年。我们从中提炼出五个关键判断，它们共同指向一个结论：海上钻井可能正在进入一个“低供给弹性、高定价权”的新阶段，而市场对这一结构性变化的定价仍不充分。

1. 深水需求正在加速，但真正支撑定价的是供给端的“不响应”

报告的供需预测是明确的：到2030年，全球海上钻井平台需求将增长约12\%，其中深水部分增长17\%，大陆架部分增长10\%。而同期供给端几乎没有新增——报告假设2027至2030年间既无冷停堆重启，也无新船交付。

这个“供给零增长”假设是整份报告最值得关注的判断基础。它不是基于对船东意愿的乐观猜测，而是基于对现有船队结构的分析：冷停堆的高规格浮式钻井平台集中掌握在少数几家公开上市的钻井承包商手中，而这些公司在经历上一轮周期惨痛教训后，对资本纪律的执行力显著提升。

这里的关键洞察是：市场可能低估了供给端行为模式的改变。过去十年，海上钻井行业经历了大规模的资产减记、破产重组和船队退役。幸存下来的船东不仅规模更小、资产负债表更干净，而且管理层对“为市场份额而牺牲定价”的策略已经彻底失去兴趣。这不是一次周期性的克制，而可能是结构性的行为转变。

报告预测，到2028年，海上钻井平台的市场化利用率将升至85\%左右的中高水平。在历史经验中，这恰恰是船东获得显著定价权的阈值区间。

更值得关注的不是绝对数字，而是达到这一水平的速度和路径。报告指出，当前深水浮式钻井平台的合同授予量已经达到近15年来的最高水平，合同前置时间和合同期限均高于正常水平。这些指标通常在市场收紧的早期阶段出现，意味着需求端的动能还没有完全反映在日费率中。

报告预计，高规格深水钻井平台的领先日费率将在未来1至2年内恢复到2023年峰值水平，随后继续上行。具体而言，6G+级钻井船和半潜式平台的日费率到2028年将分别上升约20\%至52.5万美元和47.5万美元，高规格自升式平台到2030年将上升10\%至15\%至17.5万美元。

这些数字本身并不惊人，但结合供给端的零增长假设，它们暗示了一个更重要的含义：定价权正在从船东的“供给管理”转向市场的“刚性需求”。当供给无法响应需求增长时，日费率的上涨将不再是周期性的，而可能具有更强的持续性。

报告特别指出了区域增长的显著分化。全球海上资本支出增长最快的地区集中在非洲、拉美（除巴西外）和亚太。在这些市场，报告预计2027年资本支出同比增长15\%，2026至2030年复合年增长率达到5\%至10\%。

这一判断对理解行业格局的意义在于：增长并非均匀分布，而是高度集中在那些对能源安全、供应多元化和战略储备建设有迫切需求的国家和地区。这意味着，能够在这些市场建立长期客户关系和运营网络的服务商和承包商，将获得比行业平均更强劲的增长曲线。

报告提到的几个关键信号值得关注：深水活动的早期拐点已经出现，活跃钻机数量比2025年底的低点增加了7\%；合同期限的延长说明运营商正在为长期产能锁定而提前签约；而过去几年大量老旧船队的退役，使得剩余的可动用船队质量更高、效率更强。

这些信号共同指向一个判断：

[... middle omitted ...]

认为，当前效率提升的影响正在被更强劲的宏观需求和战略储备建设的需求所“压倒”，但这一判断的验证需要时间。如果未来几年全球石油需求增长放缓，或者OPEC+的产量政策发生变化，效率提升对钻机需求的压制效应可能会重新显现。

另一个未完全展开的问题是：冷停堆重启的成本和可行性。虽然报告认为冷停堆集中在少数几家纪律严格的船东手中，但重启一座冷停堆钻井平台的实际成本、时间周期和技术风险，在不同船型和不同状态下差异巨大。报告没有详细拆解这一变量，而它恰恰是判断供给纪律能否真正维持的关键。

综合报告的判断，一个更重要的框架性启示是：对于海上钻井和相关服务板块的投资逻辑，可能需要从传统的“周期择时”转向“结构定价”。

传统上，投资者关注的是“现在处于周期的什么位置”，然后据此判断是买入还是卖出。但报告暗示，这一轮可能不是简单的周期重复，而是行业结构发生了根本性变化：供给侧的弹性大幅降低，船东的资本纪律从被迫转向主动，而需求侧的增长则受到能源安全和战略储备建设等非周期性因素的支撑。

\subsection{[R024] 世界杯与夏季旅行正成为消费分化的新分水岭，而非总量信号}
{\small\color{refgray}归类：风险偏好与资金流向：机构与散户出现显著分歧}

美国消费者正站在一个看似矛盾的路口。一方面，通胀担忧在最新一轮调查中回升至59\%，较1月低点明显反弹；另一方面，62\%的消费者计划在夏季出行，净支出意愿高达+32\%，与过去两年持平。如果只看总量数字，很容易得出“消费韧性仍在”的结论。但某外资投行最新发布的AlphaWise消费者调查揭示了更深层的结构性信号：真正重要的不是消费总量是否增长，而是增长从何而来、由谁驱动、以及哪些品类正在被重新定价。

这份覆盖约2000名美国消费者的第76轮调查，捕捉到了两个值得产业决策者高度关注的变化。第一，2026年世界杯正在创造一个新的消费场景，44\%的消费者计划参与相关活动，其中70\%预计会增加支出。第二，夏季出行计划的结构正在发生微妙但重要的迁移——露营意愿从9\%跃升至16\%，而邮轮从15\%回落至8\%。这些变化叠加在一起，指向一个核心判断：美国消费市场正在从“总量驱动”转向“事件驱动”和“体验分层”，而企业能否识别并捕获这些结构性拐点，将决定未来12个月的竞争位势。

1. 世界杯消费不是一次性脉冲，而是测试品牌能否激活“轻参与”场景的试金石

调查显示，44\%的消费者计划以某种形式参与世界杯，但参与方式高度集中在“轻参与”场景：30\%通过电视或流媒体观看，23\%通过社交媒体追踪。真正重度参与——如现场观赛（4\%）、参加观赛派对（8\%）或举办聚会（6\%）——的比例远低于预期。

这意味着什么？世界杯对消费的拉动不会以传统“超级碗式”的集中爆发呈现，而是分散在多个轻量级触点中。32\%的参与者预计会增加食品和非酒精饮料支出，28\%会增加外卖或配送支出，另有28\%可能订阅新的流媒体服务。这些品类恰恰是过去两年通胀环境下消费者压缩最明显的领域。世界杯提供了一个“合理化额外支出”的理由，让消费者在已经紧绷的预算中为体验性消费腾出空间。

对于消费品和零售企业而言，关键不在于押注大型营销活动，而在于能否在“轻参与”场景中嵌入购买闭环。例如，流媒体平台能否通过赛事内容提高订阅转化率，外卖平台能否通过观赛套餐提升客单价，食品品牌能否通过世界杯主题产品创造溢价空间。调查数据中一个值得注意的细节是，21\%的参与者表示会增加体育博彩支出——这可能是调查尚未充分展开的、对数字支付和金融科技公司的潜在利好。

2. 夏季出行意愿的结构性迁移，暴露了消费者在预算约束下的优先级重排

62\%的消费者计划夏季出行，这一比例略高于去年的60\%，但结构变化远比总量有意义。公路旅行占比从39\%升至42\%，露营从9\%跃升至16\%，而邮轮从15\%降至8\%。这组数据不只是一个出行方式偏好的波动，它反映了消费者在通胀压力下对“性价比体验”的重新定义。

露营的爆发式增长尤其值得注意。这一类别在2024年仅占13\%，2025年跌至9\%，2026年却翻倍至16\%。结合汽油价格持续高企的背景，消费者并非在减少出行，而是在选择成本更低、可控性更高的出行方式。公路旅行和露营本质上是对“固定成本出行”的替代——消费者仍然渴望旅行体验，但通过降低住宿和交通的边际成本来维持预算平衡。

这对旅游和休闲行业的启示是分层的。高端酒店和航空公司可能继续受益于高收入群体的刚性需求——调查显示，年收入15万美元以上家庭的出行意愿达78\%，显著高于低收入群体。但中端市场正在被露营、公路旅行等替代方案侵蚀。对于邮轮行业，8\%的意愿虽然低于2025年的15\%，但仍高于2024年的6\%，说明邮轮需求并未消失，只是波动性加大。关键在于，邮轮公司能否通过定价策略和航线设计，重新捕获那些转向露营的“性价比敏感型”消费者。

3. 净支出意愿+32\%看似稳健，但真正的风险隐藏在“谁在驱动”这个维度

45\%的旅行者计划增加支出，13\%计划减少，净支出意愿+32\%——这一数字与过去两年持平。但如果只看这个数字

[... middle omitted ...]

的投资逻辑需要进一步细分。依赖中低收入消费者高频消费的零售和休闲品牌，可能面临比总量数据更严峻的挑战。而定位高收入群体的奢侈品、高端体验和差异化服务，则可能继续受益于“K型消费”的深化。

4. 报告尚未完全回答的关键问题：世界杯消费的真实增量到底有多大？

调查显示70\%的参与者预计会增加支出，但这里的“增加”是相对于什么基准？是相对于日常消费，还是相对于原本就计划进行的消费？如果是前者，世界杯确实创造了增量；如果是后者，更多是消费的重新分配。报告没有区分“净新增消费”和“消费转移”，这是理解世界杯经济效应的核心盲点。

另一个未解的问题是，世界杯的消费刺激能否持续到赛事结束后。调查集中在5月下旬，此时距世界杯开幕（6月11日）仅数周，消费者的预期可能包含一定的“赛事前兴奋”成分。如果实际体验低于预期，或赛事期间出现意外事件（如天气、地缘政治等），消费落地可能低于调查数据。这些不确定性意味着，世界杯对相关股票的催化作用可能更多体现在短期情绪层面，而非长期基本面改善。

\subsection{[R025] 半导体CEO们的共识：AI需求正在重塑WFE定价权，但真正的约束是物理空间}
{\small\color{refgray}归类：AI与半导体：需求扩散至多节点，供给约束成为新定价变量}

半导体CEO们的共识：AI需求正在重塑WFE定价权，但真正的约束是物理空间

当市场还在争论半导体周期何时见顶时，Bernstein 2026年战略决策会议上五位半导体CEO的发言传递了一个更值得关注的信号：AI对半导体设备支出的拉动已从“预期”进入“执行验证期”，而行业真正的瓶颈不再是需求，而是洁净室产能。

这份研报的核心价值不在于它给出了多少数字，而在于它让市场听到了产业一线的声音——Lam Research、Applied Materials、Qualcomm、Texas Instruments和Entegris的CEO们，在同一个舞台上，几乎都指向了同一个方向：AI正在让WFE的长期可见性变得前所未有地清晰，但供给侧的物理约束正在成为新的定价变量。

1. AI对WFE的拉动已从“训练驱动”切换到“推理+存储+封装”的多波次共振

过去两年，市场对AI拉动半导体设备的叙事主要集中在训练侧对先进逻辑芯片的需求。但Lam Research CEO Tim Archer在对话中清晰地描述了需求波次的演进：AI训练需求首先拉动了先进Foundry/Logic，而AI推理的爆发正在驱动NAND和CPU需求的上升。这意味着WFE的驱动力正在从单一节点向多节点扩散。

Applied Materials CEO Gary Dickerson提供了更量化的佐证：AI驱动的先进逻辑、DRAM（HBM）和先进封装，预计将贡献2026年增量WFE支出的80\%。这不是一个“如果AI持续增长”的假设，而是一个基于客户滚动8个季度预测的既成事实。

这里的关键洞察是：AI对半导体的需求不再是“一波流”，而是形成了训练-推理-存储-封装的递进式共振。这拉长了WFE上行周期的持续时间，也增加了需求的“厚度”。对于设备厂商而言，这意味着他们不再需要押注单一技术节点，而是可以在多个技术方向上同时获得订单。

2. 洁净室产能成为WFE增长的硬约束，这意味着设备商的议价能力正在结构性提升

Lam Research和Applied Materials都提到了同一个瓶颈：洁净室空间。Lam明确表示，洁净室空间是当前WFE最大的约束，2027年也将如此。Applied Materials虽然没有给出WFE预测，但客户提供的滚动8个季度预测本身就说明需求远超当前产能规划。

这是报告中最容易被低估的信息。洁净室不是可以快速扩张的产能——从规划到投产通常需要2-3年。当需求增长快于洁净室扩张速度时，设备商的订单可见性会显著延长，同时设备本身的定价权也会增强。

Applied Materials CEO已经明确表示，公司看到“利润池比以往任何时候都大”，并且正在通过定价来捕捉其为客户创造的价值。Lam Research的毛利率已经进入“50s neighborhood”，而且是在中国业务占比下降的情况下实现的——这说明定价能力的提升不是靠区域套利，而是靠技术价值。

对于投资者来说，这意味着设备商的利润率可能比市场预期的更有韧性，甚至可能随着技术节点升级而持续改善。

3. 中国业务不再是增长引擎，但也不再是拖累——设备商正在适应“去中国化”的新常态

中国业务是过去几年半导体设备股波动的主要来源之一。但这次CEO们的表态显示，中国正在从一个“不确定变量”变成一个“可预测常量”。

Lam Research预计中国业务今年持平或略有增长，因为公司只能在中国参与成熟节点的竞争。Applied Materials的中国业务占比约为中20\%左右，主要来自ICAPS（物联网、通信、汽车、电源和传感器），预计持平或微增。Entegris的情况更特殊——中国限制主要针对大型设备，材料领域的限制相对宽松，公司甚至可以通过本地化供应来服务中国市场。

这里的含义

[... middle omitted ...]

nstruments CEO Haviv Ilan的发言，展示了AI需求如何从设备端向下游芯片设计公司和模拟芯片公司传导。

Qualcomm的核心叙事是“AI手机升级周期”。CEO认为当前手机出货受到内存短缺的人为压制，但AI代理对算力的需求远超当前手机能力，因此换机周期是确定的。同时，数据中心业务预计在FY27贡献数十亿美元收入，且对运营利润率有正向贡献。

TI的故事则更接地气。数据中心业务虽然只占收入的12\%，但同比增长90\%+，已经成为一个20亿美元年化收入的业务。TI认为自己在电源、传感、冷却和保护等领域的广泛产品组合，使其在数据中心TAM（约125亿美元）中具备份额增长空间。

这两个案例说明：AI对半导体的拉动正在从“设备投资”向“芯片出货”传导，但路径不同。Qualcomm依赖的是终端设备的AI化升级，TI依赖的是数据中心基础设施的电力与热管理需求。对于投资者而言，这意味着AI的受益面正在扩大，但每个细分赛道的驱动力和竞争格局差异较大，不能一概而论。

\subsection{[R026] 肺癌治疗正在经历一场“分子分型”驱动的定价重构}
{\small\color{refgray}归类：医疗健康：双抗改写肺癌标准治疗，口服GLP-1竞赛参数被重新定义}

市场对肿瘤药的投资叙事，长期被“PD-1之后下一个是什么”的宏大问题主导。但某外资投行在ASCO 2026现场发布的KOL深度访谈纪要，揭示了一个更具体、也更值得投资者关注的信号：晚期肺癌正在向慢性病转化，而真正改变竞争格局的力量，不是某个单一靶点或机制，而是分子分型在早期干预、耐药管理和联合治疗三个维度上同时加速落地。

这份报告的价值不在于复述数据，而在于一位顶级临床医生对每项关键数据的“信任度评分”——哪些信号足以改变处方行为，哪些只是统计学上的显著性，背后是临床实践的真实门槛。以下是我们从中提炼的五个核心洞察。

1. RET辅助治疗的数据令人震撼，但真正的瓶颈不在药物，在检测

LIBRETTO-432研究中，selpercatinib在RET阳性早期非小细胞肺癌辅助治疗中取得了0.17的风险比——在肿瘤辅助治疗领域，这是一个几乎“好到不真实”的信号。报告中的KOL明确指出，这是ASCO周末最令人信服的数据之一，延续了ADAURA在EGFR突变和ALINA在ALK阳性疾病中建立的模式。

但这里有一个被市场普遍低估的结构性障碍：约70\%的RET阳性患者是从不吸烟者，他们不在现行肺癌筛查标准之内。这意味着，即使药物疗效再强，如果患者无法被识别，商业潜力也就被天然锁死。KOL估计，大约30\%的RET阳性患者在1B-3A期被诊断，但这个数字可能因其所在的纽约执业环境而偏高。

所以呢？对于投资者而言，selpercatinib的峰值销售预测，不应仅基于疗效数据，而必须纳入分子检测渗透率这一硬约束。目前早期肺癌的液体活检可靠性低于晚期，组织检测仍是金标准，但即使在资源充足的机构，组织获取和分子分型也远未普及。这给所有依赖生物标志物驱动的辅助治疗药物，设定了一个共同的“天花板”。

2. HARMONi-6的OS数据是积极的，但“中国单中心”的全球可平移性仍是关键悬疑

Ivonescimab在鳞状非小细胞肺癌中给出了0.66的总生存风险比，且出血风险可控——这比市场此前对VEGF抑制剂在鳞癌中安全性的担忧要好。但KOL的评论值得反复咀嚼：他不会根据一个风险比数字来决定处方行为。

原因有三。第一，这是一项中国单中心研究，患者年龄上限为75岁，入排标准的细节以描述性而非分析性方式呈现，KOL明确表示中国人群可能与西方临床实践不完全对应。第二，鳞癌患者中超过一半年龄在65岁以上，这部分人群的VEGF抑制风险-获益计算需要个体化评估。第三，真正的全球验证数据——HARMONi-3的OS结果——预计要到2026年下半年才能读出。

这里有一个更深的投资含义：市场对Summit Therapeutics的定价，实际上已经隐含了HARMONi-3成功的预期。但KOL的评论暗示，即使HARMONi-3达到统计学显著性，临床接纳仍将是一个渐进过程，而非开关式切换。FDA批准和全球数据的一致性，是药物获得临床地位的“两个必要条件”，而非充分条件。

3. 洛拉替尼的CROWN数据设定了极高基准，但Nuvalent的路径比市场想象的要清晰

CROWN研究的7年更新数据令人震撼：中位无进展生存期仍未达到。KOL回忆，2000年他刚到美国时，ALK阳性患者的生存期只有6-7个月。现在，这些患者可以活7-10年，肺癌正在接近慢性病范式。

但正是这个“极高基准”，反而为Nuvalent的第四代ALK抑制剂创造了清晰的市场切入点。KOL指出，并非所有患者都是洛拉替尼的候选者——神经毒性和合并症使部分患者无法使用。Nuvalent药物在洛拉替尼耐药患者中约26\%的响应率，已经定义了明确的未满足需求。

更值得关注的是KOL对临床惯性的描述：从第二代阿来替尼转向第三代洛拉替尼，花了数年时间才克服文化惯性。这意味着，即使Nuvalent的数

[... middle omitted ...]

”。Revolution Medicines在这一领域的布局正在加速，尽管完整数据仍需等待。

这里有一个对投资框架的含义：KRAS领域的竞争将从二线治疗快速前移至一线，而一线市场的定价权和商业规模，与二线完全不在一个量级。能够在一线建立证据的药物，将获得显著的先发优势。但这也意味着，任何在二线看似有竞争力的数据，都可能被一线格局的变化所颠覆。

5. Sac-TMT的数据令人信服，但“全人群vs生物标志物”的哲学冲突尚未解决

Sacituzumab tirumotecan在TROP2 ADC中的表现，被KOL认为是目前非小细胞肺癌中最令人信服的TROP2药物。但KOL作为生物标志物研究者的背景，让他对“全人群策略”感到哲学上的冲突。

这个冲突的实质是：没有伴随诊断的ADC，其临床定位和商业推广将更加依赖医生对“毒性-疗效比”的个体化判断。KOL指出，ADC的毒性相对于化疗很难一概而论——患者反应差异很大。临床试验常规排除体能状态差的患者，这限制了真实世界的外推。

\subsection{[R027] 市场低估的不是Broadcom的短期业绩，而是其2027年XPU份额的再定价空间}
{\small\color{refgray}归类：AI与半导体：需求扩散至多节点，供给约束成为新定价变量}

市场低估的不是Broadcom的短期业绩，而是其2027年XPU份额的再定价空间

市场对Broadcom的认知，正在经历一次关键的框架切换。过去几个季度，投资者习惯于围绕“AI收入是否超预期”来交易这只股票。但某外资投行在最新一期半导体周报中，提出了一个值得深究的信号：2026年下半年的业绩曲线可能不如市场期待的那样陡峭，但真正决定公司估值中枢的，是2027年XPU（自定义芯片）的份额增长，而非当前季度的beat幅度。

这份报告的判断，指向一个容易被忽略的结论：Broadcom当前约10\%的AI市场份额，才是未来两年最值得被重新定价的变量。市场如果只盯着季度环比增速的起伏，就会错过一个结构性叙事正在成熟的事实。

报告最引人注目的调整，是对2026年下半年AI收入节奏的重新判断。分析师明确表示，由于机架部署相关的供应链问题，原本预期集中在2026年下半年的约210亿美元机会，大部分将推迟至2027年。这直接导致其对7月季度AI收入的预测从此前的195亿美元下调至169亿美元，环比增速从80\%降至56\%。

但关键在于，报告强调这并非需求消失。订单没有取消，产能被重新分配给其他TPU客户。这意味着，Broadcom的客户基础正在扩大，而不仅仅是单一超级客户的集中采购。从框架角度看，这不是一个坏消息——它说明Broadcom在TPU生态中的角色正在从“单一供应商”向“平台型产能分配者”演进。

对于投资者而言，这意味着2026年下半年的业绩可能不会出现市场此前预期的“超级爆发”，但2027年的收入基数将被显著抬高。如果市场因为2H26增速放缓而下调估值，反而可能是一个重新入场的机会。

2. 2027年的核心叙事，不是收入规模，而是份额从10\%向上突破

报告中最有洞察力的部分，是对2027年机会的量化框架。公司目前看到的客户计算容量接近10GW，分析师用每GW约200亿美元的收入框架进行测算。虽然报告谨慎地指出这个数字更接近牛市情形，且每GW的内容价值因客户而异，但这一框架本身揭示了机会的量级。

更值得关注的是，Broadcom当前的AI市场份额仅为约10\%。报告认为这一比例将上升。在AI基础设施投资持续膨胀的背景下，10\%的份额意味着公司仍处于渗透率早期。如果份额提升至15\%或20\%，其收入弹性将远超当前市场一致预期。

这里隐含的投资逻辑是：市场目前对Broadcom的估值，更多反映的是其现有订单的可见性，而非份额提升的期权价值。而后者，恰恰是2027年最有可能超预期的来源。

关于CoT（竞争对手威胁）的讨论，报告给出了一个冷静的判断：联发科确实是一个风险，但并非颠覆性因素。分析师估计，真正处于风险中的TPU内容占比仅为10\%-15\%，Broadcom仍将保留大部分TPU份额。

这个判断的支撑逻辑在于：在如此大规模的量产节点上更换在位供应商，难度远超市场想象。这不仅涉及技术验证、产能爬坡，还包括与客户系统架构的深度绑定。报告认为Broadcom的TPU地位比许多投资者认为的更具韧性。

对于产业观察者而言，这提醒我们：在AI芯片领域，技术能力只是入场券，真正的竞争壁垒在于与客户共同演进的能力。Broadcom与谷歌TPU的深度合作关系，不是一朝一夕可以被替代的。

4. 报告尚未完全回答的关键问题：份额增长的驱动力和可持续性

尽管报告对2027年持乐观态度，但一个核心问题并未被充分展开：Broadcom的AI市场份额从10\%向上突破，具体驱动力是什么？是现有TPU客户的扩产，还是新客户的导入？如果是新客户，其规模和节奏如何？这些问题的答案，将直接影响份额提升的速度和确定性。

此外，报告提到的“每GW约200亿美元”框架，虽然提供了一个粗略的参考，但内容价值因客户而异。如果新客户的每GW内容价值显

[... middle omitted ...]

张节奏。报告提到公司看到接近10GW的客户容量，但这一数字的更新频率和客户构成，是判断需求真实性的先行指标。如果每个季度的容量指引持续上调，说明需求端比预期更强。

第二，非AI半导体业务的底部确认。报告指出，目前尚未看到广泛半导体组合的周期性复苏，但趋势已经稳定。Broadcom的非AI业务虽然“稳定但不令人兴奋”，但其触底回升的时点，将决定公司整体收入结构的健康度。

第三，VMware整合带来的现金流贡献。报告将基础设施软件视为稳定增长和现金流的来源。VMware的运营优化能否持续超预期，将直接影响公司整体估值中的软件部分。

这三个信号，比单个季度的AI收入数字更能反映公司的长期价值走向。

报告将目标价从470美元上调至485美元，核心驱动是2027年EPS估值倍数从27倍提升至28倍。这个倍数调整虽然幅度不大，但方向值得注意。在AI芯片领域，市场通常愿意为高增速支付高估值，但Broadcom的情况略有不同——它的估值溢价更多来自增长的可预见性，而非爆发性。

\subsection{[R028] 市场真正低估的不是需求，而是供给侧的再定价}
{\small\color{refgray}归类：区域市场：亚洲央行鹰派转向，澳洲投资逻辑切换}

澳大利亚的宏观叙事正在经历一次罕见的切换。过去几个月，投资者习惯性地将注意力集中在需求端——消费何时复苏、通胀是否见顶、RBA何时转向。但某外资投行在最新发布的年中展望中提出了一个值得认真对待的判断：澳大利亚的投资逻辑已经发生了根本性变化，而市场对这一变化的定价远未充分。

这份报告的核心洞察并不是某个单一数据的超预期或不及预期，而是一个结构性判断——澳大利亚正在进入一个“全新安排”的阶段。在这个阶段，决定资产价格走向的关键变量，从需求周期的波动，转向了供给侧的结构性再定价。具体来说，就是资本支出管道、劳动力市场调整、以及能源冲击的二次效应，正在取代传统的利率预期和消费数据，成为影响市场方向的主导力量。

这不是一个短期的交易主题。它意味着投资者需要重新审视自己分析框架中的权重分配。那些在过去几年中被视为次要变量的因素——比如国防开支的长期规划、数据中心建设的结构性需求、资源行业的资本纪律——现在正在成为决定市场走向的主线。

以下是我们从这份报告中提炼出的五个关键层次，它们共同指向同一个结论：市场的下一个拐点，不会来自需求端的意外，而是来自供给侧的重新定价。

报告开篇就点明了当前投资者争论的五个核心议题，其中最容易被低估的是资本支出管道的韧性。市场普遍担忧高利率和财政紧缩会压缩企业投资意愿，但报告指出，当前的资本支出驱动力与过去几轮周期有本质区别。

传统的资本支出周期高度依赖企业盈利预期和商业信心，而这一轮的结构性支撑来自三个几乎不受短期利率影响的来源：国防开支（联邦预算已追加530亿澳元）、能源转型与安全（150亿澳元能源安全包）、以及数据中心建设（电脑设备进口的异常飙升已经验证了这一趋势）。报告预测，六大关键领域的资本支出将从FY25的2250亿澳元增长至FY30的3260亿澳元。

这意味着什么？对于投资者而言，判断一家公司是否值得持有的标准正在发生变化。过去，市场更关注一家公司的需求弹性——它在消费疲软时能否守住营收。现在，更重要的指标变成了它的供给能力——它能否在结构性投资浪潮中获得订单、锁定产能、并最终将规模优势转化为定价权。那些能够嵌入国防、能源、数据中心供应链的公司，将在未来几年获得超越周期的增长。

2. RBA的困境不是要不要降息，而是能不能在通胀黏性中保持沉默

报告对RBA政策路径的判断值得细读。4月CPI从4.6\%降至4.2\%，低于预期，但核心CPI（trimmed mean）反而从3.3\%微升至3.4\%。临时性的燃料和公共交通补贴是主要驱动因素，而核心通胀仍然顽固地高于目标区间。

报告的核心判断是：RBA更可能选择“按兵不动”，而不是进一步加息。理由是国内需求放缓，尤其是房地产市场走弱——这是RBA认为最重要的货币政策传导渠道——给了央行更多“空间”来维持现状。但与此同时，RBA的沟通仍将保持鹰派基调，因为通胀黏性不允许它转向鸽派。报告明确表示，RBA转向更温和立场的前提不仅是增长放缓，还需要通胀趋势实质性走软，而这一点预计要到明年才会出现。

这对市场的含义是：利率预期的波动性将维持在高位，但方向性交易的空间反而收窄。投资者不应该押注RBA在短期内转向降息，也不应该过度担忧进一步加息。更务实的做法是接受一个“高利率更久”的环境，并据此调整组合的久期暴露和行业配置。那些对利率敏感但现金流稳定的资产，比如基础设施和部分REITs，可能比纯粹的成长股更具相对吸引力。

报告在盈利预期方面给出了一个清晰的警告。当前市场共识预期ASX 200在FY26和FY27的EPS增速分别为11.9\%和12.3\%，但报告自己的自上而下模型指向一个更为保守的展望——FY27增速可能只有约6\%。两者之间存在显著差距。

报告没有简单地说“共识太乐观”，而是提供了一个重要的分析框架：盈利下行的风险分布

[... middle omitted ...]

动力市场的分析指向一个容易被忽视的变量：不是失业率会上升到什么水平，而是劳动力市场走弱会如何改变企业和消费者的行为。报告预测失业率将在年底升至4.7\%，并在2027年继续上升。但这组数字本身可能并不是最值得关注的。

更重要的信号来自领先指标。报告引用了就业意向和失业预期的数据，两者都在转负。同时，报告提到了一个关键变量：AI对劳动力市场的影响。在收入环境走弱的情况下，企业可能加速采用AI来替代人工，这将在结构性层面改变劳动力市场的供需平衡。

劳动力市场的二阶效应有两个层面。第一，消费者行为将变得更加谨慎。过去几年，强劲的就业市场是支撑消费韧性的核心支柱。一旦这根支柱开始动摇，消费端的放缓可能比数据所显示的更为剧烈。第二，企业成本结构将发生变化。劳动力成本高企是过去两年企业利润率的主要压力来源，但如果就业市场走弱，工资增长放缓，企业的利润率反而可能获得喘息。这意味着，劳动力市场走弱对股市的影响并不是单向的——它对消费类股票是利空，但对部分工业和服务业股票可能是利好。

\subsection{[R029] 市场低估的不是零食需求，而是成本周期与渠道红利的共振}
{\small\color{refgray}归类：工业科技：全球自动化周期出现“双顶”，中国以外复苏加速}

二季度历来是中国零食行业的传统淡季。过去一个多月，覆盖的几家零食公司股价下跌了近9\%。市场似乎在用短期数据推演全年，但一份最新发布的投行研报给出了一个与股价走势相反的判断：真正的机会不在需求弹性的故事里，而在成本端结构性改善与新兴渠道持续放量共同作用下的利润弹性。

这份报告的核心信号很清晰：二季度业绩大概率是“混合体”——有亮点也有拖累，但下半年将迎来更强的恢复。研报维持对盐津铺子和卫龙美味的“买入”评级，同时小幅下调了卫龙盈利预测0-2\%。这种看似矛盾的调整，恰恰揭示了当前市场定价中尚未充分反映的结构性变量。

1. 量贩渠道的开店速度仍在超预期，这是行业增长最确定的引擎

市场对零食股的担忧之一，是量贩折扣渠道的红利是否已经见顶。但研报提供的数据给出了相反的答案：截至二季度至今，主要量贩系统的年初至今净开店数已达到约2500家/系统，而全年目标仅为 家。这意味着全年门店数量增速有望达到34-41\%。

更关键的是，盐津铺子和卫龙美味对量贩渠道的销售敞口分别高达约38\%和35\%。当渠道本身还在以超过30\%的速度扩张时，这两个品牌天然受益于“渠道铺货面积扩大”这一最朴素的增长逻辑。二季度虽然是淡季，但暑期消费旺季即将到来，量贩渠道的夏季高峰将直接转化为这两家公司的收入增量。

传统渠道和电商的表现则更为分化。盐津铺子在二季度至今看到传统散装产品的下滑幅度在收窄，而卫龙美味则因短期的组织架构调整出现疲软。电商方面，盐津铺子二季度电商收入同比下降约30\%，但研报预计下半年将转为同比增长。卫龙则预计2026年上半年电商增速将明显快于2025年下半年。

这里的核心含义是：渠道结构的此消彼长并不意味着增长消失，而是增长动能的切换正在发生。那些能够快速将资源从传统渠道转向量贩和电商的公司，将在这一轮渠道变迁中获得超额收益。

2. 魔芋产品是品类层面最确定的增长极，竞争格局比表面看起来更稳固

魔芋零食是近年来中国休闲食品行业少有的高增长细分品类。研报预计，盐津铺子和卫龙美味在二季度/上半年的魔芋产品收入增速将分别达到30\%以上和约20\%。

但真正值得关注的数据不是增速本身，而是市场份额的变化。截至2026年4月，卫龙美味在线下魔芋零食市场的份额为52.9\%，盐津铺子旗下“大魔王”品牌为11.0\%，而去年同期这两个数字分别是51.9\%和8.3\%。这意味着，即便市场上不断有新的竞争者进入，头部品牌的市场份额仍在扩大。

从月度数据看，卫龙的市场份额在2025年下半年一度回落至50\%以下，最低点出现在2025年11月的49\%。但随后快速回升，2026年2月甚至冲高到61\%。这种波动性说明竞争确实存在，但卫龙通过品牌力、渠道覆盖和产品迭代，有能力在竞争中重新夺回份额。盐津铺子的份额则从去年同期的8.3\%稳步提升至11.0\%，在10\%的水平上保持了近一年的稳定性。

这份数据合在一起意味着：魔芋品类的高增长并非昙花一现的品类红利，而是由头部品牌通过持续的市场份额扩张来驱动的。对于投资者而言，盯住市场份额的边际变化，比盯住品类增速更有意义。

如果说渠道和产品是收入端的逻辑，那么原材料成本就是利润端最重要的变量。研报明确指出，两家公司都有望受益于魔芋原材料价格的下行趋势。

魔芋原材料成本在2024年曾同比飙升约60\%，达到历史高点。但根据研报的数据，目前价格已从高点回落约25\%，且下行趋势有望延续至2026年下半年。研报给出的基准假设是2026年全年魔芋原材料成本同比下降约15\%。

对于零食企业而言，原材料成本是毛利率最敏感的变量。当成本下降时，企业通常面临两种选择：一是维持终端售价，直接释放利润弹性；二是适度降价以抢占市场份额。从目前的市场格局来看，头部品牌更可能选择前者——因为它们的市场份额仍在扩大，不需要通过价格战来

[... middle omitted ...]

一利润弹性的存在。

从估值角度看，盐津铺子当前交易在2026年预期市盈率约18倍，卫龙美味约11倍。后者还提供了约5\%的股息率（基于60\%的派息率）。

将这两家公司放在全球零食品牌的估值框架中来看：全球零食品牌平均2026年预期市盈率约17倍，中国零食品牌平均约15倍。卫龙11倍的估值不仅低于中国同行，也低于全球同行，而它的收入增速（ 年收入复合增长率约11\%）和净利润增速（约13\%）都显著高于全球均值。

盐津铺子的18倍估值看起来略高于中国同行均值，但考虑到其更高的增速（收入复合增长率约13\%，净利润增速约20\%）和更强的渠道布局，这一估值并非没有支撑。

更重要的是，这两家公司的估值已经反映了二季度淡季的疲软，以及市场对竞争加剧的担忧。但研报给出的下半年净利润增速指引——盐津铺子同比增长27\%，卫龙美味同比增长12\%——尚未被充分定价。如果这一增速在下半年兑现，当前的估值水平将显得过于保守。

5. 报告尚未完全回答的关键问题：竞争烈度何时会实质性侵蚀利润率

\subsection{[R030] 口服GLP-1的竞赛逻辑已经变了：中国资产正在定义下一轮竞争参数}
{\small\color{refgray}归类：医疗健康：双抗改写肺癌标准治疗，口服GLP-1竞赛参数被重新定义}

口服GLP-1的竞赛逻辑已经变了：中国资产正在定义下一轮竞争参数

距离美国糖尿病协会年会还有不到两周，市场关注的焦点早已不是“谁会发布数据”，而是“谁的数据能改写竞争格局”。

某外资投行最新发布的医疗健康研报，围绕即将到来的ADA会议，系统梳理了中国GLP-1管线的最新进展。这份报告的核心判断值得每个关注代谢疾病赛道的投资者认真对待：口服GLP-1的市场竞赛正在经历一次结构性转变，而中国生物科技公司正在从“跟随者”变成“参数定义者”。

这不是一个关于“中国公司有没有机会”的问题——机会窗口已经打开。真正的问题是：哪些公司能在这个窗口期内完成从临床数据到全球商业化的跨越。

1. 口服GLP-1的2030年市场规模已经清晰，但竞争参数正在被重新定义

研报给出了一个值得反复推敲的框架性判断：到2030年，口服减肥药市场将占据整个肥胖药物总可寻址市场的约35\%，规模超过350亿美元。这个数字本身并不令人意外——真正重要的是支撑这个预测的逻辑链。

驱动口服GLP-1市场爆发的两个关键变量已经落地：Wegovy口服片和orforglipron（Foundayo）的获批。这两个药物的上市，不仅验证了口服GLP-1的临床可行性，更重要的是，它们为后来者设定了必须超越的基准线。

研报特别强调了一个容易被忽视的细节：口服GLP-1的竞争焦点并不是减重幅度，而是耐受性。Wegovy口服片和orforglipron在减重上的表现——13.9\%和11.5\%的安慰剂调整后减重——已经将门槛推到了足够高的位置。对于初级诊疗场景，超过10\%的减重效果已经可以被视为“可接受”。

真正决定胜负的，是呕吐、恶心等胃肠道不良反应的发生率。研报明确指出，任何新进入的口服GLP-1候选药物，都必须将呕吐率控制在20\%以下，才能与现有药物竞争。这一点，正在成为整个赛道的“隐形准入门槛”。

这意味着什么？意味着后来者不能仅仅依靠“更强的减重效果”来获得差异化——耐受性曲线上的每一个百分点改善，都可能比减重绝对值上的两三个百分点更有商业价值。

2. 信达生物的IBI3032：早期数据惊艳，但真正的考验在临床执行

在即将召开的ADA会议上，信达生物的口服小分子GLP-1药物IBI3032将是市场最关注的数据之一。公司此前已经透露，IBI3032在4周治疗中实现了10\%的减重效果。

这个数字放在任何语境下都是引人注目的。如果将其与现有药物的减重曲线对比，IBI3032的早期数据已经显示出潜在的差异化优势——它可能是继Structure Therapeutics的aleniglipron之后，第二个在减重幅度上明显超越现有基准的口服GLP-1候选药物。

第一，样本量。信达生物表示IBI3032的I期临床共入组了超过100名患者，但市场更关心的是ADA上会公布多少数据。如果公布20-30例的多次给药数据，这个样本量已经足够减少随机变异，提供有意义的参考。

第二，剂量依赖关系。10\%的减重是在哪个剂量水平达到的？这个问题的答案将直接决定IBI3032的治疗窗口和商业化潜力。如果高剂量才能达到这个效果，那么耐受性问题就会凸显；如果中低剂量就能实现，那将是强有力的差异化信号。

第三，单次给药与多次给药的耐受性差异。早期临床中，不良反应率的波动通常较大，但市场会关注多次给药后耐受性是否改善——这将是判断未来剂量滴定方案设计的重要线索。

研报给出了一个更为关键的判断框架：IBI3032的管线价值，最终取决于两个变量——与后期阶段同类药物的差异化程度，以及全球临床的执行能力。按时间推算，IBI3032比aleniglipron落后约2.5年，比Ascletis的ASC30落后约1.5年。在口服GLP-1这个赛道上，时间本身就是一种成本。信达生物能否

[... middle omitted ...]

的临床数据——这标志着恒瑞正在将GLP-1管线从单纯的减重和糖尿病适应症，向更广泛的代谢疾病领域拓展。

另一个值得关注的信号是HR17031——GLP-1与胰岛素的固定剂量复方制剂。在血糖控制不佳的糖尿病患者中与甘精胰岛素的头对头比较数据，将揭示恒瑞在联合治疗领域的布局深度。

研报还提及了恒瑞的下一代多靶点药物HRS-4729（GLP-1/GIP/GCG三靶点），在12周治疗中实现了16\%的减重效果（安慰剂调整后10.6\%）。虽然非头对头比较显示其疗效略低于United Labs/Novo Nordisk的UBT251，但恒瑞正在构建的是一个从单靶点到三靶点、从口服到注射、从单一适应症到多适应症的“全谱系”组合。

这种布局策略的商业逻辑是清晰的：在GLP-1这个赛道上，单一产品的领先优势可能难以持久，真正的护城河来自于管线深度和组合协同。恒瑞的策略，本质上是在赌一个判断——GLP-1市场的终局竞争，不会是“一款药赢者通吃”，而是“一个产品组合覆盖所有细分场景”。

\subsection{[R031] 韩国出口的“分裂式繁荣”：半导体孤军深入，而结构性风险正在积累}
{\small\color{refgray}归类：全球贸易与供应链：补库信号已现，有效关税税率差异重塑格局}

韩国出口的“分裂式繁荣”：半导体孤军深入，而结构性风险正在积累

韩国5月出口数据再度超出市场预期，同比增速攀升至53.2\%，连续第六个月实现环比正增长。某外资投行在最新研报中确认了这一趋势，并指出其高于自身已属乐观的预测。对于关注全球贸易与供应链的决策者而言，这份数据传递的信号远比表面数字复杂：技术产品出口连续九个月上升，半导体更是保持强劲；但非技术类出口的持续萎缩、以及对单一品类依赖度的加深，正在为这轮出口繁荣埋下隐忧。真正值得追问的问题不是“韩国出口还能涨多久”，而是“当半导体周期拐点来临时，韩国经济的缓冲垫在哪里”。

5月韩国半导体出口环比增长3.9\%，延续了数月来的强势。但与此同时，计算机、手机、家电等其他科技产品出口环比下降4.7\%，这是自2025年6月以来首次出现下滑。非科技类出口更是集体走弱，船舶、汽车及零部件、机械、能源相关产品无一幸免，环比跌幅从前月的1.9\%扩大至2.8\%。

这意味着韩国出口的“增长面”正在急剧收窄。半导体已经从一个重要引擎变成了几乎唯一的引擎。研报数据显示，半导体出口的持续增长足以抵消其他所有品类的疲弱，但这也意味着一旦半导体需求出现波动，整体出口将面临断崖式调整的风险。对于投资者而言，关注韩国出口不应再笼统地看总量，而应拆解为“半导体”和“其他”两个独立变量来跟踪。

2. 出口目的地分化：美国与新兴市场成为双引擎，欧洲意外失速

从目的地来看，5月韩国出口的增长是广泛但非均匀的。对东盟出口环比飙升17.4\%，几乎贡献了全部增量；对中国和日本出口也分别实现5.3\%和6.8\%的环比增长。对美出口连续第八个月上升，环比增长3.0\%，势头稳健。

唯一显著掉队的是欧盟，韩国对欧盟出口环比下降4.2\%。这一反差值得警惕：在全球贸易再平衡的背景下，美国和中国市场对韩国产品的需求仍在扩张，而欧洲却出现了萎缩。报告没有深入探讨这一分化的原因，但结合当前欧洲制造业PMI持续低迷、能源成本高企的背景，韩国对欧出口的弱势可能不是一个短期波动，而是结构性竞争力的反映。对于关注全球贸易格局的读者，这是一个需要持续追踪的信号。

3. 贸易顺差创历史新高，但“数量”与“价格”的贡献尚未厘清

5月韩国贸易顺差达到269亿美元，超过前月的238亿美元，再度刷新历史纪录。表面看，这是一个极为亮眼的数字。但仔细拆解会发现，顺差的扩大并非完全来自出口量的强劲增长。

进口方面，能源相关进口环比增长11.3\%，延续了前月10.6\%的强势。这意味着韩国正在以更高的成本进口能源，而出口端的高增长中，有多少来自价格上涨、有多少来自实际出货量的增加，报告并未给出清晰拆解。对于贸易顺差的质量判断，价格效应和数量效应的区分至关重要。如果出口增长主要由价格驱动而非量驱动，那么顺差的可持续性将大打折扣。这一点，报告留下了空白。

4. 一个尚未被充分讨论的关键问题：半导体的“超级周期”是否已经透支了需求？

报告确认半导体出口持续强劲，但并未深入分析需求端的结构。当前全球AI基础设施投资热潮推动了对高端存储芯片的旺盛需求，韩国半导体企业显然是这一周期的最大受益者之一。但这一需求是否可持续？企业客户是否在提前囤货？下游终端消费电子需求是否真的同步复苏？

这些问题报告没有回答，但它们是判断出口前景的核心。如果半导体需求中有相当一部分是“预防性备货”或“政策驱动”的短期行为，那么后续增速放缓的幅度可能超出预期。而一旦半导体增速回落，韩国整体出口将失去唯一的支撑点，非科技类出口又无法接棒，届时贸易顺差的收缩将比当前扩张更剧烈。

5. 对读者的观察框架：从“总量”到“结构”再到“拐点信号”

基于这份报告的解析，我们建议读者建立以下观察框架，以替代简单的“出口增速”跟踪：

第一，关注半导体出口环比增速的变化趋势，而非同比

[... middle omitted ...]

可能是全球贸易格局变化的一个先行指标。如果欧洲需求持续疲弱，韩国企业可能被迫将更多产能转向美国和中国市场，加剧竞争。

第四，留意贸易顺差中“价格效应”与“数量效应”的拆分。任何关于韩国出口前景的深度分析，都必须回答“多出来的顺差到底来自哪里”。

这份某外资投行的研报提供了扎实的数据基础，但正如我们上文指出的，许多关键判断仍有待验证。半导体周期的可持续性、非科技出口能否触底、贸易顺差的真实质量，这些问题的答案不在当前的公开数据中，而需要更深入的行业调研和交叉验证。

如果你对这些未解问题感兴趣，欢迎来星球微信群里继续讨论。我们会持续跟踪韩国出口的结构性变化，并在关键拐点出现时提供更及时的分析。完整研报的原始图表和更多解读，也将在社群内分享。

Personal reading notes and learning share only. Not investment advice.

韩国5月出口同比+53.2\%，连续6个月增长，比预期更强。但拆开看，全靠芯片在扛。

\subsection{[R032] 韩国半导体出口的真相：市场低估了供给约束的持久性}
{\small\color{refgray}归类：AI与半导体：需求扩散至多节点，供给约束成为新定价变量}

2026年5月，韩国半导体出口录得255\%的同比增长，DRAM出口更是创下370\%的同比增速，这是自2008年有追踪数据以来的最高记录。但真正值得关注的，不是这个数字本身，而是它背后的结构含义：这一轮高增长并非需求端的全面爆发，而是在供给端持续收紧、资本开支谨慎的背景下，由AI驱动的结构性需求与有限产能之间的错配所推动。市场如果把这份数据简单理解为“周期上行”，就可能低估了供给约束的持久性，以及它对定价权的重塑作用。

某外资投行最新发布的韩国科技出口追踪报告，通过海关数据、进口中间品追踪、设备进口趋势等角度，勾勒出一个比表面数字更复杂的图景。以下是我们从这份报告中提炼出的五个关键洞察，以及一个尚未被充分讨论的问题。

1. 连续四个月200\%以上的同比增长，背后是两重力量的叠加而非单一驱动力

5月韩国内存出口同比增长255\%，这是连续第四个月增速超过200\%。DRAM出口增长370\%，NAND芯片出口增长207\%，SSD出口增长338\%——几乎所有内存子品类都在高速运转。

但“为什么”比“多少”更重要。报告中的数据显示，这一轮高增长是两重力量叠加的结果：一是AI服务器对HBM和高端DRAM的持续消耗，二是传统内存产品在低基数上的补涨。三星电子的内存收入预计在2Q26同比增长452\%，增速高于SK海力士的271\%，原因正是三星对传统内存产品的暴露更高，且在HBM收入上有一个较低的基数。

这意味着，市场不能简单地用“AI需求强劲”来解释一切。传统内存的爆发式增长中，有一部分是补库存和价格反弹的贡献，而这部分增长的可持续性，取决于下游终端需求的真实恢复情况。报告提到，韩国产业通商资源部将整体强劲归因于美国和中国大型科技公司的强劲需求——但这里的“大型科技公司”究竟是在囤货还是真实消耗，报告并没有完全展开。

2. 日本进口的封装材料数据，暴露了HBM产能爬坡的节奏差异

这份报告最有价值的洞察之一，是它通过追踪日本向韩国特定地区出口的封装材料，来推断SK海力士和三星电子的HBM出货量。这是一个典型的“高频替代数据”思路：SK海力士的HBM主要使用MR-MUF工艺，其关键材料环氧树脂主要从日本Namics公司采购，进口至利川/清州；三星电子的HBM主要使用TC-NCF工艺，其关键材料塑料薄膜主要从日本Resonac采购，进口至华城/平泽。

报告显示，4月SK海力士的环氧树脂进口同比增长57\%，三星电子的塑料薄膜进口同比增长96\%。从历史数据看，这两项进口数据与两家公司的HBM出货量有很强的相关性。但更值得关注的是节奏差异：三星电子的进口增速远高于SK海力士，这与其在HBM市场的追赶策略一致。

然而，这里有一个尚未解答的问题：进口数据反映的是备货行为还是实际生产？报告本身也承认，这些数据是“良好代理变量”，但并非精确对应。从Exhibit 1和Exhibit 2的数据看，1Q26 SK海力士的环氧树脂进口环比下降，而同期其HBM出货量仍在增长——这说明库存消耗或工艺改进可能正在发生。对于试图通过高频数据提前判断HBM出货节奏的投资者来说，这个差异值得深入拆解。

5月韩国半导体设备进口同比增长99\%，出口同比增长66\%。报告认为，设备进口的强劲增长来自内存供应商为缓解供应紧张而增加的资本开支，而出口的强劲则来自全球半导体产能建设的旺盛需求。

这个数据点的重要性在于：它暗示了韩国内存厂商正在加速扩产。但扩产的方向是什么？是HBM专用产能，还是传统DRAM产能？报告没有给出明确的拆分。从逻辑上推断，如果扩产主要集中在HBM，那么传统DRAM的供给约束可能持续更久，价格支撑更强；如果扩产覆盖传统DRAM，那么当前的高增速可能提前见顶。

这个问题的答案，将直接影响对2027年内存定价的判断。报告中的设备进口数据提

[... middle omitted ...]

这些领域的判断是：OLED的增长来自新手机产品的拉动，但终端市场需求因半导体成本上升而受到抑制；MLCC的增长主要由AI服务器和汽车用MLCC驱动；锂离子电池的增长则来自EV和ESS电池出口的增加以及关键矿产价格上涨带来的ASP提升。

这些领域的共同特征是：增长有支撑，但不够强劲，且高度依赖特定细分市场。对于非AI领域的科技公司来说，当前的环境更像是“结构性分化”而非“全面复苏”。那些能够切入AI供应链的公司（如三星电机在AI服务器MLCC上的布局）表现更好，而那些依赖消费电子和传统汽车的公司则面临更大的不确定性。

5. 报告尚未完全回答的关键问题：高增长的可持续性取决于两个变量的走向

这份报告提供了丰富的数据和清晰的逻辑框架，但有两个关键问题没有给出明确答案：

第一，内存出口的高增速还能持续多久？当前的增长中，价格因素和量的因素各占多少？报告给出了出口金额，但没有拆分ASP和出货量。如果当前的高增长主要来自ASP上涨，那么一旦价格拐点出现，增速可能迅速回落。

\subsection{[R033] 锂价反弹的真正考验：供给侧正在醒来，但市场低估了重启的节奏}
{\small\color{refgray}归类：能源与大宗：供给恢复慢于预期，锂矿重启信号密集}

过去两个月，锂辉石精矿价格从每吨不足800美元回升至1500美元以上。对于经历了2024年剧烈去库存的锂行业而言，这无疑是久违的暖意。但真正的问题不是价格是否还能涨，而是当价格涨到某个临界点，那些在低谷期被迫关停的产能，会以多快的速度重新回到市场。

某外资投行最新发布的澳大利亚材料行业研报，捕捉到了一个正在发生的结构性变化：澳大利亚的锂矿重启信号正在密集出现。Bald Hill、Mt Marion、Mt Holland、Pioneer Dome——这些名字在短短两周内重新进入市场视野。它们不是遥远的勘探故事，而是已经或即将做出最终投资决策的现实产能。

这份报告的核心判断是：锂价上涨正在激励供给端的响应，但市场对响应的速度和规模可能准备不足。2026年全球锂市场预计仍有约11.7万吨碳酸锂当量的缺口，但这并不意味着价格可以线性外推。重启产能的累积效应，可能在 年将市场推向小幅过剩。而真正决定价格走势的，不是需求的故事，是供给的节奏。

1. 澳大利亚矿企正在用“低成本重启”回应价格上涨，而不是大规模新建产能

这轮供给响应的最显著特征，是“重启”而非“新建”。矿商们没有急于批准绿地项目的数十亿美元资本开支，而是优先恢复那些已经完成采矿、选矿基础设施的停产项目。

报告列举了过去两周内的关键动态：Mineral Resources的Bald Hill重启，将增加14万吨/年的SC6精矿产能，约合1.8万吨LCE，首批货物最早今年7月即可发出。Mt Marion的浮选厂和地下扩建项目，将增加10万吨/年产能，约合1.3万吨LCE，目标在2028财年投产。Mt Holland获得监管批准，将矿石加工能力翻倍至440万吨/年，潜在新增约4.5万吨LCE。Pioneer Dome的DSO项目同样在推进，可能增加约2.1万吨LCE。

这些项目的共同点是什么？资本支出低、投产周期短、且大部分已有基础设施。这不是 年那种“先讲故事、再融资、再建设”的模式，而是“先有矿、再重启、快速出矿”的现实主义路径。这意味着供给响应的速度可能比市场预期的更快。

对于投资者而言，这提出了一个关键问题：当市场还在用2026年的供需缺口来论证价格上行空间时，2027年之后的供给曲线已经在被重新绘制。如果重启产能的累积效应超出预期，当前的价格可能已经包含了部分未来的供给溢价。

2. 2026年的缺口并不稀缺，稀缺的是对2027年之后供给曲线的正确判断

报告给出的市场平衡表显示，2026年全球锂市场预计存在约11.7万吨LCE的缺口，但2027年将转为5万吨过剩，2028年进一步扩大至9.6万吨。到2030年，市场将重新回到小幅短缺状态。

这一组数字本身并不令人意外。真正值得深究的是背后的假设。报告明确提示，其2026年供需平衡模型假设中国批发BEV增长率为9\%，而年初至今的实际运行率仅为约3\%。如果需求端不及预期，缺口的规模可能被高估。反之，如果供给端重启的速度快于预期，2027年之后的过剩规模也可能被低估。

报告还特别提到了钠离子电池的潜在影响。如果钠离子电池订单持续增长，可能会侵蚀磷酸铁锂电池在储能、入门级乘用车和轻型商用车市场的份额。这是一个尚未被充分定价的长期变量。

对于产业决策者而言，这意味着：当前锂价的上涨，本质上是对2024年过度去库存的修正，而不是新一轮超级周期的开始。真正的供需平衡点，取决于两个变量的博弈——需求增长的真实斜率，以及低成本重启产能的释放节奏。

3. 报告尚未完全回答的关键问题：重启产能的成本曲线究竟在哪里？

这是整份报告中最值得追问的问题。报告列举了多个重启项目，但并未给出每个项目的完全成本曲线。Bald Hill和Ngungaju这样的项目，其重启成本可能低于行业平均，因为它们不需要重新建设矿

[... middle omitted ...]

升至1500美元/吨以上，中国锂盐厂是否会重新开始补库？如果补库需求集中释放，可能会在短期内推高价格，从而进一步激励澳大利亚矿商加速重启。这是一个自我强化的循环，其持续时间取决于下游的真实需求和库存水平。

对于持仓者而言，这些不确定性意味着：当前的价格区间可能不是稳态。基于2026年缺口的看多逻辑是成立的，但需要警惕2027年之后的供给过剩风险。真正重要的是判断重启产能的边际成本曲线，以及中国需求的实际增速是否能够消化新增供给。

基于报告的供给侧分析，可以建立一个简易的观察框架，帮助判断锂价是否正在接近周期性顶部。

第一个信号是澳大利亚重启项目的FID密集度。如果未来3个月内，有超过3个类似Bald Hill或Pioneer Dome规模的项目宣布FID，说明供给响应的速度正在加速。报告中已经列出了Finiss、Ngungaju、Kathleen Valley、Pioneer Dome、Mt Marion等多个潜在项目，关键看它们是否会在同一时间窗口内做出决策。

\subsection{[R034] 欧洲钢铁业正在经历一场被低估的供给侧重定价}
{\small\color{refgray}归类：能源与大宗：供给恢复慢于预期，锂矿重启信号密集}

欧洲钢铁市场正在发生一个被多数投资者低估的结构性转折。这不是一个周期性的需求回暖故事，而是一套由政策驱动的供给侧框架正在从根本上改写进口平价、利润池分配和资产定价逻辑。

某外资投行最新发布的欧洲钢铁行业研报指出，欧盟新钢贸政策将进口配额较2024年水平削减约47\%，超出配额的关税税率提升至50\%，并计划从2026年7月起实施。与此同时，碳边境调节机制也将在2026年进入正式执行阶段，对进口钢材的碳合规成本施加更大压力。这些政策叠加在一起，意味着欧洲钢铁利润池正在向本土钢厂迁移，而市场对这一结构性变化的定价远未完成。

报告的核心判断是：欧洲钢铁股的风险回报比已经显著改善，但投资者仍在使用旧框架来评估新格局。以下是我们从这份报告中提炼出的五个关键洞察，以及一个尚未被充分回答的问题。

市场通常将欧盟钢贸政策理解为“加关税”，但报告揭示的机制远比关税复杂。进口配额削减47\%叠加50\%的额外关税，本质上是将进口钢材的成本底线系统性抬高。这意味着欧洲本土钢厂不再需要与全球边际产能竞争定价，而是可以基于新的进口平价来设定售价。

更重要的是，碳边境调节机制的执行进一步挤压了低成本地区对欧洲出口的经济性。碳合规成本正在成为进口钢材的一项固定附加成本，且这一成本会随碳价上升而递增。对于依赖高炉工艺的海外钢厂而言，这一成本结构的变化是永久性的，而非周期性的。

报告明确指出，这些措施共同支持了欧洲本土钢价，并带来了对市场一致预期的盈利上调空间。但这里需要强调的是：市场对“价格支持”的认知仍然停留在短期政策刺激层面，而报告的逻辑指向的是一个更长期的利润池再分配——从海外进口商转向欧洲本土产能。

2. 利润池迁移的核心受益者是拥有本地资产和运营灵活性的综合钢厂

政策变化的赢家并非所有欧洲钢厂，而是那些能够在新的定价环境下实现产能利用率和盈利杠杆双提升的企业。报告特别指出，拥有本地资产、装运灵活性和高运营杠杆的综合碳钢企业将是这一轮利润池迁移的最大受益者。

安赛乐米塔尔被列为最受偏好的碳钢标的。报告认为，该公司对欧洲盈利复苏具有强劲的运营杠杆，同时运营和资产负债表风险有限。其上行情景下的企业价值对息税折旧摊销前利润倍数压缩至约3.9倍，而历史周期均值约为7.1倍，这意味着45\%-50\%的折价。报告还指出，盈利上行将支持公司在基础股息之外加速股票回购，进一步强化股权回报。

瑞典钢铁公司也被列为超配评级。报告认为市场尚未充分反映美国钢板业务动能、特种钢盈利韧性以及能源成本敞口下降这三者的组合效应。尽管奥克瑟勒松德和吕勒奥的产能转型仍存在执行风险，但近期的盈利动能正在改善。

这里的关键判断是：市场对欧洲钢企的估值折价，部分源于对需求端的悲观预期，但报告暗示，如果供给侧政策能够持续支撑价格和利用率，即使需求不出现显著复苏，钢企的盈利中枢也可能上移至高于市场预期的水平。

不锈钢与碳钢面临的政策驱动力并不完全相同，但报告对不锈钢板块同样持建设性展望。欧盟不锈钢价差今年以来已经扩大，主要得益于政策支持、进口显著回落——碳边境调节机制和预期中的保障措施是主要推手。美国不锈钢价差也在需求部分恢复的背景下走阔。

奥珀姆被列为最受偏好的不锈钢标的。报告指出，该公司业务模式多元化，涵盖巴西、回收和增值产品业务，这为其盈利韧性提供了支撑。更重要的是，其欧洲不锈钢业务当前产能利用率仅为65\%，这意味着一旦欧洲需求出现任何边际改善，其盈利弹性将相当显著。

阿塞里诺克斯同样被超配。报告认为该公司拥有美国业务和合金业务带来的盈利韧性，通过收购海恩斯合金公司实现了美国扩张和合金业务增长，尽管去杠杆限制了未来两年的股东回报，但估值相对于自身历史水平和合金同行仍然具有吸引力。

这里值得注意的一个微妙判断是：不锈钢的盈利驱动因素中，政策对进口的抑制效果比

[... middle omitted ...]

息税折旧摊销前利润倍数仅为4.7倍，而周期均值约为7.0倍。报告认为，随着现金消耗下降和执行风险缓和，估值修复空间仍然存在。

对于蒂森克虏伯，报告指出其股价已从近期高点回落，投资者似乎正在折价计入更严峻的宏观环境、能源成本担忧以及钢欧洲业务战略选项的执行风险。但报告也暗示，如果钢欧洲业务能够实现价值增值的交易或重组，电梯业务持股能够兑现价值，以及材料服务业务按计划分拆，那么投资案例仍存在多个关键催化节点。

整体而言，报告传递的信号是：市场对欧洲钢企的盈利预期可能仍然保守，尤其是在 年政策效果逐步显现之后，盈利上修的概率正在上升，但市场尚未对此充分定价。

5. 报告尚未完全回答的关键问题：需求侧能否配合供给侧的重构

报告最值得注意的克制之处在于，它并未对欧洲钢铁需求给出乐观判断。从图表数据来看，德国建筑许可、欧洲建筑业采购经理人指数和建筑业信心指标均处于低位，尚未出现明确的复苏信号。全球钢铁产量增速也在放缓， 年的复合年增长率仅为1.1\%，远低于 年的5.4\%。

\subsection{[R035] 市场真正低估的不是需求，而是供给恢复的速度}
{\small\color{refgray}归类：能源与大宗：供给恢复慢于预期，锂矿重启信号密集}

当前布伦特原油价格约92美元/桶，经通胀调整后恰好处于过去二十年的中位数水平。这个数字本身并不引人注目，真正引人注目的是它所隐含的预期：市场正在对一个“干净且近在咫尺”的供应恢复进行定价。但某外资投行最新发布的研报提出了一个截然不同的判断——供应恢复的速度将远比市场预期的更慢、更曲折，而当前价格已经计入了过于乐观的假设。

这份由该行首席商品策略师领衔的报告，调整了其对霍尔木兹海峡重新开放及后续产量恢复路径的建模假设。核心结论是：即便海峡在7月下旬重新开放，产量恢复到正常水平也需要至少4个月，且恢复过程将呈现明显的“前慢后快”特征。在此基础上，该行维持三季度布伦特100美元/桶的预测，并将四季度预测从90美元上调至95美元，2027年一季度从80美元上调至85美元。

在92美元的现货价格面前，维持近端预测本身就是看多的信号。但这篇报告真正值得关注的地方，不在于价格预测本身，而在于它揭示了一个被市场系统性低估的结构性问题：大规模断供后的供给恢复，其复杂性和时间成本远超大多数投资者的直觉。

1. 市场的定价逻辑建立在“快速恢复”这一尚未验证的假设之上

当前油价最令人困惑的特征不是价格水平，而是价格曲线的形态。自5月中旬以来，布伦特近月期货的价差结构已经大幅走平。近月布伦特CFD——最能反映现货市场紧张程度的指标——从5.7美元/桶降至2.6美元/桶，M1-M2价差从3.8美元/桶降至不足1美元。这意味着，市场不仅没有对供应风险进行重新定价，反而在过去三周里变得更加乐观。

这种乐观情绪的逻辑链条是清晰的：市场参与者相信，一旦政治协议达成，供应将迅速恢复。但这份报告指出，这条逻辑链上的每一个环节都存在显著的不确定性。

报告将其对海峡重新开放的建模假设从5月底推迟到7月下旬。这并非简单的日期推迟，而是反映了两个层面的判断：第一，截至报告撰写时，美伊之间尚未达成任何协议或谅解备忘录，谈判可能还需要数周时间；第二，即便政治协议达成，排雷、商业信心重建、保险条款重新协商等操作层面的问题，还需要额外数周时间。

这里的关键洞察是：政治协议与供应恢复之间，存在一个被市场严重低估的“执行鸿沟”。市场将“达成协议”等同于“供应恢复”，但实际的时间线要复杂得多。对于资产定价而言，这意味着当前油价所隐含的“V型恢复”路径，可能面临系统性向上修正的风险。

报告详细列举了供应恢复过程中需要克服的六重障碍，这些障碍的叠加效应是理解供给恢复速度的关键。

第一重是安全与保险问题。即便海峡被宣布重新开放，船东和保险公司需要看到切实的证据，证明通行是安全的。排雷本身就需要数周到数月的时间。

第二重是油轮错配。过去数月间，无法进入海峡的油轮已经改道至美国墨西哥湾、西非、巴西等地。沙特阿美CEO在一季度业绩会上将其称为“最大的问题”，并指出即使在最乐观的情景下，供应链也需要数月时间才能恢复正常。虽然海峡外有少量空置油轮可以立即投入使用，但大规模的油轮重新部署需要时间。

第三重是存储空间。油井无法在储油设施已满的情况下恢复生产。这对于伊拉克和科威特等缺乏管道出口能力的国家尤为重要，它们的海运出口实际上已降至零。

第四重是油田重启的技术难度。根据Rystad Energy的数据，海湾六国约有1万口油井处于停产状态。停产时间越长，重启难度越大。报告引用了Rystad的估算：停产两个月后面临重启约束的油井约2800口，四个月后增至4100口，六个月后超过5000口。报告假设7月下旬恢复生产意味着停产约4.5至5个月，对应约4000至5000口油井存在重启困难。

第五重是基础设施损坏。物理损坏主要集中在炼油和出口终端，而非原油生产环节。受损炼油能力约430万桶/日，其中约250万桶/日仍处于停运状态。修复这些设施可能需要2至3个季度。

第六重

[... middle omitted ...]

国原油出口量在冲突期间保持高位，部分抵消了中东出口的下降。中国的低进口则减少了全球市场的需求压力——中国在冲突前就降低了原油采购量，这使得原本可能流向中国的油轮得以转向其他市场。

但问题在于，这三个缓冲机制都有其局限性。美国出口不可能无限增长，中国的进口需求也不会永远低迷，而“预期重新开放”这一缓冲本身就是不稳定的——它建立在尚未被验证的假设之上。

报告特别指出，三季度市场将保持与二季度同样紧张的状态，库存将进一步下降。这意味着，当前的缓冲机制可能难以在整个三季度维持。

这份报告的分析框架主要集中在供给端，对于需求侧的分析相对有限。这是报告本身明确承认的局限，也是读者需要自行判断的关键变量。

一个尚未被充分回答的问题是：高油价对需求的影响是否被低估？当前92美元/桶的布伦特价格，经通胀调整后处于20年中位数水平，但考虑到全球经济增长的放缓，这个价格水平对需求的实际抑制效应可能被低估。特别是在欧洲和部分新兴市场，能源成本已经对工业生产和消费者支出产生了显著影响。

\subsection{[R036] ASCO 2026最值得关注的信号：双抗在肺癌领域首次击败PD-1标准疗法}
{\small\color{refgray}归类：医疗健康：双抗改写肺癌标准治疗，口服GLP-1竞赛参数被重新定义}

ASCO 2026最值得关注的信号：双抗在肺癌领域首次击败PD-1标准疗法

每年ASCO（美国临床肿瘤学会）会议结束后，市场往往会被大量数据淹没。但今年会议周末传出的一个信号，值得产业决策者和投资者停下来重新审视：在一项头对头的III期研究中，PD-1xVEGF双抗联合化疗首次在统计学上显著优于PD-1单抗联合化疗——这是该领域过去十年未曾被打破的格局。

这份某外资投行在ASCO会议期间发布的研报，围绕这一核心信号展开，并提供了多个细分领域的交叉验证。我们提炼出五个关键洞察，以帮助读者理解这些数据背后的产业含义，以及尚未被市场完全定价的变量。

1. 肺癌一线治疗格局可能正在被重新定义，双抗的“头对头胜出”是历史性突破

在ASCO全体会议上，Summit Therapeutics（SMMT）的合作伙伴康方生物公布了HARMONi-6研究的总体生存期（OS）数据。这项在中国开展的III期试验，将依沃西单抗（ivonescimab，一种PD-1xVEGF双抗）联合化疗，与替雷利珠单抗（抗PD-1）联合化疗进行头对头比较，用于一线鳞状非小细胞肺癌（NSCLC）患者。

结果：依沃西单抗组中位OS为27.9个月，对照组为23.7个月，风险比（HR）为0.66（p=0.0017），死亡风险降低34\%。这一数据不仅超过了市场此前预期的0.72的HR门槛，更重要的是，这是首个在III期试验中证明比PD-1联合化疗更优的方案。

此前市场对PD-1xVEGF双抗的质疑主要集中在：PD-1与VEGF的双重靶向是否会带来额外的毒性，以及能否在OS上真正超越已成熟的PD-1疗法。HARMONi-6的数据给出了明确回答。值得注意的是，这一结果也为SMMT正在进行的全球III期HARMONi-3研究提供了正向信号——该研究将依沃西单抗与Keytruda（帕博利珠单抗）进行头对头比较，预计将于今年下半年公布最终PFS和中期OS数据。

*产业含义**：如果HARMONi-3能够复现HARMONi-6的结果，那么全球一线NSCLC的标准治疗将面临实质性挑战。对于拥有PD-1产品的药企而言，这不仅是市场份额的威胁，更意味着其核心产品的生命周期可能被压缩。而对于正在开发双抗或联合疗法的公司，这一结果提供了关键的临床验证。

除了肺癌，SMMT在ASCO上还展示了一项针对转移性结直肠癌（mCRC）的II期研究海报数据。尽管样本量有限（49例患者），但结果值得关注：在两个剂量组（20mg/kg和10mg/kg）中，客观缓解率（ORR）均达到70.8\%，疾病控制率（DCR）为100\%，9个月无进展生存率分别为76.1\%和70.1\%。

更重要的是，高剂量组（20mg/kg）的缓解持续时间（DoR）显著优于低剂量组（9个月DoR估计值为79.1\% vs. 41.5\%）。尽管高剂量组的不良事件发生率略高，但研究者指出，由于当前mCRC一线标准治疗已包含贝伐珠单抗（抗VEGF），医生对管理此类毒性已有经验。

*竞争格局视角**：该数据与辉瑞（PFE）的同类PD-1xVEGF分子PF-4404相比，在ORR和DCR上均展现出优势（PF-4404的II期ORR为57.1\%，DCR为95.2\%）。虽然跨试验比较需要谨慎，但这一信号表明，依沃西单抗在结直肠癌这一大适应症上同样具备竞争力。

*尚未完全回答的问题**：目前的数据主要来自亚洲患者（含美国9例），种族差异对疗效的影响仍需更大规模全球III期研究（HARMONi-GI3）验证。此外，KRAS/BRAF突变亚组的疗效数据尚未充分披露，而这部分患者通常预后较差。

3. 硬纤维瘤领域：新药在关键终点上全面超越现有标准，但最佳治疗时长仍是未知数

Immunome（IMNM）的III期RINGSIDE研

[... middle omitted ...]

stment advice.

封面：ASCO 2026 副标题：肺癌、卵巢癌、肠癌新药数据一览

ASCO 周末结束，几个关键数据点值得反复看。我整理了最核心的几条逻辑链：

1️⃣ 肺癌：PD-1/VEGF 双抗改写标准治疗？ 投行研报指出，SMMT 合作伙伴在 ASCO 全体会议报告了 Ph3 HARMONi-6 研究。在肺鳞癌一线，双抗+化疗 vs PD-1+化疗： 中位 OS：27.9个月 vs 23.7个月 HR=0.66（死亡风险降34\%） 这是首个在 Ph3 中击败 PD-1+化疗方案的药物 亚组分析显示 PD-L1 阴性患者同样获益 安全性可管理，无新信号

2️⃣ 肠癌：早期数据支持全球 Ph3 Ph1/2 研究评估双抗+mFOLFOX6 方案： ORR：两剂量组均为70.8\% DCR：100\% 9个月无进展生存率：高剂量组76.1\% 研究包含部分美国患者，疗效与亚洲患者一致 KOL 认为数据支持正在进行的全球 Ph3 HARMONi-GI3 研究

\subsection{[R037] 日本IT服务业的真正拐点，藏在两家非覆盖公司的财报细节里}
{\small\color{refgray}归类：区域市场：亚洲央行鹰派转向，澳洲投资逻辑切换}

当市场还在为宏观利率和地缘政治焦虑时，日本IT服务板块的底层需求正在发生结构性的变化。某外资投行近期发布了一份特殊的研报——它并非针对其覆盖的头部公司，而是通过调研两家未覆盖的中间层系统集成商和网络集成商，得出了一个对全行业都有意义的判断：网络基础设施建设的结构性需求正在成为驱动利润增长的新引擎，而这一趋势被市场系统性低估。

报告调研的两家公司分别是互联网倡议公司（IIJ）和电通综研（Dentsu Soken）。前者在FY3/26业绩未达指引，但FY3/27指引却给出了两位数的利润增长预期；后者在汽车板块承压的背景下，凭借金融和自有软件业务实现了强劲开局。这两家公司的财报细节，揭示了一个正在发生的行业转折点。

1. 网络基础设施建设的订单爆发，不是周期性的，而是结构性的

IIJ的系统集成业务在4Q3/26的订单同比增长51\%，这是一个值得高度关注的信号。更关键的是，订单积压同比增长37.7\%，达到1589.1亿日元。这不是一次性项目带来的脉冲，而是来自政府机构和金融部门的持续需求。

报告明确指出，网络基础设施建设的需求是“结构性”的。这意味着推动力不是某个单一政策或预算周期，而是企业数字化转型和网络安全需求带来的长期资本开支。IIJ管理层预计这一趋势将在FY3/27持续，且由于维护业务的权重增加，毛利率有望改善。

对于投资者而言，这意味着：在IIJ身上验证的需求，很可能会传导至整个IT服务产业链。网络基础设施的订单增长，本质上是对数字基础设施投资的先行指标。当一家中型网络集成商的订单增速达到50\%以上时，头部公司的受益程度只会更大。

IIJ的FY3/26业绩低于指引，这看起来是一个负面信号。但报告的核心洞察在于：网络服务业务的盈利能力正在触底。FY3/27指引预计营业利润385亿日元，同比增长11\%，而这一增长将由网络基础设施业务驱动。

为什么利润率的底部确认如此关键？因为IIJ正在采取两个关键行动：一是对企事业客户实施提价，预计利润影响接近8亿日元；二是控制其他成本。这意味着公司不再单纯依赖规模扩张，而是开始通过定价权来修复利润率。

在MVNO/MVNE业务方面，IIJ目前没有提价计划，而是依靠低价策略获取市场份额。但这里有一个值得关注的张力：当内存价格上涨导致组件采购困难时，公司从1月份开始提前采购。在系统集成业务中，成本转嫁相对容易，但在网络服务业务中，由于合同周期限制，转嫁需要时间。这意味着短期内利润率可能承压，但中期来看，定价能力的提升将逐步显现。

3. 金融行业的IT支出正在加速，而汽车板块的疲软是结构性的

电通综研的数据提供了另一个维度。其1Q12/26订单同比增长20\%，主要驱动力来自金融板块，尤其是大型银行的系统开发需求。其自有软件产品，如合并会计软件STRAVIS和人力资源管理软件POSITIVE，在商社和电力公司等客户中获得了强劲需求。

与此同时，汽车板块的投资意愿因关税和业绩不佳而缺乏动力。电通综研正在将战略重心从汽车转向半导体。这一转向本身就说明了问题：汽车行业的IT支出疲软并非短期现象，而是受到贸易政策和行业结构性调整的双重影响。

这对整个IT服务行业意味着：依赖汽车客户的系统集成商可能面临持续的压力，而金融和半导体领域的IT支出则显示出更强的韧性。电通综研的订单积压同比增长33\%，管理层对达成指引的信心增强，这进一步验证了金融板块需求的可持续性。

4. 报告尚未完全回答的关键问题：成本转嫁的时滞与竞争格局的变化

尽管IIJ和电通综研的数据都指向积极方向，但报告中有两个关键问题没有完全展开。

第一个是成本转嫁的时滞问题。IIJ提到内存价格上涨导致组件采购困难，虽然系统集成业务可以相对顺畅地转嫁成本，但网络服务业务由于合同期限问题，转嫁需要时间。这意味着在

[... middle omitted ...]

务板块的投资机会：

第一，关注订单积压的增速而非单季度订单。IIJ的订单积压同比增长37.7\%，电通综研为33\%，这比单季度订单波动更能反映需求的持续性。

第二，区分结构性需求与周期性需求。网络基础设施、金融系统开发、自有软件产品属于结构性需求，而依赖汽车客户或单一项目的业务则可能面临周期性风险。

第三，追踪成本转嫁能力。在通胀环境下，能够快速将成本上涨转嫁给客户的公司将享有更好的利润率保护。系统集成业务通常比网络服务业务更容易转嫁成本。

第四，关注客户结构的多样性。电通综研从汽车向半导体的战略转移表明，客户集中度高的公司面临更大的下行风险。

*顺着这些未解问题继续深入，你会发现：真正的投资机会不在于判断行业趋势本身，而在于识别哪些公司拥有定价权和客户粘性，能够在结构性需求中扩大份额。欢迎来星球微信群里继续讨论，我们会在那里分享完整的研报解读和原始图表，深入拆解IIJ和电通综研的财务细节，以及这些数据对BIPROGY、NEC、富士通等头部公司的具体含义。**

\section{Disclaimer}
{\scriptsize\color{refgray}
This community is a paid learning and information-sharing community created for general educational purposes only. The membership fee is charged solely for access to community discussions, educational content, organization, and operational costs. It is not an advisory fee, management fee, performance fee, brokerage fee, or compensation for personalized financial, investment, legal, tax, or professional advice.\par
I participate and share information solely as an individual and for educational discussion among community members. I am not acting as a financial adviser, investment adviser, broker, dealer, portfolio manager, legal adviser, tax adviser, accountant, fiduciary, or any other licensed professional adviser. Nothing shared in this community constitutes, and should not be interpreted as, financial advice, investment advice, legal advice, tax advice, a recommendation, solicitation, offer, endorsement, instruction, or guarantee to buy, sell, hold, short, trade, or invest in any security, cryptocurrency, fund, derivative, strategy, or other financial product.\par
All content shared in this community, including messages, discussions, links, files, charts, examples, market commentary, watchlists, AI-generated or AI-polished summaries, and third-party materials, is general, impersonal, and educational in nature. It does not take into account any individual's financial situation, investment objectives, risk tolerance, investment horizon, liquidity needs, tax status, legal restrictions, or suitability requirements. No content should be relied upon as suitable for any specific person, account, portfolio, or financial decision.\par
Information shared in this community may be incomplete, inaccurate, outdated, speculative, AI-generated, AI-edited, or based on third-party internet sources that have not been independently verified. I make no representation or warranty as to the accuracy, completeness, timeliness, reliability, or suitability of any information shared.\par
Financial markets involve substantial risk, including the possible loss of principal. Past performance, historical data, hypothetical examples, simulated results, backtests, personal opinions, or educational examples do not guarantee future results. Each member is solely responsible for conducting their own research, exercising independent judgment, and making their own decisions. Before making any financial, investment, legal, or tax decision, members should consult qualified and licensed professionals.\par
Participation in this community, payment of any membership fee, or communication with me or other members does not create any adviser-client, fiduciary, professional, contractual, agency, partnership, employment, or other advisory relationship. By joining or participating in this community, you acknowledge and agree that you are solely responsible for your own decisions, actions, transactions, profits, losses, risks, and consequences, and that I shall not be liable for any direct, indirect, incidental, consequential, financial, legal, tax, or other loss or damage arising from or related to any content, discussion, information, or material shared in this community.\par
}
\end{document}
